\documentclass{ctexbook}
%\usepackage{ctexcap}
\usepackage{amsmath, amsthm, amssymb,cases,mathrsfs,diagbox,mathtools}
\numberwithin{equation}{chapter}
\usepackage{enumerate}
\usepackage{extarrows}
\usepackage{pifont}
\usepackage{paralist}
\usepackage[all]{xy}
\usepackage{xcolor, tikz,tikz-3dplot}
\usepackage{cancel}
\usepackage{physics}
%\usepackage{bbding}
%\usepackage{draftwatermark}
%\SetWatermarkScale{1}
\usepackage{hyperref}
\usepackage{pgfplots}
\pgfplotsset{compat=1.18}
\usepgfplotslibrary{groupplots}
\usepackage{makeidx}
\usepackage{esdiff}
\usepackage{listings}  % 渲染代码块
\usepackage{xcolor}  % 用到了 \color 命令

\definecolor{codegreen}{rgb}{0,0.6,0}
\definecolor{codegray}{rgb}{0.5,0.5,0.5}
\definecolor{codepurple}{rgb}{0.58,0,0.82}
\definecolor{backcolour}{rgb}{0.95,0.95,0.92}

% 定义 listings 风格，可以定义多个
\lstdefinestyle{mystyle}{
	% backgroundcolor=\color{backcolour},  % 背景色
	commentstyle= \color{red!50!green!50!blue!50},  % 注释的颜色
	keywordstyle= \color{blue!70},  % 关键字/程序语言中的保留字颜色
	numberstyle=\tiny\color{codegray},  % 左侧行号显示的颜色
	stringstyle=\color{codepurple},
	basicstyle=\ttfamily\footnotesize,
	breakatwhitespace=false,
	breaklines=true,  % 对过长的代码自动换行
	captionpos=b,
	keepspaces=true,
	numbers=left,  % 在左侧显示行号
	numbersep=5pt,
	showspaces=false,
	showstringspaces=false,  % 不显示字符串中的空格
	showtabs=false,
	tabsize=2,
	frame=single  % [none | single | shadowbox] 显示边框
}

\lstset{style=mystyle}  % 使用 listings 风格

\makeindex
\pgfplotsset{compat=1.18}
%\newtheorem*{solution}{解}
\hypersetup{
	pdftitle={大学数学},
	pdfauthor={唐家俊},
	pdfsubject={},
	pdfkeywords={},
	pdfcreator={},
	pdfstartpage={} }

\theoremstyle{plain}
\newtheorem{thm}{定理}[section]
\newtheorem{prop}{命题}[section]
\newtheorem{lem}{引理}[section]
\newtheorem{cor}{推论}[section]
\newtheorem{app}{应用}[section]


\theoremstyle{definition}
\newtheorem{defn}{定义}[section]
\newtheorem{exam}{例}[section]
\newtheorem{exer}{练习}[section]

\theoremstyle{remark}
\newtheorem{rem}{注记}[section]

\newenvironment{solution}{\begin{proof}[解]}{\end{proof}}
\newcommand*{\ow}{\overrightarrow}
\newcommand*{\Cov}{\mathrm{Cov}}

\newcommand{\R}{\mathbb{R}}
\newcommand{\T}{\mathcal{T}}
\newcommand{\C}{\mathcal{C}}
\newcommand{\osc}{\operatorname{osc}}


\allowdisplaybreaks[4]

%带圆圈数字命令
\newcommand*{\circled}[1]{\lower.7ex\hbox{\tikz\draw (0pt, 0pt) circle (.5em) node {\makebox[1em][c]{\small #1}};}}

\title{\bfseries 积分论(校验版)}
\author{T. Claesson, L. H\"ormander ~ 著}
\date{\today}

%❈%❈%❈%❈%❈%❈%❈%❈%❈%❈%❈%❈%❈%❈%❈%❈%❈%❈%❈%❈%❈%❈%❈%❈%❈%❈%❈%❈%❈%❈%❈%❈
\begin{document}
	\frontmatter
	\maketitle
	
	\renewcommand\contentsname{目~录}
	\tableofcontents
	
	\mainmatter
	
	\chapter*{前言}
	\addcontentsline{toc}{chapter}{前言}
	
	本书是在作者 1959 年秋季于斯德哥尔摩大学讲授课程的讲义基础上整理而成的。原讲义中, 第一章和第二章属于学士学位考试课程的一部分; 而全书内容则包含在硕士学位的必修课程之中, 至少在隆德大学当时是如此。
	
	我们希望这个新版本中的第一章和第二章, 可以作为学士学位阶段一门扩充课程的教材, 大约相当于 \(6\) 个学分; 同时也希望本书能够用于博士阶段的相关课程。由于目前基础课程应包含哪些内容并不完全明确, 本书增加了关于 Riemann 积分的一章作为第一章, 这样就不必把 Riemann 积分预设为读者已经熟悉的知识。
	
	习题部分补充了隆德大学考试中使用过的一些题目, 也收录了讨论班中提出的一些问题。附录中汇集了若干定义、记号和注释, 这样做的目的是避免正文部分变得过于冗长。
	
	关于 \(L^p\) 空间中函数与范数的若干命题, 应当结合 Banach 空间及其对偶空间中弱拓扑的学习来处理。为此, 读者还应阅读泛函分析教材中的相关内容。若想获得更完整的整体图景, 也应参考 Bourbaki 学派关于积分理论的著作。本讲义在整体体系上大体沿用了该书的处理方式。
	
	\begin{flushright}
		T. Claesson\\
		L. H\"ormander\\
		1970 年于隆德
	\end{flushright}
	
	\chapter{黎曼积分}

\section{黎曼积分的定义}

本章先在一个有限闭区间上建立黎曼积分. 这里``区间"指
\(\mathbb{R}^d\) 中边与坐标轴平行的闭长方体, 即
\[
I=[a_1,b_1]\times\cdots\times [a_d,b_d]
=\{x=(x_1,\ldots,x_d): a_j\leq x_j\leq b_j,\ 1\leq j\leq d\}.
\]
它的体积定义为
\[
m(I)=\prod_{j=1}^{d}(b_j-a_j).
\]
对于这类区间的子区间, 体积也按同样方式定义.

\subsection{分片常数函数及其积分}

设 \(f\) 是定义在 \(I\) 上的实值函数. 若存在 \(I\) 的一个有限剖分
\[
I=\bigcup_{k=1}^{N}I_k,
\]
其中各 \(I_k\) 的内部互不相交, 且 \(f\) 在每个 \(I_k\) 上为常数, 则称
\(f\) 为 \(I\) 上的分片常数函数. 记所有这类函数构成的集合为 \(\T\).
若 \(f\) 在 \(I_k\) 上的常值为 \(c_k\), 则定义
\[
\int_I f(x)\,\dd x=\sum_{k=1}^{N}c_k m(I_k).
\]
这个定义与剖分的选取无关. 事实上, 如果换用另一个剖分, 只要取两个
剖分的公共加细, 上式右端便化为同一个有限和.

由定义立刻得到以下两个基本性质:
\begin{enumerate}
	\item $\int_I f(x)\,\dd x\geq 0$, $f\in\T,\ f\geq 0$, \label{eq:step-positive}
	\item $\int_I \bigl(af(x)+bg(x)\bigr)\,\dd x
	=a\int_I f(x)\,\dd x+b\int_I g(x)\,\dd x$, $f,g\in\T,\ a,b\in\R$. \label{eq:step-linear}
\end{enumerate}
也就是说, 在分片常数函数上, 积分已经是一个正线性泛函.

\subsection{上积分与下积分}

现在令 \(f\) 为 \(I\) 上的有界实值函数. 用分片常数函数从上、从下夹住
\(f\), 定义它的上黎曼积分和下黎曼积分分别为
\begin{align}
	\overline{\int_I} f(x)\,\dd x
	&=\inf\left\{\int_I h(x)\,\dd x: h\in\T,\ f\leq h\right\},\label{eq:upper-riemann}\\
	\underline{\int_I} f(x)\,\dd x
	&=\sup\left\{\int_I h(x)\,\dd x: h\in\T,\ h\leq f\right\}.\label{eq:lower-riemann}
\end{align}
其中, 上确界与下确界是对所有满足相应不等式的分片常数函数取的.

显然总有
\[
\underline{\int_I} f(x)\,\dd x
\leq
\overline{\int_I} f(x)\,\dd x.
\]
事实上, 若 \(h_1,h_2\in\T\) 且 \(h_1\leq f\leq h_2\), 则
\[
0\leq \int_I \bigl(h_2(x)-h_1(x)\bigr)\,\dd x
=\int_I h_2(x)\,\dd x-\int_I h_1(x)\,\dd x.
\]
这说明任何从下逼近的积分值都不超过任何从上逼近的积分值.

\begin{defn}[黎曼可积]
	若
	\[
	\underline{\int_I} f(x)\,\dd x
	=
	\overline{\int_I} f(x)\,\dd x,
	\]
	则称 \(f\) 在 \(I\) 上黎曼可积. 此时二者的共同值记作
	\[
	\int_I f(x)\,\dd x,
	\]
	称为 \(f\) 在 \(I\) 上的黎曼积分.
\end{defn}

上述定义等价于下面的 \(\varepsilon\)-夹逼形式: \(f\) 黎曼可积当且仅当对任意
\(\varepsilon>0\), 存在 \(h_1,h_2\in\T\), 使得
\[
h_1\leq f\leq h_2,
\qquad
\int_I \bigl(h_2(x)-h_1(x)\bigr)\,\dd x<\varepsilon.
\]
这一定义的意义在于: 若一个有界函数可以被两个分片常数函数以任意小的
积分误差夹住, 那么它的积分就是确定的.

\begin{prop}
	若 \(f\) 在 \(I\) 上连续, 则 \(f\) 在 \(I\) 上黎曼可积.
\end{prop}

\begin{proof}
	由于 \(I\) 是紧集, 连续函数 \(f\) 在 \(I\) 上一致连续. 给定
	\(\varepsilon>0\), 取足够细的剖分 \(I=\bigcup I_k\), 使得 \(f\) 在每个
	\(I_k\) 上的振幅都小于 \(\varepsilon/m(I)\). 在每个 \(I_k\) 上分别取
	\[
	m_k=\inf_{x\in I_k}f(x),\qquad M_k=\sup_{x\in I_k}f(x),
	\]
	并令 \(h_1=m_k\), \(h_2=M_k\) 于 \(I_k\) 上. 则
	\(h_1,h_2\in\T\), \(h_1\leq f\leq h_2\), 且
	\[
	\int_I(h_2-h_1)\,\dd x
	=\sum_k(M_k-m_k)m(I_k)<\varepsilon.
	\]
	故 \(f\) 黎曼可积.
\end{proof}

进一步, 当剖分的细度趋于零时, 对任意选取的点 \(\xi_k\in I_k\), 连续函数
\(f\) 的黎曼和满足
\[
\sum_k f(\xi_k)m(I_k)\longrightarrow \int_I f(x)\,\dd x.
\]
这说明上下积分定义与通常的黎曼和定义相容.

\subsection{黎曼可积函数的线性性质}

若 \(a,b\geq 0\), 且 \(f,g\) 为有界实值函数, 则由上下积分的定义可得
\begin{align}
	\overline{\int_I} (af+bg)\,\dd x
	&\leq a\overline{\int_I} f\,\dd x
	+b\overline{\int_I} g\,\dd x,\label{eq:upper-sublinear}\\
	\underline{\int_I} (af+bg)\,\dd x
	&\geq a\underline{\int_I} f\,\dd x
	+b\underline{\int_I} g\,\dd x.\label{eq:lower-superlinear}
\end{align}
因此, 若 \(f\) 与 \(g\) 均黎曼可积, 则 \(af+bg\) 也黎曼可积, 且
\[
\int_I(af+bg)\,\dd x
=a\int_I f\,\dd x+b\int_I g\,\dd x.
\]
再由 \(-f\) 的情形可知, 上式对任意 \(a,b\in\R\) 都成立.

\begin{thm}[黎曼积分的线性性]
	设 \(f,g\) 在 \(I\) 上黎曼可积, \(a,b\in\R\). 则 \(af+bg\) 在 \(I\) 上黎曼可积, 且
	\[
	\int_I(af+bg)(x)\,\dd x
	=a\int_I f(x)\,\dd x+b\int_I g(x)\,\dd x.
	\]
\end{thm}

\begin{exer}
	证明: 若 \(f\) 或 \(g\) 中至少有一个是黎曼可积的, 则式
	\eqref{eq:upper-sublinear} 中相应的不等式可以强化为含有一个等号的形式.
	进一步证明: 若 \(f\) 为有界函数, 且对任意有界函数 \(g\) 都有
	\[
	\overline{\int_I}(f+g)\,\dd x
	=\overline{\int_I}f\,\dd x+
	\overline{\int_I}g\,\dd x,
	\]
	则 \(f\) 必为黎曼可积.
\end{exer}

\subsection{用连续函数逼近}

前面用分片常数函数定义黎曼积分. 事实上, 也可以用连续函数来定义同样的
上下积分. 记 \(\C(I)\) 为 \(I\) 上的连续函数全体, 则有
\begin{align*}
	\overline{\int_I} f\,\dd x
	&=\inf\left\{\int_I h\,\dd x: h\in\C(I),\ f\leq h\right\},\\
	\underline{\int_I} f\,\dd x
	&=\sup\left\{\int_I h\,\dd x: h\in\C(I),\ h\leq f\right\}.
\end{align*}

简要说明如下. 一方面, 连续函数也是可积函数, 所以用连续函数从上或从下逼近不会给出比原定义更强的结果. 另一方面, 任意分片常数函数都可以在积分意义下由连续函数逼近. 比如, 在每个小区间内部保持常值, 在边界附近作线性过渡, 由于边界附近可以取得任意薄, 因而由此带来的积分误差可以任意小.

因此 \(f\) 黎曼可积当且仅当对任意 \(\varepsilon>0\), 存在
\(h_1,h_2\in\C(I)\), 使得
\[
h_1\leq f\leq h_2,
\qquad
\int_I(h_2-h_1)\,\dd x<\varepsilon.
\]
由此还可推出一个常用判别准则:

\begin{prop}
	有界函数 \(f\) 在 \(I\) 上黎曼可积的充分必要条件是: 对任意
	\(\varepsilon>0\), 存在 \(g\in\C(I)\), 使得
	\[
	\overline{\int_I}|f-g|\,\dd x<\varepsilon.
	\]
\end{prop}

\begin{exer}[Darboux 定理]
	设 \(f\) 是 \(I\) 上有界函数. 对剖分 \(\mathcal{P}: I=\bigcup I_k\), 记
	\[
	M_k=\sup_{x\in I_k}f(x),\qquad m_k=\inf_{x\in I_k}f(x).
	\]
	证明: 当剖分细度趋于零时,
	\[
	\sum_k M_km(I_k)\longrightarrow \overline{\int_I} f\,\dd x,
	\qquad
	\sum_k m_km(I_k)\longrightarrow \underline{\int_I} f\,\dd x.
	\]
	提示: 可先对连续函数证明, 再利用连续函数逼近一般有界函数.
\end{exer}

\subsection{Jordan 测度}

黎曼积分还给出了测量集合大小的一种方法. 设 \(E\subset I\), 令
\(\chi_E\) 为 \(E\) 的特征函数, 即
\[
\chi_E(x)=
\begin{cases}
	1, & x\in E,\\
	0, & x\notin E.
\end{cases}
\]
定义 \(E\) 的 Jordan 外测度与 Jordan 内测度分别为
\[
m^*(E)=\overline{\int_I}\chi_E(x)\,\dd x,
\qquad
m_*(E)=\underline{\int_I}\chi_E(x)\,\dd x.
\]
若二者相等, 即 \(\chi_E\) 黎曼可积, 则称 \(E\) 为 Jordan 可测集, 并定义
\[
m(E)=\int_I\chi_E(x)\,\dd x.
\]
更一般地, 对任意有界函数 \(f\), 可定义它在 \(E\) 上的积分为
\[
\int_E f(x)\,\dd x:=\int_I f(x)\chi_E(x)\,\dd x,
\]
只要右端积分存在.

\begin{exam}[有理点集不是 Jordan 可测]
	令 \(E\) 为区间 \(I\) 中所有坐标均为有理数的点组成的集合. 任意非空开子区间
	都同时含有 \(E\) 中的点和 \(I\setminus E\) 中的点. 因而
	\[
	\overline{\int_I}\chi_E\,\dd x=m(I),
	\qquad
	\underline{\int_I}\chi_E\,\dd x=0.
	\]
	所以 \(E\) 不是 Jordan 可测集.
\end{exam}

\begin{exam}[Cantor 集]
	从 \([0,1]\) 出发, 第一步删去中间三分之一开区间, 得
	\[
	E_1=\left[0,\frac13\right]\cup\left[\frac23,1\right].
	\]
	第二步从剩下的两个闭区间中各删去中间三分之一开区间, 得
	\[
	E_2=\left[0,\frac19\right]\cup\left[\frac29,\frac13\right]
	\cup\left[\frac23,\frac79\right]\cup\left[\frac89,1\right].
	\]
	如此继续, 令
	\[
	C=\bigcap_{n=1}^{\infty}E_n.
	\]
	则 \(C\) 为 Cantor 集. 因为 \(E_n\) 由 \(2^n\) 个长度为 \(3^{-n}\) 的闭区间组成,
	所以
	\[
	m^*(C)\leq m(E_n)=\left(\frac23\right)^n,
	\qquad n=1,2,\ldots.
	\]
	令 \(n\to\infty\), 得 \(m^*(C)=0\). 因此 \(C\) 是 Jordan 可测集且
	\[
	m(C)=0.
	\]
	另一方面, Cantor 集不可数. 例如, 每个二进制序列 \((\varepsilon_n)\),
	\(\varepsilon_n\in\{0,2\}\), 都对应一个三进制展开
	\[
	x=\sum_{n=1}^{\infty}\frac{\varepsilon_n}{3^n},
	\]
	这些点均属于 \(C\), 由此可见 \(C\) 至少与二进制序列全体等势, 因而不可数.
\end{exam}

\subsection{一致极限与积分}

最后给出一个关于积分号下取极限的基本事实.

\begin{thm}
	设 \(f_n\) 是 \(I\) 上的黎曼可积函数, 且 \(f_n\) 在 \(I\) 上一致收敛到
	\(f\). 则 \(f\) 也黎曼可积, 并且
	\[
	\lim_{n\to\infty}\int_I f_n(x)\,\dd x
	=\int_I f(x)\,\dd x.
	\]
\end{thm}

\begin{proof}
	给定 \(\varepsilon>0\). 由于 \(f_n\to f\) 一致收敛, 存在 \(N\), 使得当
	\(n\geq N\) 时,
	\[
	f_n(x)-\varepsilon\leq f(x)\leq f_n(x)+\varepsilon,
	\qquad x\in I.
	\]
	于是
	\[
	\underline{\int_I} f_n\,\dd x-\varepsilon m(I)
	\leq
	\underline{\int_I} f\,\dd x
	\leq
	\overline{\int_I} f\,\dd x
	\leq
	\overline{\int_I} f_n\,\dd x+\varepsilon m(I).
	\]
	因为 \(f_n\) 黎曼可积, 上式左右两端相差不超过 \(2\varepsilon m(I)\).
	令 \(\varepsilon\to0\), 得 \(f\) 黎曼可积. 同时由
	\[
	\left|\int_I f\,\dd x-\int_I f_n\,\dd x\right|
	\leq \varepsilon m(I)
	\]
	可知积分收敛.
\end{proof}

\clearpage

\section{正项级数的积分}
\label{sec:positive-series-integration}

在处理非负项级数时, 有一个基本事实: 若 \(a_{mn}\geq 0\), 则
\[
\sum_{n=1}^{\infty}\sum_{m=1}^{\infty}a_{mn}
=
\sum_{m=1}^{\infty}\sum_{n=1}^{\infty}a_{mn}.
\]
这里允许两端同时取值为 \(+\infty\). 这是因为正项级数的和可以看作有限部分和的上确界, 而对非负数求和时, 有限求和的次序不会影响结果.

自然地, 我们希望把上式中的某一个求和号替换成积分号. 但是, 如果仍然局限于黎曼积分, 情形就没有这么完美. 本节的定理正是说明这一点: 对非负黎曼可积函数列, 逐项积分在“下积分”意义下仍然成立; 但即使每一项都是连续函数, 级数和也未必黎曼可积. 这正是引入勒贝格积分的一个重要动机.

\subsection{Dini 引理}

先给出一个将在证明中使用的紧致性结论.

\begin{lem}[Dini 引理]
	\label{lem:dini-riemann}
	设 \(f_n\) 是 \(I\) 上的非负连续函数. 若对每个 \(x\in I\),
	\[
	f_1(x)\geq f_2(x)\geq \cdots \geq 0,
	\qquad
	\lim_{n\to\infty}f_n(x)=0,
	\]
	则 \(f_n\) 在 \(I\) 上一致收敛于 \(0\). 因而
	\[
	\lim_{n\to\infty}\int_I f_n(x)\,\dd x=0.
	\]
\end{lem}

\begin{proof}
	给定 \(\varepsilon>0\). 对每个 \(x\in I\), 由 \(f_n(x)\downarrow 0\) 可取整数
	\(N_x\), 使得
	\[
	f_{N_x}(x)<\varepsilon.
	\]
	由于 \(f_{N_x}\) 连续, 存在 \(x\) 的一个邻域 \(U_x\), 使得当 \(y\in U_x\) 时,
	\[
	f_{N_x}(y)<\varepsilon.
	\]
	又由单调性, 当 \(n\geq N_x\) 且 \(y\in U_x\) 时,
	\[
	f_n(y)\leq f_{N_x}(y)<\varepsilon.
	\]
	这些 \(U_x\) 覆盖紧集 \(I\). 由有限覆盖定理, 存在有限个点
	\(x_1,\ldots,x_k\), 使得
	\[
	I\subset U_{x_1}\cup\cdots\cup U_{x_k}.
	\]
	令
	\[
	N=\max\{N_{x_1},\ldots,N_{x_k}\},
	\]
	则当 \(n\geq N\) 时, 对任意 \(y\in I\), 都有 \(f_n(y)<\varepsilon\). 因此
	\(f_n\to 0\) 一致收敛. 由上一节的一致极限定理,
	\[
	\int_I f_n(x)\,\dd x\longrightarrow 0.
	\]
\end{proof}

\subsection{正项级数的逐项积分}

本节中会遇到非负函数的无穷和, 它未必有界. 因此, 我们需要把下积分的定义稍微扩展一下: 若 \(F\) 是 \(I\) 上的下方有界函数, 定义
\[
\underline{\int_I}F(x)\,\dd x
=
\sup\left\{
\int_I h(x)\,\dd x:
h\in \C(I),\ h\leq F
\right\}.
\]
当 \(F\) 有界时, 这与上一节的下积分定义一致.

\begin{thm}[非负项级数的下积分]
	\label{thm:positive-series-riemann}
	设 \(u_n\) 是 \(I\) 上的非负黎曼可积函数, \(n=1,2,\ldots\). 令
	\[
	U(x)=\sum_{n=1}^{\infty}u_n(x),
	\]
	其中允许 \(U(x)=+\infty\). 则
	\[
	\underline{\int_I}U(x)\,\dd x
	=
	\sum_{n=1}^{\infty}\int_I u_n(x)\,\dd x.
	\]
	但是, 即使所有 \(u_n\) 都连续, 并且右端为有限值, 仍可能出现
	\(\sum_{n=1}^{\infty}u_n\) 不是黎曼可积的情形.
\end{thm}

\begin{proof}
	先证明等式.
	
	记
	\[
	S(x)=\sum_{n=1}^{\infty}u_n(x).
	\]
	由于对任意 \(N\),
	\[
	\sum_{n=1}^{N}u_n(x)\leq S(x),
	\]
	且有限和 \(\sum_{n=1}^{N}u_n\) 黎曼可积, 所以由下积分的定义可得
	\[
	\underline{\int_I}S(x)\,\dd x
	\geq
	\int_I \sum_{n=1}^{N}u_n(x)\,\dd x
	=
	\sum_{n=1}^{N}\int_I u_n(x)\,\dd x.
	\]
	令 \(N\to\infty\), 得
	\[
	\underline{\int_I}S(x)\,\dd x
	\geq
	\sum_{n=1}^{\infty}\int_I u_n(x)\,\dd x.
	\]
	
	下面证明反向不等式. 给定 \(\varepsilon>0\). 由下积分的定义, 可取连续函数
	\(H\in\C(I)\), 使得
	\[
	H(x)\leq S(x),
	\qquad
	\underline{\int_I}S(x)\,\dd x
	<
	\int_I H(x)\,\dd x+\frac{\varepsilon}{2}.
	\]
	由于每个 \(u_n\) 黎曼可积, 又由上一节的连续函数逼近准则, 对每个 \(n\) 可取连续函数 \(h_n\in\C(I)\), 使得
	\[
	u_n(x)\leq h_n(x),
	\qquad
	\int_I h_n(x)\,\dd x
	<
	\int_I u_n(x)\,\dd x+\frac{\varepsilon}{2^{n+1}}.
	\]
	于是对每个 \(x\in I\),
	\[
	\sum_{n=1}^{\infty}h_n(x)\geq \sum_{n=1}^{\infty}u_n(x)=S(x)\geq H(x).
	\]
	定义
	\[
	F_N(x)
	=
	\max\left\{
	H(x)-\sum_{n=1}^{N}h_n(x),\,0
	\right\}.
	\]
	则 \(F_N\) 是 \(I\) 上的非负连续函数, 且 \(F_N(x)\downarrow 0\). 由引理
	\ref{lem:dini-riemann},
	\[
	\int_I F_N(x)\,\dd x\longrightarrow 0.
	\]
	另一方面,
	\[
	H(x)\leq \sum_{n=1}^{N}h_n(x)+F_N(x),
	\]
	从而
	\[
	\int_I H(x)\,\dd x
	\leq
	\sum_{n=1}^{N}\int_I h_n(x)\,\dd x
	+
	\int_I F_N(x)\,\dd x.
	\]
	于是
	\begin{align*}
		\underline{\int_I}S(x)\,\dd x
		&<
		\int_I H(x)\,\dd x+\frac{\varepsilon}{2}\\
		&\leq
		\sum_{n=1}^{N}\int_I h_n(x)\,\dd x
		+
		\int_I F_N(x)\,\dd x
		+\frac{\varepsilon}{2}\\
		&<
		\sum_{n=1}^{N}\int_I u_n(x)\,\dd x
		+
		\sum_{n=1}^{N}\frac{\varepsilon}{2^{n+1}}
		+
		\int_I F_N(x)\,\dd x
		+\frac{\varepsilon}{2}.
	\end{align*}
	令 \(N\to\infty\), 得
	\[
	\underline{\int_I}S(x)\,\dd x
	\leq
	\sum_{n=1}^{\infty}\int_I u_n(x)\,\dd x+\varepsilon.
	\]
	由于 \(\varepsilon>0\) 任意, 反向不等式成立. 因此
	\[
	\underline{\int_I}\sum_{n=1}^{\infty}u_n(x)\,\dd x
	=
	\sum_{n=1}^{\infty}\int_I u_n(x)\,\dd x.
	\]
	
	下面构造反例, 说明级数和不一定黎曼可积. 为简单起见, 假设 \(m(I)>0\). 将 \(I\) 中所有坐标均为有理数的点排成一个序列
	\[
	y_1,y_2,\ldots.
	\]
	对每个 \(n\), 取连续函数 \(f_n\in\C(I)\), 使得
	\[
	0\leq f_n\leq 1,
	\qquad
	f_n(y_n)=1,
	\qquad
	\int_I f_n(x)\,\dd x<\frac{m(I)}{2^{n+1}}.
	\]
	这样的函数可以取为以 \(y_n\) 为中心、支集足够小的连续尖峰函数.
	
	定义
	\[
	u_1=f_1,
	\qquad
	u_n=f_n(1-f_1)(1-f_2)\cdots(1-f_{n-1}),
	\quad n\geq 2.
	\]
	则每个 \(u_n\) 都是非负连续函数, 且
	\[
	0\leq u_n\leq f_n,
	\]
	因此
	\[
	\sum_{n=1}^{\infty}\int_I u_n(x)\,\dd x
	\leq
	\sum_{n=1}^{\infty}\int_I f_n(x)\,\dd x
	<
	\frac{m(I)}{2}.
	\]
	另一方面, 对任意 \(N\),
	\[
	\sum_{n=1}^{N}u_n(x)
	=
	1-\prod_{n=1}^{N}\bigl(1-f_n(x)\bigr).
	\]
	因此
	\[
	0\leq \sum_{n=1}^{\infty}u_n(x)\leq 1.
	\]
	并且, 如果 \(x=y_k\), 则 \(f_k(y_k)=1\), 所以
	\[
	\sum_{n=1}^{\infty}u_n(y_k)=1.
	\]
	也就是说, 级数和
	\[
	U(x)=\sum_{n=1}^{\infty}u_n(x)
	\]
	在所有有理坐标点上都等于 \(1\).
	
	由于 \(I\) 的任一非空子区间都含有有理坐标点, 所以在任一这样的子区间上 \(U\) 的上确界都是 \(1\). 因而
	\[
	\overline{\int_I}U(x)\,\dd x=m(I).
	\]
	但是由已经证明的逐项积分公式,
	\[
	\underline{\int_I}U(x)\,\dd x
	=
	\sum_{n=1}^{\infty}\int_I u_n(x)\,\dd x
	<
	\frac{m(I)}{2}.
	\]
	于是
	\[
	\underline{\int_I}U(x)\,\dd x
	<
	\overline{\int_I}U(x)\,\dd x,
	\]
	所以 \(U\) 不是黎曼可积的.
\end{proof}

\begin{rem}
	\label{rem:riemann-defect-positive-series}
	定理 \ref{thm:positive-series-riemann} 表明, 对非负函数列逐项积分时, 黎曼积分存在一个本质缺陷: 级数和这一端不得不用下积分来理解, 而逐项积分那一端却是通常的黎曼积分. 因此, 等式虽然成立, 但两端的“积分”并不是同一种意义下的积分.
	
	下一章将通过推广积分概念来弥补这一缺陷. 在勒贝格积分框架下, 非负函数列的逐项积分将具有更自然的形式.
\end{rem}

\clearpage

\section{累次积分}
\label{sec:iterated-riemann-integrals}

设 \(I\subset\R^d\) 和 \(J\subset\R^e\) 都是有限闭区间. 我们分别用 \(x\) 和 \(y\) 表示其中的变量, 即
\[
I\times J=\bigl\{(x,y):x\in I,\ y\in J\bigr\}\subset\R^{d+e}.
\]
若 \(f\) 是 \(I\times J\) 上的有界实值函数, 则对固定的 \(x\in I\), 可以把
\(y\mapsto f(x,y)\) 看成 \(J\) 上的有界函数, 从而可以讨论它关于 \(y\) 的上积分和下积分. 因此得到两个关于 \(x\) 的函数:
\[
\Phi(x)=\overline{\int_J} f(x,y)\,\dd y,
\qquad
\phi(x)=\underline{\int_J} f(x,y)\,\dd y.
\]
一般来说, \(\Phi\) 与 \(\phi\) 未必相等; 也就是说, 即使 \(f\) 在乘积区间上表现良好, 固定 \(x\) 后得到的截面函数也未必黎曼可积. 本节的目的就是说明这种现象, 并给出黎曼积分框架下能够成立的累次积分结论.

\subsection{上下积分的不等式}

先证明一个基本不等式. 它是后面所有结论的出发点.

\begin{prop}\label{prop:upper-lower-iterated-riemann}
	设 \(f\) 是 \(I\times J\) 上的有界实值函数. 则
	\begin{align}
		\overline{\int_I}\left(\overline{\int_J}f(x,y)\,\dd y\right)\dd x
		&\leq
		\overline{\int_{I\times J}} f(x,y)\,\dd x\dd y,\label{eq:upper-iterated-riemann}\\
		\underline{\int_I}\left(\underline{\int_J}f(x,y)\,\dd y\right)\dd x
		&\geq
		\underline{\int_{I\times J}} f(x,y)\,\dd x\dd y.\label{eq:lower-iterated-riemann}
	\end{align}
\end{prop}

\begin{proof}
	先证 \eqref{eq:upper-iterated-riemann}. 若 \(h\) 是 \(I\times J\) 上的分片常数函数且 \(f\leq h\), 则对每个固定的 \(x\), 有
	\[
	\overline{\int_J}f(x,y)\,\dd y
	\leq
	\int_J h(x,y)\,\dd y.
	\]
	于是
	\[
	\overline{\int_I}\left(\overline{\int_J}f(x,y)\,\dd y\right)\dd x
	\leq
	\int_I\left(\int_J h(x,y)\,\dd y\right)\dd x.
	\]
	当 \(h\) 为分片常数函数时, 累次积分与乘积区间上的积分显然相等, 即
	\[
	\int_I\left(\int_J h(x,y)\,\dd y\right)\dd x
	=
	\int_{I\times J}h(x,y)\,\dd x\dd y.
	\]
	这里用到的只是体积公式
	\[
	m(I_k\times J_l)=m(I_k)m(J_l).
	\]
	对所有满足 \(f\leq h\) 的分片常数函数 \(h\) 取下确界, 即得
	\[
	\overline{\int_I}\left(\overline{\int_J}f(x,y)\,\dd y\right)\dd x
	\leq
	\overline{\int_{I\times J}} f(x,y)\,\dd x\dd y.
	\]
	这就是 \eqref{eq:upper-iterated-riemann}.
	
	至于 \eqref{eq:lower-iterated-riemann}, 可以仿照上面的论证直接证明; 也可以将 \eqref{eq:upper-iterated-riemann} 应用于 \(-f\), 再利用上下积分之间的关系
	\[
	\overline{\int}(-f)=-\underline{\int}f.
	\]
\end{proof}

\subsection{连续函数的累次积分}

对于连续函数, 上述不等式立即给出通常形式的累次积分公式.

\begin{thm}[连续函数的累次积分]\label{thm:continuous-fubini-riemann}
	若 \(f\) 是 \(I\times J\) 上的连续函数, 则
	\[
	F(x)=\int_J f(x,y)\,\dd y
	\]
	是 \(I\) 上的连续函数, 并且
	\[
	\int_I\left(\int_J f(x,y)\,\dd y\right)\dd x
	=
	\int_{I\times J}f(x,y)\,\dd x\dd y.
	\]
\end{thm}

\begin{proof}
	因为 \(I\times J\) 是紧集, \(f\) 在 \(I\times J\) 上一致连续. 若 \(x_n\to x\), 则
	\[
	f(x_n,y)\longrightarrow f(x,y)
	\]
	关于 \(y\in J\) 一致成立. 由上一节关于一致极限与积分交换的定理可知
	\[
	\int_J f(x_n,y)\,\dd y
	\longrightarrow
	\int_J f(x,y)\,\dd y.
	\]
	因此 \(F\) 连续.
	
	由于连续函数黎曼可积, 对每个 \(x\in I\) 都有
	\[
	\overline{\int_J}f(x,y)\,\dd y
	=
	\underline{\int_J}f(x,y)\,\dd y
	=
	\int_J f(x,y)\,\dd y.
	\]
	同时 \(f\) 在 \(I\times J\) 上黎曼可积. 将命题 \ref{prop:upper-lower-iterated-riemann} 中的两个不等式合在一起, 就得到
	\[
	\int_I F(x)\,\dd x
	=
	\int_{I\times J}f(x,y)\,\dd x\dd y.
	\]
\end{proof}

\begin{rem}
	由对称性, 定理 \ref{thm:continuous-fubini-riemann} 中也可以先对 \(x\) 积分再对 \(y\) 积分. 因而连续函数满足
	\[
	\int_I\left(\int_J f(x,y)\,\dd y\right)\dd x
	=
	\int_J\left(\int_I f(x,y)\,\dd x\right)\dd y
	=
	\int_{I\times J}f(x,y)\,\dd x\dd y.
	\]
\end{rem}

\subsection{一般黎曼可积函数的情形}

现在设 \(f\) 是 \(I\times J\) 上的黎曼可积函数. 这时不能贸然断言
\(y\mapsto f(x,y)\) 对每个 \(x\) 都黎曼可积, 因而不能直接写出普通意义下的累次积分. 但是上下积分仍然给出一个可用的结论.

\begin{thm}[乘积可积函数的上下累次积分]\label{thm:riemann-integrable-iterated-upper-lower}
	设 \(f\) 在 \(I\times J\) 上黎曼可积. 定义
	\[
	\Phi(x)=\overline{\int_J}f(x,y)\,\dd y,
	\qquad
	\phi(x)=\underline{\int_J}f(x,y)\,\dd y.
	\]
	则 \(\Phi\) 与 \(\phi\) 都在 \(I\) 上黎曼可积, 且
	\[
	\int_I \Phi(x)\,\dd x
	=
	\int_I \phi(x)\,\dd x
	=
	\int_{I\times J}f(x,y)\,\dd x\dd y.
	\]
	特别地, 若对每个 \(x\in I\), 函数 \(y\mapsto f(x,y)\) 都在 \(J\) 上黎曼可积, 则
	\[
	\int_I\left(\int_J f(x,y)\,\dd y\right)\dd x
	=
	\int_{I\times J}f(x,y)\,\dd x\dd y.
	\]
\end{thm}

\begin{proof}
	由命题 \ref{prop:upper-lower-iterated-riemann} 得
	\[
	\overline{\int_I}\Phi(x)\,\dd x
	\leq
	\overline{\int_{I\times J}}f(x,y)\,\dd x\dd y,
	\]
	和
	\[
	\underline{\int_I}\phi(x)\,\dd x
	\geq
	\underline{\int_{I\times J}}f(x,y)\,\dd x\dd y.
	\]
	由于 \(f\) 在 \(I\times J\) 上黎曼可积, 乘积区间上的上积分与下积分都等于
	\[
	A=\int_{I\times J}f(x,y)\,\dd x\dd y.
	\]
	另一方面, 对每个 \(x\in I\), 都有 \(\phi(x)\leq \Phi(x)\), 因此
	\[
	\underline{\int_I}\phi(x)\,\dd x
	\leq
	\overline{\int_I}\phi(x)\,\dd x
	\leq
	\overline{\int_I}\Phi(x)\,\dd x.
	\]
	综合以上不等式, 得
	\[
	A
	\leq
	\underline{\int_I}\phi(x)\,\dd x
	\leq
	\overline{\int_I}\phi(x)\,\dd x
	\leq
	\overline{\int_I}\Phi(x)\,\dd x
	\leq A.
	\]
	于是所有项都等于 \(A\). 特别地, \(\phi\) 黎曼可积且
	\[
	\int_I\phi(x)\,\dd x=A.
	\]
	又由于 \(\phi\leq \Phi\), 且 \(\overline{\int_I}\Phi\leq A\), 同时
	\[
	A=\int_I\phi(x)\,\dd x\leq \underline{\int_I}\Phi(x)\,\dd x,
	\]
	所以 \(\Phi\) 也黎曼可积, 且
	\[
	\int_I\Phi(x)\,\dd x=A.
	\]
	若每个截面函数 \(y\mapsto f(x,y)\) 均黎曼可积, 则
	\[
	\Phi(x)=\phi(x)=\int_J f(x,y)\,\dd y,
	\]
	从而得到最后的公式.
\end{proof}

\begin{rem}
	定理 \ref{thm:riemann-integrable-iterated-upper-lower} 的重点在于: 乘积区间上的黎曼可积性并不保证每个截面都黎曼可积. 因此, 在黎曼积分理论中, 累次积分公式必须谨慎表述. 后面的勒贝格积分会把这件事处理得更自然.
\end{rem}

\subsection{一个说明截面不可积的例子}

下面的例子说明: 即使 \(f\) 在 \(I\times J\) 上黎曼可积, 对某些固定的 \(x\), 内层积分
\[
\int_J f(x,y)\,\dd y
\]
仍然可能没有定义.

\begin{exam}\label{exam:riemann-product-integrable-bad-sections}
	令 \(I=J=[0,1]\). 定义
	\[
	g(x)=
	\begin{cases}
		\dfrac{1}{q}, & x=\dfrac{p}{q},\ p,q\text{互质},\ q>0,\\[0.6em]
		0, & x\text{为无理数},
	\end{cases}
	\]
	以及
	\[
	h(y)=
	\begin{cases}
		1, & y\text{为有理数},\\
		0, & y\text{为无理数}.
	\end{cases}
	\]
	则 \(g\) 是黎曼可积函数, 且
	\[
	\int_0^1 g(x)\,\dd x=0.
	\]
	另一方面,
	\[
	\overline{\int_0^1}h(y)\,\dd y=1,
	\qquad
	\underline{\int_0^1}h(y)\,\dd y=0.
	\]
	令
	\[
	f(x,y)=g(x)h(y).
	\]
	由于
	\[
	0\leq f(x,y)\leq g(x),
	\]
	而 \((x,y)\mapsto g(x)\) 在 \([0,1]\times[0,1]\) 上黎曼可积且积分为 \(0\), 所以 \(f\) 在 \(I\times J\) 上黎曼可积, 并且
	\[
	\int_{I\times J}f(x,y)\,\dd x\dd y=0.
	\]
	但是, 固定 \(x\) 后,
	\[
	\overline{\int_0^1}f(x,y)\,\dd y
	-
	\underline{\int_0^1}f(x,y)\,\dd y
	=g(x).
	\]
	因此, 当 \(x\) 为有理数时, \(g(x)\neq 0\), 函数 \(y\mapsto f(x,y)\) 不是黎曼可积的. 于是
	\[
	\int_0^1 f(x,y)\,\dd y
	\]
	对这些 \(x\) 没有定义. 这说明在一般黎曼积分理论中, 公式
	\[
	\int_{I\times J}f(x,y)\,\dd x\dd y
	=
	\int_I\left(\int_J f(x,y)\,\dd y\right)\dd x
	\]
	未必有意义.
\end{exam}

\begin{proof}
	这里只补充说明 \(g\) 的可积性. 显然 \(g\geq0\), 且任意子区间都含有无理点, 所以任意剖分对应的下和都是 \(0\), 从而
	\[
	\underline{\int_0^1}g(x)\,\dd x=0.
	\]
	给定 \(\varepsilon>0\). 取 \(Q\) 足够大, 使 \(1/Q<\varepsilon/2\). 在 \([0,1]\) 中, 分母 \(q\leq Q\) 的既约有理数只有有限多个. 用总长度小于 \(\varepsilon/2\) 的有限个小区间覆盖这些点. 在余下部分, 若 \(x=p/q\) 且 \(q>Q\), 则
	\[
	g(x)<\frac{1}{Q}<\frac{\varepsilon}{2};
	\]
	若 \(x\) 为无理数, 则 \(g(x)=0\). 因此可使上和任意小, 从而
	\[
	\overline{\int_0^1}g(x)\,\dd x=0.
	\]
	故 \(g\) 黎曼可积且积分为 \(0\).
\end{proof}

\begin{exer}\label{exer:product-function-upper-lower-equality}
	设
	\[
	f(x,y)=g(x)h(y),
	\]
	其中 \(g\) 和 \(h\) 分别是 \(I\) 与 \(J\) 上的非负有界函数. 证明命题 \ref{prop:upper-lower-iterated-riemann} 中的两个不等式实际上都是等式.
	
	特别地, 若取 \(g=\chi_E\), \(h=\chi_F\), 其中 \(E\subset I\), \(F\subset J\), 则
	\[
	m^*(E\times F)=m^*(E)m^*(F),
	\qquad
	m_*(E\times F)=m_*(E)m_*(F).
	\]
	若 \(E\) 和 \(F\) 都是 Jordan 可测集, 则
	\[
	m(E\times F)=m(E)m(F).
	\]
\end{exer}

\begin{exer}\label{exer:nonnegative-assumption-essential}
	上一题中条件 \(g,h\geq0\) 是本质的. 令 \(I=[0,1]\), \(J=[0,2]\), 并定义
	\[
	g(x)=
	\begin{cases}
		1, & x\text{为有理数},\\
		-x, & x\text{为无理数},
	\end{cases}
	\qquad
	h(y)=
	\begin{cases}
		-1, & y\text{为有理数},\\
		y, & y\text{为无理数}.
	\end{cases}
	\]
	令
	\[
	f(x,y)=g(x)h(y)+c,
	\]
	其中 \(c\) 为任意常数. 证明下面三个上积分互不相等:
	\[
	\overline{\int_{I\times J}} f(x,y)\,\dd x\dd y,
	\qquad
	\overline{\int_I}\left(\overline{\int_J}f(x,y)\,\dd y\right)\dd x,
	\qquad
	\overline{\int_J}\left(\overline{\int_I}f(x,y)\,\dd x\right)\dd y.
	\]
\end{exer}

\clearpage

\section{变量替换}
\label{sec:change-of-variables-riemann}

在本节中, 我们把黎曼积分从有限区间上的函数推广到整个空间中具有紧致支集的连续函数, 然后讨论变量替换公式. 这种处理方式的好处是: 平移、线性变换和一般的微分同胚都可以放在统一的框架中讨论.

记 \(\C_0(\R^d)\) 为 \(\R^d\) 上所有具有紧致支集的连续实值函数全体, 记 \(\C_0^+(\R^d)\) 为其中的非负函数全体. 若 \(f\in \C_0(\R^d)\), 取一个有限闭区间 \(I\subset\R^d\), 使得
\[
\operatorname{supp} f\subset I.
\]
定义
\[
\int_{\R^d}f(x)\,\dd x:=\int_I f(x)\,\dd x.
\]
这个定义与 \(I\) 的选择无关. 事实上, 若 \(I_1\) 与 \(I_2\) 都包含 \(\operatorname{supp} f\), 则在一个共同包含二者的区间上比较即可.

由有限区间上黎曼积分的性质, 立即得到
\begin{align}
	\int_{\R^d}\bigl(af(x)+bg(x)\bigr)\,\dd x
	&=a\int_{\R^d}f(x)\,\dd x
	+b\int_{\R^d}g(x)\,\dd x,
	\label{eq:c0-integral-linear}\\
	\int_{\R^d}f(x)\,\dd x&\geq 0,
	\qquad f\in\C_0^+(\R^d),
	\label{eq:c0-integral-positive}
\end{align}
其中 \(a,b\in\R\), \(f,g\in\C_0(\R^d)\).

\subsection{仿射变换与积分}

设
\[
\phi:\R^d\longrightarrow \R^d
\]
为一个仿射变换, 即
\[
\phi(x)=Ax+x_0,
\label{eq:affine-map}
\]
其中 \(x_0\in\R^d\), \(A\in \operatorname{GL}(d,\R)\). 若 \(f\) 是 \(\R^d\) 上的函数, 我们记
\[
\phi^*f=f\circ\phi,
\qquad
(\phi^*f)(x)=f(\phi(x)).
\]
显然, 若 \(f\in\C_0(\R^d)\), 则
\[
\phi^*f\in\C_0(\R^d).
\]

先考虑最简单的情形: \(A\) 为对角矩阵. 此时若 \(f\) 是一个区间的特征函数, 则 \(\phi^*f\) 仍是一个区间的特征函数. 例如, 若
\[
I=\{x:a_j\leq x_j\leq b_j,
\quad j=1,\ldots,d\},
\]
且 \(f=\chi_I\), 而 \(A=\operatorname{diag}(\alpha_1,\ldots,\alpha_d)\), 则 \(\phi^{-1}(I)\) 仍是一个坐标区间. 由区间体积的定义可得
\[
\int_{\R^d}\phi^*f(x)\,\dd x
=\frac{1}{|\det A|}\int_{\R^d}f(x)\,\dd x.
\]
也就是说
\begin{equation}
	|\det A|\int_{\R^d}\phi^*f(x)\,\dd x
	=
	\int_{\R^d}f(x)\,\dd x.
	\label{eq:affine-change-diagonal}
\end{equation}
由线性性, 公式 \eqref{eq:affine-change-diagonal} 对所有分片常数函数成立. 再由连续函数的黎曼积分定义, 它也对所有 \(f\in\C_0(\R^d)\) 成立.

特别地, 若
\[
f_h(x)=f(x-h),
\qquad h\in\R^d,
\]
则由 \eqref{eq:affine-change-diagonal} 在 \(A=I_d\) 的情形得到平移不变性:
\begin{equation}
	\int_{\R^d}f_h(x)\,\dd x
	=
	\int_{\R^d}f(x)\,\dd x.
	\label{eq:translation-invariance-riemann}
\end{equation}

下面说明: 线性性、正性和平移不变性几乎完全刻画了 \(\C_0(\R^d)\) 上的黎曼积分.

\begin{thm}\label{thm:translation-invariant-positive-functional}
	设 \(L:\C_0(\R^d)\to\R\) 是一个泛函, 满足:
	\begin{enumerate}
		\item \(L(af+bg)=aL(f)+bL(g)\), 对任意 \(f,g\in\C_0(\R^d)\) 与 \(a,b\in\R\) 成立;
		\item 若 \(f\in\C_0^+(\R^d)\), 则 \(L(f)\geq 0\);
		\item \(L(f_h)=L(f)\), 对任意 \(f\in\C_0(\R^d)\) 与 \(h\in\R^d\) 成立, 其中 \(f_h(x)=f(x-h)\).
	\end{enumerate}
	则存在常数 \(C\), 使得
	\[
	L(f)=C\int_{\R^d}f(x)\,\dd x,
	\qquad f\in\C_0(\R^d).
	\]
\end{thm}

\begin{proof}
	先注意, 由正性和线性性可知, 若 \(f\leq g\), 则
	\[
	L(f)\leq L(g).
	\]
	因此 \(L\) 关于函数的逐点大小是单调的.
	
	进一步, 若一列函数 \(f_n\in\C_0(\R^d)\) 的支集都包含在同一个紧集 \(K\) 中, 且 \(f_n\) 一致收敛到 \(f\), 则
	\[
	L(f_n)\longrightarrow L(f).
	\]
	事实上, 取 \(g\in\C_0^+(\R^d)\), 使得 \(g\geq 1\) 于 \(K\) 上. 令
	\[
	\varepsilon_n=\sup_{x\in\R^d}|f_n(x)-f(x)|.
	\]
	则
	\[
	-\varepsilon_n g\leq f_n-f\leq \varepsilon_n g,
	\]
	从而
	\[
	|L(f_n)-L(f)|\leq \varepsilon_n L(g)\to0.
	\]
	
	现在取 \(f,g\in\C_0(\R^d)\), 且假设它们的支集都包含于某个区间 \(I\) 中. 设
	\[
	I=\bigcup_{j=1}^{N}I_j
	\]
	是一个剖分, 各 \(I_j\) 没有公共内点. 在每个 \(I_j\) 中任取一点 \(h_j\), 并定义
	\[
	F(x)=\sum_{j=1}^{N}f(x-h_j)g(h_j)m(I_j).
	\]
	由平移不变性和线性性,
	\begin{equation}
		L(F)=L(f)\sum_{j=1}^{N}g(h_j)m(I_j).
		\label{eq:functional-riemann-sum}
	\end{equation}
	当剖分充分细时, \(F(x)\) 逼近卷积函数
	\[
	(f*g)(x)=\int_{\R^d}f(x-h)g(h)\,\dd h.
	\]
	更具体地说, 因为 \((x,h)\mapsto f(x-h)g(h)\) 在一个紧集上一致连续, 所以当剖分细度趋于零时,
	\[
	F\longrightarrow f*g
	\]
	一致收敛. 令剖分细度趋于零, 由 \eqref{eq:functional-riemann-sum} 得
	\begin{equation}
		L(f*g)=L(f)\int_{\R^d}g(h)\,\dd h.
		\label{eq:functional-convolution-1}
	\end{equation}
	
	另一方面, 由变量替换 \(h=x-y\), 可知
	\[
	(f*g)(x)=(g*f)(x).
	\]
	将 \eqref{eq:functional-convolution-1} 中 \(f\) 与 \(g\) 对调, 得
	\begin{equation}
		L(f)\int_{\R^d}g(x)\,\dd x
		=
		L(g)\int_{\R^d}f(x)\,\dd x.
		\label{eq:functional-convolution-2}
	\end{equation}
	固定一个 \(g\in\C_0(\R^d)\), 使得
	\[
	\int_{\R^d}g(x)\,\dd x\neq0.
	\]
	于是由 \eqref{eq:functional-convolution-2} 得
	\[
	L(f)=C\int_{\R^d}f(x)\,\dd x,
	\qquad
	C=\frac{L(g)}{\int_{\R^d}g(x)\,\dd x}.
	\]
\end{proof}

\begin{rem}
	定理 \ref{thm:translation-invariant-positive-functional} 的意思是: 除了一个常数因子以外, \(\R^d\) 上的黎曼积分正是 \(\C_0(\R^d)\) 上的平移不变正线性泛函. 这说明平移不变性不是积分的附属性质, 而是积分概念的核心结构之一.
\end{rem}

现在回到一般可逆线性变换. 对 \(A\in\operatorname{GL}(d,\R)\), 考虑泛函
\[
L_A(f)=\int_{\R^d}f(Ax)\,\dd x,
\qquad f\in\C_0(\R^d).
\]
它显然是正线性的. 而且
\[
L_A(f_h)=\int_{\R^d}f(Ax-h)\,\dd x
=\int_{\R^d}f(A(x-A^{-1}h))\,\dd x
=L_A(f),
\]
所以它也平移不变. 由定理 \ref{thm:translation-invariant-positive-functional}, 存在只依赖于 \(A\) 的常数 \(C(A)\), 使得
\begin{equation}
	\int_{\R^d}f(Ax)\,\dd x
	=C(A)\int_{\R^d}f(x)\,\dd x.
	\label{eq:linear-transform-constant}
\end{equation}

若 \(A,B\in\operatorname{GL}(d,\R)\), 则
\[
\int_{\R^d}f(ABx)\,\dd x
=C(B)\int_{\R^d}f(Ax)\,\dd x
=C(A)C(B)\int_{\R^d}f(x)\,\dd x,
\]
故
\begin{equation}
	C(AB)=C(A)C(B).
	\label{eq:C-multiplicative}
\end{equation}
若 \(A\) 是正交矩阵, 则 \(C(A)=1\). 事实上, 对形如
\[
f(x)=\psi(\|x\|)
\]
的非零函数, 有 \(f(Ax)=f(x)\), 从而由 \eqref{eq:linear-transform-constant} 得 \(C(A)=1\). 又前面已经知道, 当 \(D\) 为可逆对角矩阵时,
\[
C(D)=\frac{1}{|\det D|}.
\]

任意 \(A\in\operatorname{GL}(d,\R)\) 都可以写成
\[
A=O_1DO_2,
\]
其中 \(O_1,O_2\) 为正交矩阵, \(D\) 为可逆对角矩阵. 因此
\[
C(A)=C(O_1)C(D)C(O_2)=\frac{1}{|\det D|}
=\frac{1}{|\det A|}.
\]
所以对任意 \(A\in\operatorname{GL}(d,\R)\), 有
\begin{equation}
	\int_{\R^d}f(Ax)\,\dd x
	=\frac{1}{|\det A|}\int_{\R^d}f(x)\,\dd x.
	\label{eq:linear-change-c0}
\end{equation}
更一般地, 对仿射变换 \(\phi(x)=Ax+x_0\), 有
\begin{equation}
	|\det A|\int_{\R^d}\phi^*f(x)\,\dd x
	=
	\int_{\R^d}f(x)\,\dd x,
	\qquad f\in\C_0(\R^d).
	\label{eq:affine-change-c0}
\end{equation}

\begin{rem}
	上面用到的分解 \(A=O_1DO_2\) 可由如下事实得到: 正对称矩阵 \(A^TA\) 可以通过正交变换化为正对角矩阵. 即存在正交矩阵 \(O_2\) 和正对角矩阵 \(D^2\), 使得
	\[
	A^TA=O_2^TD^2O_2.
	\]
	令
	\[
	O_1=AO_2^TD^{-1},
	\]
	则 \(O_1\) 也是正交矩阵, 且 \(A=O_1DO_2\).
\end{rem}

公式 \eqref{eq:affine-change-c0} 也适用于具有紧致支集的有界函数的上积分和下积分. 因而, 若 \(u\) 为这类函数, 则 \(\phi^*u\) 黎曼可积当且仅当 \(u\) 黎曼可积. 特别地, 若 \(E\subset\R^d\) 是 Jordan 可测集, 则 \(A(E)\) 也是 Jordan 可测集, 且
\begin{equation}
	m(A(E))=|\det A|m(E).
	\label{eq:det-volume-ratio}
\end{equation}
这严格解释了行列式作为体积比的意义.

若 \(\det A=0\), 公式 \eqref{eq:det-volume-ratio} 对区间 \(E=I\) 仍然成立. 此时 \(A(I)\) 包含在某个超平面中, 通过正交变换可以把这个超平面化为坐标超平面 \(x_d=0\), 所以它的 Jordan 测度为零.

\subsection{一个体积估计}

下面的估计将在处理一般连续可微映射时用到.

设矩阵范数定义为
\[
\|A\|=\max_{x\neq0}\frac{\|Ax\|}{\|x\|}.
\]
若 \(A\) 可逆, 由上面的正交--对角--正交分解可知
\[
|\det A|\leq \|A\|^d.
\]
所以线性像 \(A(I)\) 的体积不超过 \(\|A\|^d m(I)\). 若 \(I\) 是边长为 \(\delta\) 的立方体, 则 \(A(I)\) 被包含在一个边长为 \(\delta\sqrt d\|A\|\) 的立方体中.

\begin{lem}\label{lem:linear-image-neighborhood-estimate}
	设 \(A\) 为一个 \(d\times d\) 实矩阵, \(I\subset\R^d\) 是边长为 \(\delta\) 的立方体. 对 \(\eta>0\), 记
	\[
	A(I)_\eta=\{x\in\R^d:\operatorname{dist}(x,A(I))\leq \eta\}.
	\]
	则
	\begin{equation}
		m^*\bigl(A(I)_\eta\bigr)
		\leq
		\delta^d|\det A|
		+4d\eta\bigl(2\eta+\delta\sqrt{d-1}\|A\|\bigr)^{d-1}.
		\label{eq:linear-image-neighborhood-estimate}
	\end{equation}
\end{lem}

\begin{proof}
	若 \(\det A\neq0\), 则 \(A\) 把 \(I\) 的内点映为 \(A(I)\) 的内点, 因而 \(A(I)\) 的边界包含在 \(A(\partial I)\) 中. 若 \(\det A=0\), 则过 \(I\) 中每一点都有一条直线被 \(A\) 压缩为一点, 这条直线必与 \(\partial I\) 相交, 从而
	\[
	A(I)=A(\partial I).
	\]
	所以, 除去 \(A(I)\) 本身的体积 \(\delta^d|\det A|\), 只需估计 \(A(\partial I)\) 的 \(\eta\)-邻域.
	
	边界 \(\partial I\) 由 \(2d\) 个 \((d-1)\)-维面组成. 取其中一个面 \(I'\). 由正交不变性和平移不变性, 可设 \(A(I')\) 位于超平面 \(x_d=0\) 内. 又 \(A(I')\) 包含在该超平面内一个边长不超过 \(\delta\sqrt{d-1}\|A\|\) 的 \((d-1)\)-维立方体中.
	
	若点 \(x=(x',x_d)\) 到 \(A(I')\) 的距离不超过 \(\eta\), 则 \(|x_d|\leq \eta\), 且 \(x'\) 到 \(A(I')\) 在超平面中的距离不超过 \(\eta\). 因此这些点被包含在一个集合中, 其体积不超过
	\[
	2\eta\bigl(2\eta+\delta\sqrt{d-1}\|A\|\bigr)^{d-1}.
	\]
	将 \(2d\) 个面相加, 得到
	\[
	4d\eta\bigl(2\eta+\delta\sqrt{d-1}\|A\|\bigr)^{d-1}.
	\]
	再加上 \(A(I)\) 自身的体积上界 \(\delta^d|\det A|\), 即得 \eqref{eq:linear-image-neighborhood-estimate}.
\end{proof}

\subsection{连续可微映射下的体积估计}

现在讨论一般的连续可微映射. 这里暂时不要求映射是一一的.

\begin{thm}\label{thm:c1-map-upper-estimate}
	设 \(\Omega_1,\Omega_2\subset\R^d\) 是开集, \(\phi:\Omega_1\to\Omega_2\) 是连续可微映射. 若 \(K\subset\Omega_1\) 为紧集, 则对任意 \(u\in\C_0^+(\Omega_2)\), 有
	\begin{equation}
		\overline{\int_{\R^d}}\chi_{\phi(K)}(y)u(y)\,\dd y
		\leq
		\int_K u(\phi(x))|\det\phi'(x)|\,\dd x.
		\label{eq:c1-map-upper-estimate}
	\end{equation}
	这里
	\[
	\phi'(x)=\left(\frac{\partial\phi_j}{\partial x_k}(x)\right)_{j,k=1}^{d}.
	\]
\end{thm}

\begin{proof}
	给定 \(\varepsilon>0\). 因为 \(K\) 紧且 \(\phi'\) 连续, \(\phi'\) 在 \(K\) 上一致连续. 取足够细的立方体剖分 \(\{I_j\}\) 覆盖 \(K\), 并在每个与 \(K\) 相交的 \(I_j\) 中取一点 \(y_j\in K\cap I_j\). 当 \(x\in I_j\) 时, 由 Taylor 公式和一致连续性, 有
	\[
	\|\phi(x)-\phi(y_j)-\phi'(y_j)(x-y_j)\|
	\leq \varepsilon\delta,
	\]
	其中 \(\delta\) 为立方体边长.
	于是
	\[
	\phi(I_j\cap K)\subset \bigl(\phi(y_j)+\phi'(y_j)(I_j-y_j)\bigr)_{C\varepsilon\delta},
	\]
	这里常数 \(C\) 与维数和 \(K\) 附近的有界性有关.
	
	由引理 \ref{lem:linear-image-neighborhood-estimate}, 当 \(\delta\) 足够小时,
	\[
	m^*(\phi(I_j\cap K))
	\leq
	m(I_j)\bigl(|\det\phi'(y_j)|+C\varepsilon\bigr).
	\]
	在每个 \(I_j\) 中再取 \(y_j\) 使得 \(u(\phi(y_j))\) 接近该小块上的最大值, 则
	\[
	\overline{\int_{\R^d}}\chi_{\phi(K)}(y)u(y)\,\dd y
	\leq
	\sum_j m(I_j)\bigl(|\det\phi'(y_j)|+C\varepsilon\bigr)u(\phi(y_j)).
	\]
	令剖分细度趋于零, 再令 \(\varepsilon\to0\), 右端趋于
	\[
	\int_K u(\phi(x))|\det\phi'(x)|\,\dd x.
	\]
	于是得到 \eqref{eq:c1-map-upper-estimate}.
\end{proof}

\begin{cor}[Morse--Sard 型推论]
	\label{cor:morse-sard-c1-riemann}
	在定理 \ref{thm:c1-map-upper-estimate} 的假设下, 集合
	\[
	\bigl\{\phi(x):x\in K,\ \det\phi'(x)=0\bigr\}
	\]
	的 Jordan 外测度为零.
\end{cor}

\begin{proof}
	记
	\[
	E=\{x\in K:\det\phi'(x)=0\}.
	\]
	对任意 \(n\), 令
	\[
	E_n=\{x\in K:|\det\phi'(x)|<1/n\}.
	\]
	则 \(E\subset E_n\). 由定理 \ref{thm:c1-map-upper-estimate}, 对任意 \(u\in\C_0^+(\Omega_2)\) 有
	\[
	\overline{\int_{\R^d}}\chi_{\phi(E)}u\,\dd y
	\leq
	\overline{\int_{\R^d}}\chi_{\phi(E_n)}u\,\dd y
	\leq
	\int_{E_n}u(\phi(x))|\det\phi'(x)|\,\dd x.
	\]
	由于 \(u\) 有界, 右端不超过
	\[
	\frac{\|u\|_\infty}{n}m(K).
	\]
	令 \(n\to\infty\), 得
	\[
	\overline{\int_{\R^d}}\chi_{\phi(E)}u\,\dd y=0.
	\]
	取 \(u\) 在 \(\phi(K)\) 的一个邻域上恒等于 \(1\), 即得 \(m^*(\phi(E))=0\).
\end{proof}

\begin{rem}
	推论 \ref{cor:morse-sard-c1-riemann} 说明, 在一个紧集上, 雅可比行列式为零的点所产生的像集体积为零. 因此, 对绝大多数 \(y\in\Omega_2\), 若 \(y=\phi(x)\), 则在相应点处 \(\det\phi'(x)\neq0\), 可以使用反函数定理.
\end{rem}

\subsection{变量替换公式}

现在可以证明多重积分的变量替换公式.

\begin{thm}[变量替换公式]
	\label{thm:change-of-variables-riemann}
	设 \(\Omega_1,\Omega_2\subset\R^d\) 是开集, \(\phi:\Omega_1\to\Omega_2\) 是一一映射, 且 \(\phi\) 与 \(\phi^{-1}\) 都连续可微. 若 \(u\) 是 \(\Omega_2\) 上的有界函数, 且在 \(\Omega_2\) 的某个紧子集之外为零, 则 \(u\) 黎曼可积的充分必要条件是
	\[
	(\phi^*u)|\det\phi'|
	\]
	在 \(\Omega_1\) 上黎曼可积. 此时
	\begin{equation}
		\int_{\Omega_2}u(y)\,\dd y
		=
		\int_{\Omega_1}u(\phi(x))|\det\phi'(x)|\,\dd x.
		\label{eq:change-of-variables-riemann}
	\end{equation}
\end{thm}

\begin{proof}
	先设 \(u\in\C_0^+(\Omega_2)\). 取紧集 \(K\subset\Omega_1\), 使得在 \(\phi(K)\) 之外有 \(u=0\). 对 \(\phi\) 应用定理 \ref{thm:c1-map-upper-estimate}, 得
	\[
	\int_{\Omega_2}u(y)\,\dd y
	\leq
	\int_{\Omega_1}u(\phi(x))|\det\phi'(x)|\,\dd x.
	\]
	再对 \(\phi^{-1}\) 应用同一结论, 并注意
	\[
	\det (\phi^{-1})'(\phi(x))\cdot \det\phi'(x)=1,
	\]
	得到反向不等式. 因而当 \(u\in\C_0^+(\Omega_2)\) 时, \eqref{eq:change-of-variables-riemann} 成立.
	
	由线性性, 公式对所有 \(u\in\C_0(\Omega_2)\) 成立. 进一步, 与前面由连续函数逼近黎曼可积函数的论证相同, 该公式对有界且具有紧致支集的函数的上积分和下积分都成立. 因此, \(u\) 黎曼可积当且仅当 \((\phi^*u)|\det\phi'|\) 黎曼可积, 并且此时公式 \eqref{eq:change-of-variables-riemann} 成立.
\end{proof}

\begin{rem}
	公式 \eqref{eq:change-of-variables-riemann} 中的因子 \(|\det\phi'(x)|\) 可以理解为 \(\phi\) 在点 \(x\) 附近的局部体积伸缩率. 仿射情形中它退化为常数 \(|\det A|\), 而一般情形则是把每一点处的线性近似累积起来.
\end{rem}

\subsection{关于非紧支集函数的补充}

若希望把本节的变量替换公式推广到不一定具有紧致支集的函数, 可以按如下方式扩展黎曼积分.

设 \(\Omega\subset\R^d\) 为开集. 若 \(f\geq0\) 是 \(\Omega\) 上的连续函数, 定义
\[
\int_\Omega f(x)\,\dd x
=
\sup\left\{
\int_\Omega g(x)\,\dd x:
0\leq g\leq f,
\ g\in\C_0(\Omega)
\right\}.
\]
这个值允许为 \(+\infty\). 当 \(f\in\C_0(\Omega)\) 时, 这个定义与前面的定义一致.

若 \(f\geq0\) 在 \(\Omega\) 上局部有界, 即在每个紧子集上有界, 则可以定义它的上积分为
\[
\overline{\int_\Omega}f(x)\,\dd x
=
\inf\left\{
\int_\Omega g(x)\,\dd x:
f\leq g,
\ g\geq0,
\ g\text{连续}
\right\}.
\]
最后, 对一般局部有界函数 \(f\), 若对任意 \(\varepsilon>0\), 存在 \(g_\varepsilon\in\C_0(\Omega)\), 使得
\[
\overline{\int_\Omega}|f(x)-g_\varepsilon(x)|\,\dd x<\varepsilon,
\]
则称 \(f\) 在 \(\Omega\) 上黎曼可积, 并把
\[
\int_\Omega f(x)\,\dd x
:=
\lim_{\varepsilon\to0}
\int_\Omega g_\varepsilon(x)\,\dd x
\]
作为它的积分. 这里极限与逼近函数的选取无关.

在这个扩展意义下, 定理 \ref{thm:change-of-variables-riemann} 可以推广到任意局部有界且可积的函数. 本书下一章的重点是勒贝格积分, 因而这里不再展开这些细节.
	
	\chapter{勒贝格积分}

\section{下方半连续正函数的积分}
\label{sec:lsc-positive-integral}

在第一章中, 我们已经把黎曼积分从有限区间推广到 \(\C_0(\R^d)\) 中的函数. 后面将经常只使用下面两个基本性质:
\begin{align}
	\int_{\R^d}\bigl(af(x)+bg(x)\bigr)\,\dd x
	&=a\int_{\R^d}f(x)\,\dd x
	+b\int_{\R^d}g(x)\,\dd x,
	\label{eq:c0-linear-lebesgue-start}\\
	\int_{\R^d}f(x)\,\dd x&\geq 0,
	\qquad f\in \C_0^+(\R^d),
	\label{eq:c0-positive-lebesgue-start}
\end{align}
其中 \(a,b\in\R\), \(f,g\in\C_0(\R^d)\). 由这两个性质又立即得到单调性: 若 \(f,g\in\C_0(\R^d)\) 且 \(f\leq g\), 则
\[
\int_{\R^d}f(x)\,\dd x
\leq
\int_{\R^d}g(x)\,\dd x.
\]
请注意, 到目前为止这里并没有用到黎曼积分关于坐标平移的不变性. 这一点在第三章讨论更一般的 Radon 测度时会变得很重要.

本节的目标是把积分从 \(\C_0^+(\R^d)\) 扩展到更大的函数类. 这个函数类正是下方半连续的非负函数.

\begin{defn}[下方半连续函数]
	\label{def:lsc-function}
	设 \(f\) 是定义在 \(\R^d\) 上、取值于 \((-\infty,+\infty]\) 的函数. 若对任意 \(x\in\R^d\) 和任意实数 \(a<f(x)\), 都存在 \(x\) 的一个邻域 \(U\), 使得
	\[
	f(y)>a,
	\qquad y\in U,
	\]
	则称 \(f\) 为下方半连续函数. 若 \(-f\) 为下方半连续函数, 则称 \(f\) 为上方半连续函数.
	
	所有非负下方半连续函数构成的集合记作 \(\mathcal I^+\).
\end{defn}

这里允许函数取值 \(+\infty\). 以后采用通常约定:
\[
a<+\infty \quad (a\in\R),
\]
并且
\[
(+\infty)+a=+\infty,
\qquad
(+\infty)+(+\infty)=+\infty,
\]
其中 \(-\infty<a\leq+\infty\). 若 \(0<a<+\infty\), 则
\[
a\cdot(+\infty)=(+\infty)\cdot a=+\infty.
\]
允许 \(+\infty\) 作为函数值的一个好处是: 非空实数集总有上确界, 这样在处理极限和上确界时记号会更自然.

由开集的定义可知, 定义 \ref{def:lsc-function} 等价于说: 对任意实数 \(a\), 集合
\[
\{x\in\R^d:f(x)>a\}
\]
都是开集.

\begin{lem}\label{lem:lsc-stability}
	设 \(\{f_\alpha\}_{\alpha\in A}\) 是一族下方半连续函数. 则
	\[
	\sup_{\alpha\in A}f_\alpha
	\]
	仍是下方半连续函数. 若 \(A\) 为有限集, 则
	\[
	\inf_{\alpha\in A}f_\alpha
	\]
	也是下方半连续函数.
\end{lem}

\begin{proof}
	若
	\[
	F(x)=\sup_{\alpha\in A}f_\alpha(x),
	\]
	则
	\[
	\{x:F(x)>a\}
	=\bigcup_{\alpha\in A}\{x:f_\alpha(x)>a\},
	\]
	右端是开集的并, 因而是开集. 所以 \(F\) 下方半连续.
	
	若 \(A=\{1,\ldots,N\}\), 且
	\[
	G(x)=\min_{1\leq j\leq N}f_j(x),
	\]
	则
	\[
	\{x:G(x)>a\}
	=\bigcap_{j=1}^{N}\{x:f_j(x)>a\},
	\]
	右端是有限个开集的交, 因而是开集. 所以 \(G\) 下方半连续.
\end{proof}

下面的定理说明, \(\mathcal I^+\) 中函数正好可以看成 \(\C_0^+(\R^d)\) 中函数的非负级数和.

\begin{thm}\label{thm:lsc-as-sum-c0-positive}
	设 \(f:\R^d\to[0,+\infty]\). 则 \(f\in\mathcal I^+\) 的充分必要条件是存在 \(u_n\in\C_0^+(\R^d)\), 使得
	\begin{equation}
		f(x)=\sum_{n=1}^{\infty}u_n(x),
		\qquad x\in\R^d.
		\label{eq:lsc-sum-c0-positive}
	\end{equation}
\end{thm}

\begin{proof}
	若 \(f\) 具有表示式 \eqref{eq:lsc-sum-c0-positive}, 则它是部分和
	\[
	S_N(x)=\sum_{n=1}^{N}u_n(x)
	\]
	的上确界. 每个 \(S_N\) 连续, 因而下方半连续. 由引理 \ref{lem:lsc-stability}, \(f=\sup_N S_N\) 下方半连续.
	
	反过来, 设 \(f\in\mathcal I^+\). 我们先构造一列非负连续函数, 它单调上升到 \(f\). 对 \(n=1,2,\ldots\), 定义
	\begin{equation}
		s_n(x)=\min\left\{n,\inf_{z\in\R^d}\bigl(f(z)+n\|z-x\|\bigr)\right\}.
		\label{eq:lsc-inf-convolution}
	\end{equation}
	这里加入外层的 \(\min\{n,\cdot\}\), 是为了允许 \(f\) 在某些点甚至所有点取值为 \(+\infty\). 显然
	\[
	0\leq s_n(x)\leq f(x),
	\qquad x\in\R^d.
	\]
	并且 \(s_n\) 关于 \(n\) 单调递增.
	
	又因为对任意 \(x,y,z\in\R^d\), 有
	\[
	\bigl|\|z-x\|-\|z-y\|\bigr|\leq \|x-y\|,
	\]
	所以由 \eqref{eq:lsc-inf-convolution} 可得
	\[
	|s_n(x)-s_n(y)|\leq n\|x-y\|.
	\]
	因此 \(s_n\) 是连续函数.
	
	下面证明 \(s_n(x)\uparrow f(x)\). 若 \(f(x)<+\infty\), 任取 \(a<f(x)\). 由下方半连续性, 存在 \(x\) 的邻域 \(U\), 使得当 \(z\in U\) 时
	\[
	f(z)>a.
	\]
	另一方面, 当 \(z\notin U\) 时, \(\|z-x\|\) 有正下界. 因而当 \(n\) 足够大时,
	\[
	n\|z-x\|>a,
	\qquad z\notin U.
	\]
	于是对所有 \(z\in\R^d\) 都有
	\[
	f(z)+n\|z-x\|>a,
	\]
	从而 \(s_n(x)>a\) 对充分大的 \(n\) 成立. 由 \(a<f(x)\) 的任意性, 得 \(s_n(x)\to f(x)\).
	
	若 \(f(x)=+\infty\), 对任意 \(A>0\), 同样由下方半连续性可取 \(x\) 的邻域 \(U\), 使得 \(f(z)>A\) 于 \(U\) 内. 再令 \(n\) 足够大, 使得 \(n>A\), 且 \(n\|z-x\|>A\) 对 \(z\notin U\) 成立. 于是 \(s_n(x)>A\). 因 \(A\) 任意, 得 \(s_n(x)\to+\infty=f(x)\).
	
	现在还要把这些连续函数改造成具有紧致支集的函数. 取一个函数 \(\theta\in\C_0^+(\R^d)\), 满足
	\[
	0\leq\theta\leq1,
	\qquad
	\theta(x)=1\quad (\|x\|\leq1),
	\qquad
	\theta(x)=0\quad (\|x\|\geq2),
	\]
	并且 \(\theta(x)\) 只依赖于 \(\|x\|\), 且随 \(\|x\|\) 单调下降. 令
	\[
	\widetilde S_n(x)=s_n(x)\theta\left(\frac{x}{n}\right).
	\]
	由于 \(s_n\geq0\) 关于 \(n\) 单调递增, 而 \(\theta(x/n)\) 也关于 \(n\) 单调递增并趋于 \(1\), 所以 \(\widetilde S_n\) 关于 \(n\) 单调递增, 且
	\[
	\widetilde S_n(x)\uparrow f(x).
	\]
	同时 \(\widetilde S_n\in\C_0^+(\R^d)\). 令 \(\widetilde S_0=0\), 并定义
	\[
	u_n=\widetilde S_n-\widetilde S_{n-1},
	\qquad n=1,2,\ldots.
	\]
	则 \(u_n\in\C_0^+(\R^d)\), 且
	\[
	\sum_{n=1}^{\infty}u_n(x)
	=
	\lim_{N\to\infty}\widetilde S_N(x)
	=f(x).
	\]
	定理得证.
\end{proof}

\begin{exer}
	在 \(f\) 处处有限的情形, 证明 \eqref{eq:lsc-inf-convolution} 中不加外层截断时, 下确界可以实际达到. 即证明对每个 \(x\in\R^d\) 和每个 \(n\), 存在 \(z_x\in\R^d\), 使得
	\[
	\inf_{z\in\R^d}\bigl(f(z)+n\|z-x\|\bigr)
	=f(z_x)+n\|z_x-x\|.
	\]
\end{exer}

\subsection{下方半连续正函数的上积分}

根据定理 \ref{thm:lsc-as-sum-c0-positive}, 若 \(f\in\mathcal I^+\), 它可以写成 \(\C_0^+(\R^d)\) 中函数的级数. 定理 \ref{thm:positive-series-riemann} 的证明稍加修改便说明: 对任意这样的表示式, 积分和都由 \(f\) 本身决定. 因此自然引入如下定义.

\begin{defn}[上积分]
	\label{def:upper-integral-lsc-positive}
	设 \(f\in\mathcal I^+\). 定义 \(f\) 的上积分为
	\begin{equation}
		\int_{\R^d}^{*}f(x)\,\dd x
		=
		\sup\left\{
		\int_{\R^d}g(x)\,\dd x:
		g\in\C_0(\R^d),\ 0\leq g\leq f
		\right\}.
		\label{eq:upper-integral-lsc-positive}
	\end{equation}
\end{defn}

上式右端允许取值为 \(+\infty\). 事实上, 只让 \(g\) 取遍 \(\C_0^+(\R^d)\) 即可. 以后会看到, 如果 \(f\in\mathcal I^+\) 且
\[
\int_{\R^d}^{*}f(x)\,\dd x<+\infty,
\]
那么 \(f\) 将是勒贝格可积的. 因此在这种情形下, 这个“上积分”实际上就是积分.

下面的引理是处理上积分的基本工具.

\begin{lem}\label{lem:directed-sup-integral}
	设 \(M\subset \C_0^+(\R^d)\), 并假定 \(M\) 关于向上取大是有向的: 对任意 \(g_1,g_2\in M\), 存在 \(g\in M\), 使得
	\[
	g_1\leq g,
	\qquad
	g_2\leq g.
	\]
	令
	\[
	f(x)=\sup_{g\in M}g(x).
	\]
	则 \(f\in\mathcal I^+\), 且
	\begin{equation}
		\int_{\R^d}^{*}f(x)\,\dd x
		=
		\sup_{g\in M}\int_{\R^d}g(x)\,\dd x.
		\label{eq:directed-sup-integral}
	\end{equation}
\end{lem}

\begin{proof}
	由引理 \ref{lem:lsc-stability}, \(f=\sup_{g\in M}g\) 下方半连续, 所以 \(f\in\mathcal I^+\). 公式 \eqref{eq:directed-sup-integral} 的右端显然不超过左端.
	
	为证反向不等式, 任取 \(h\in\C_0^+(\R^d)\), 且 \(h\leq f\). 再取 \(\eta\in\C_0^+(\R^d)\), 使得 \(\eta>0\) 于 \(\operatorname{supp}h\) 上. 给定 \(\varepsilon>0\). 对每个 \(x\in\operatorname{supp}h\), 因 \(h(x)\leq f(x)=\sup_{g\in M}g(x)\), 可取 \(g_x\in M\), 使得
	\[
	g_x(x)>h(x)-\varepsilon\eta(x).
	\]
	由连续性, 此不等式在 \(x\) 的某个邻域 \(U_x\) 内仍成立. 由于 \(\operatorname{supp}h\) 紧, 可取有限个点 \(x_1,\ldots,x_N\), 使得这些邻域覆盖 \(\operatorname{supp}h\). 根据 \(M\) 的有向性, 存在 \(G\in M\), 使得
	\[
	G\geq g_{x_j},
	\qquad j=1,\ldots,N.
	\]
	于是
	\[
	G\geq h-\varepsilon\eta
	\]
	在 \(\operatorname{supp}h\) 上成立, 在 \(\operatorname{supp}h\) 外也有 \(h=0\), 因而
	\[
	\int_{\R^d}G(x)\,\dd x
	\geq
	\int_{\R^d}h(x)\,\dd x
	-
	\varepsilon\int_{\R^d}\eta(x)\,\dd x.
	\]
	令 \(\varepsilon\to0\), 得
	\[
	\sup_{g\in M}\int_{\R^d}g(x)\,\dd x
	\geq
	\int_{\R^d}h(x)\,\dd x.
	\]
	再对所有 \(0\leq h\in\C_0(\R^d)\), \(h\leq f\) 取上确界, 即得反向不等式.
\end{proof}

特别地, 若 \(f_n\in\C_0^+(\R^d)\) 且 \(f_n\uparrow f\), 则
\[
\int_{\R^d}^{*}f(x)\,\dd x
=
\lim_{n\to\infty}\int_{\R^d}f_n(x)\,\dd x.
\]

\begin{thm}\label{thm:basic-properties-upper-integral-lsc}
	上积分具有如下性质:
	\begin{enumerate}
		\item 若 \(f,g\in\mathcal I^+\) 且 \(f\leq g\), 则
		\[
		\int_{\R^d}^{*}f(x)\,\dd x
		\leq
		\int_{\R^d}^{*}g(x)\,\dd x.
		\]
		\item 若 \(f\in\mathcal I^+\), \(a>0\), 则
		\[
		\int_{\R^d}^{*}af(x)\,\dd x
		=
		a\int_{\R^d}^{*}f(x)\,\dd x.
		\]
		\item 若 \(f_n\in\mathcal I^+\), 则 \(\sum_{n=1}^{\infty}f_n\in\mathcal I^+\), 且
		\begin{equation}
			\int_{\R^d}^{*}\left(\sum_{n=1}^{\infty}f_n(x)\right)\dd x
			=
			\sum_{n=1}^{\infty}\int_{\R^d}^{*}f_n(x)\,\dd x.
			\label{eq:upper-integral-countable-additivity-lsc}
		\end{equation}
		\item 若 \(f_n\in\mathcal I^+\) 且 \(f_n\uparrow f\), 则 \(f\in\mathcal I^+\), 且
		\begin{equation}
			\int_{\R^d}^{*}f(x)\,\dd x
			=
			\lim_{n\to\infty}\int_{\R^d}^{*}f_n(x)\,\dd x.
			\label{eq:upper-integral-monotone-lsc}
		\end{equation}
	\end{enumerate}
\end{thm}

\begin{proof}
	前两个结论由定义立即得到. 下面证明 \eqref{eq:upper-integral-countable-additivity-lsc}. 对每个 \(n\), 由定理 \ref{thm:lsc-as-sum-c0-positive}, 可取 \(u_{n,k}\in\C_0^+(\R^d)\), 使得
	\[
	f_n=\sum_{k=1}^{\infty}u_{n,k},
	\qquad
	\int_{\R^d}^{*}f_n\,\dd x
	=
	\sum_{k=1}^{\infty}\int_{\R^d}u_{n,k}\,\dd x.
	\]
	将双重序列 \(\{u_{n,k}\}\) 重新排成一个单序列, 并应用引理 \ref{lem:directed-sup-integral}, 即得
	\[
	\int_{\R^d}^{*}\left(\sum_{n=1}^{\infty}f_n\right)\dd x
	=
	\sum_{n=1}^{\infty}\int_{\R^d}^{*}f_n\,\dd x.
	\]
	
	最后, 若 \(f_n\uparrow f\), 则由引理 \ref{lem:lsc-stability} 可知 \(f\in\mathcal I^+\). 令
	\[
	M=\{g\in\C_0^+(\R^d):\text{存在某个 }n\text{ 使 }g\leq f_n\}.
	\]
	这个集合是向上有向的, 且 \(\sup_{g\in M}g=f\). 由引理 \ref{lem:directed-sup-integral} 得
	\[
	\int_{\R^d}^{*}f\,\dd x
	=\sup_n\int_{\R^d}^{*}f_n\,\dd x,
	\]
	这正是 \eqref{eq:upper-integral-monotone-lsc}.
\end{proof}

\clearpage

\section{正函数的上积分}
\label{sec:upper-integral-positive-functions}

上一节已经把积分定义到了下方半连续的非负函数类 \(\mathcal I^+\) 上. 现在我们进一步把上积分定义到任意非负函数上. 这一步与黎曼积分中用分片常数函数或连续函数去夹逼一个函数的思想相似; 不同的是, 这里用来控制函数的是 \(\mathcal I^+\) 中的函数.

\begin{defn}[非负函数的上积分]
	\label{def:upper-integral-positive-functions}
	设 \(f:\R^d\to[0,+\infty]\) 为任意非负函数. 定义 \(f\) 的上积分为
	\begin{equation}
		\int_{\R^d}^{*}f(x)\,\dd x
		=
		\inf\left\{
		\int_{\R^d}^{*}h(x)\,\dd x:
		h\in\mathcal I^+,\ f\leq h
		\right\}.
		\label{eq:upper-integral-positive-functions}
	\end{equation}
	这里允许右端取值为 \(+\infty\).
\end{defn}

若 \(f\in\mathcal I^+\), 则定义 \ref{def:upper-integral-positive-functions} 与上一节的定义一致. 事实上, 因为可以取 \(h=f\), 由单调性立刻得到两种定义相同.

还要注意, 若 \(f\geq0\) 有界且具有紧致支集, 则
\begin{equation}
	\int_{\R^d}^{*}f(x)\,\dd x
	\leq
	\overline{\int_{\R^d}}f(x)\,\dd x.
	\label{eq:new-upper-integral-less-riemann-upper}
\end{equation}
这是因为黎曼上积分可以用连续函数从上方逼近来定义, 而连续函数属于 \(\mathcal I^+\). 因此, 现在的上积分是在更大的上方控制函数类中取下确界, 所以不会比黎曼上积分更大.

下面的引理用于处理单调极限中出现``非递增上控制函数".

\begin{lem}
	\label{lem:upper-integral-max-majorant}
	设 \(0\leq f_i\leq g_i\), \(g_i\in\mathcal I^+\), 且
	\[
	\int_{\R^d}^{*}g_i(x)\,\dd x
	<
	\int_{\R^d}^{*}f_i(x)\,\dd x+\varepsilon_i,
	\qquad i=1,2.
	\]
	若 \(f_1\leq f_2\), 则
	\[
	\int_{\R^d}^{*}\max(g_1,g_2)(x)\,\dd x
	<
	\int_{\R^d}^{*}f_2(x)\,\dd x+
	\varepsilon_1+
	\varepsilon_2.
	\]
\end{lem}

\begin{proof}
	因为 \(g_1,g_2\in\mathcal I^+\), 所以 \(\max(g_1,g_2)\) 和 \(\min(g_1,g_2)\) 都属于 \(\mathcal I^+\). 又有恒等式
	\[
	\max(g_1,g_2)+\min(g_1,g_2)=g_1+g_2.
	\]
	由上一节关于 \(\mathcal I^+\) 中函数积分的可加性, 得
	\begin{align*}
		&\int_{\R^d}^{*}\max(g_1,g_2)(x)\,\dd x
		+
		\int_{\R^d}^{*}\min(g_1,g_2)(x)\,\dd x\\
		&\qquad=
		\int_{\R^d}^{*}g_1(x)\,\dd x
		+
		\int_{\R^d}^{*}g_2(x)\,\dd x\\
		&\qquad<
		\int_{\R^d}^{*}f_1(x)\,\dd x
		+
		\int_{\R^d}^{*}f_2(x)\,\dd x
		+
		\varepsilon_1+\varepsilon_2.
	\end{align*}
	由于 \(f_1\leq f_2\), 有
	\[
	f_1=\min(f_1,f_2)\leq \min(g_1,g_2).
	\]
	因此
	\[
	\int_{\R^d}^{*}f_1(x)\,\dd x
	\leq
	\int_{\R^d}^{*}\min(g_1,g_2)(x)\,\dd x.
	\]
	代入上式并消去公共项, 得
	\[
	\int_{\R^d}^{*}\max(g_1,g_2)(x)\,\dd x
	<
	\int_{\R^d}^{*}f_2(x)\,\dd x+
	\varepsilon_1+
	\varepsilon_2.
	\]
\end{proof}

\begin{thm}[上积分的基本性质]
	\label{thm:basic-properties-upper-integral-positive}
	设 \(f,g,f_n\) 都是 \(\R^d\) 上的非负函数. 则有:
	\begin{enumerate}
		\item 若 \(0\leq f\leq g\), 则
		\[
		\int_{\R^d}^{*}f(x)\,\dd x
		\leq
		\int_{\R^d}^{*}g(x)\,\dd x.
		\]
		\item 若 \(a>0\), 则
		\[
		\int_{\R^d}^{*}af(x)\,\dd x
		=
		a\int_{\R^d}^{*}f(x)\,\dd x.
		\]
		\item 对任意非负函数列 \(\{f_n\}\), 有
		\begin{equation}
			\int_{\R^d}^{*}\left(\sum_{n=1}^{\infty}f_n(x)\right)\dd x
			\leq
			\sum_{n=1}^{\infty}\int_{\R^d}^{*}f_n(x)\,\dd x.
			\label{eq:upper-integral-subadditivity-positive}
		\end{equation}
		\item 若 \(0\leq f_n\uparrow f\), 即 \(f_n(x)\uparrow f(x)\) 对每个 \(x\) 成立, 则
		\begin{equation}
			\int_{\R^d}^{*}f_n(x)\,\dd x
			\longrightarrow
			\int_{\R^d}^{*}f(x)\,\dd x.
			\label{eq:upper-integral-monotone-positive}
		\end{equation}
	\end{enumerate}
\end{thm}

\begin{proof}
	前两个结论由定义直接得到.
	
	下面证明 \eqref{eq:upper-integral-subadditivity-positive}. 给定 \(\varepsilon>0\). 对每个 \(n\), 由定义可取 \(g_n\in\mathcal I^+\), 使得
	\[
	f_n\leq g_n,
	\qquad
	\int_{\R^d}^{*}g_n(x)\,\dd x
	<
	\int_{\R^d}^{*}f_n(x)\,\dd x+\frac{\varepsilon}{2^n}.
	\]
	令
	\[
	g(x)=\sum_{n=1}^{\infty}g_n(x).
	\]
	由上一节的定理, \(\sum_{n=1}^{\infty}g_n\in\mathcal I^+\), 且
	\[
	\int_{\R^d}^{*}g(x)\,\dd x
	=
	\sum_{n=1}^{\infty}\int_{\R^d}^{*}g_n(x)\,\dd x.
	\]
	又因为
	\[
	\sum_{n=1}^{\infty}f_n(x)\leq g(x),
	\]
	由定义可得
	\begin{align*}
		\int_{\R^d}^{*}\left(\sum_{n=1}^{\infty}f_n(x)\right)\dd x
		&\leq
		\int_{\R^d}^{*}g(x)\,\dd x\\
		&=
		\sum_{n=1}^{\infty}\int_{\R^d}^{*}g_n(x)\,\dd x\\
		&<
		\sum_{n=1}^{\infty}\int_{\R^d}^{*}f_n(x)\,\dd x+\varepsilon.
	\end{align*}
	令 \(\varepsilon\to0\), 得到 \eqref{eq:upper-integral-subadditivity-positive}.
	
	最后证明单调收敛性质 \eqref{eq:upper-integral-monotone-positive}. 显然, 由 \(f_n\leq f\) 可知
	\[
	\int_{\R^d}^{*}f_n(x)\,\dd x
	\leq
	\int_{\R^d}^{*}f(x)\,\dd x,
	\]
	所以只需证明反向估计.
	
	给定 \(\varepsilon>0\). 对每个 \(n\), 取 \(g_n\in\mathcal I^+\), 使得
	\[
	f_n\leq g_n,
	\qquad
	\int_{\R^d}^{*}g_n(x)\,\dd x
	<
	\int_{\R^d}^{*}f_n(x)\,\dd x+\frac{\varepsilon}{2^n}.
	\]
	由于 \(g_n\) 不一定是递增的, 令
	\[
	G_n=\max(g_1,\ldots,g_n).
	\]
	反复应用引理 \ref{lem:upper-integral-max-majorant}, 得
	\begin{equation}
		\int_{\R^d}^{*}G_n(x)\,\dd x
		<
		\int_{\R^d}^{*}f_n(x)\,\dd x+
		\varepsilon.
		\label{eq:Gn-majorant-estimate}
	\end{equation}
	显然 \(G_n\in\mathcal I^+\), 且 \(G_n\uparrow G\in\mathcal I^+\). 又因为 \(f_n\leq G_n\), 对每个 \(x\) 令 \(n\to\infty\), 得
	\[
	f(x)\leq G(x).
	\]
	于是由定义和上一节的单调收敛性质,
	\begin{align*}
		\int_{\R^d}^{*}f(x)\,\dd x
		&\leq
		\int_{\R^d}^{*}G(x)\,\dd x\\
		&=
		\lim_{n\to\infty}
		\int_{\R^d}^{*}G_n(x)\,\dd x\\
		&\leq
		\lim_{n\to\infty}
		\int_{\R^d}^{*}f_n(x)\,\dd x+
		\varepsilon.
	\end{align*}
	因为 \(\varepsilon>0\) 任意, 得
	\[
	\int_{\R^d}^{*}f(x)\,\dd x
	\leq
	\lim_{n\to\infty}
	\int_{\R^d}^{*}f_n(x)\,\dd x.
	\]
	反向不等式已经由单调性得到. 所以 \eqref{eq:upper-integral-monotone-positive} 成立.
\end{proof}

\begin{rem}
	不等式 \eqref{eq:upper-integral-subadditivity-positive} 和单调收敛性质 \eqref{eq:upper-integral-monotone-positive} 是黎曼上积分所不具备的. 这正是本节引入的上积分比黎曼上积分更适合处理极限过程的地方.
\end{rem}

\clearpage

\section{零集合和零函数}
\label{sec:null-sets-and-null-functions}

在本节中, 我们研究这样一个问题: 在哪些点集上改变一个函数的取值, 不会改变它的积分? 这将引出零函数、零集合以及``几乎处处''这些基本概念.

\begin{defn}[零函数]
	\label{def:null-function}
	设 \(f\geq0\). 若
	\[
	\int_{\R^d}^{*}f(x)\,\dd x=0,
	\]
	则称 \(f\) 为零函数.
\end{defn}

由上积分的单调性与可数次可加性, 立刻得到下面的结论.

\begin{lem}\label{lem:null-function-basic}
	若 \(0\leq f\leq g\), 且 \(g\) 是零函数, 则 \(f\) 也是零函数. 若 \(f_1,f_2,\ldots\) 都是零函数, 则
	\[
	\sum_{n=1}^{\infty}f_n
	\]
	也是零函数.
\end{lem}

\begin{proof}
	第一部分由单调性直接得到. 第二部分由上积分的次可加性得到:
	\[
	\int_{\R^d}^{*}\left(\sum_{n=1}^{\infty}f_n(x)\right)\dd x
	\leq
	\sum_{n=1}^{\infty}\int_{\R^d}^{*}f_n(x)\,\dd x=0.
	\]
	因此 \(\sum_{n=1}^{\infty}f_n\) 是零函数.
\end{proof}

\begin{thm}\label{thm:null-function-supported-on-null-function}
	设 \(f,g\geq0\), 且 \(g\) 是零函数. 若
	\[
	f(x)\neq0\quad \Longrightarrow\quad g(x)\neq0,
	\]
	则 \(f\) 是零函数.
\end{thm}

\begin{proof}
	由引理 \ref{lem:null-function-basic}, 对每个 \(n\), 函数 \(ng\) 是零函数. 因而
	\[
	G(x)=\sum_{n=1}^{\infty}g(x)
	\]
	也是零函数. 但当 \(g(x)>0\) 时, 有 \(G(x)=+\infty\); 当 \(g(x)=0\) 时, 有 \(G(x)=0\). 由假设, 若 \(f(x)\neq0\), 则 \(g(x)\neq0\), 从而 \(G(x)=+\infty\geq f(x)\); 若 \(f(x)=0\), 则当然有 \(G(x)\geq f(x)\). 因此
	\[
	0\leq f\leq G.
	\]
	再由引理 \ref{lem:null-function-basic}, \(f\) 是零函数.
\end{proof}

定理 \ref{thm:null-function-supported-on-null-function} 表明, 判断一个非负函数是否是零函数, 只需要看它在哪些点上不为零. 因此自然引入零集合的概念.

\begin{defn}[零集合]
	\label{def:null-set}
	设 \(E\subset\R^d\). 若 \(E\) 的特征函数 \(\chi_E\) 是零函数, 即
	\[
	\int_{\R^d}^{*}\chi_E(x)\,\dd x=0,
	\]
	则称 \(E\) 为零集合.
\end{defn}

\begin{thm}\label{thm:null-set-basic}
	零集合具有如下性质:
	\begin{enumerate}
		\item 零集合的任意子集仍是零集合;
		\item 可数多个零集合的并仍是零集合;
		\item 非负函数 \(f\) 是零函数的充分必要条件是集合
		\[
		\{x:f(x)\neq0\}
		\]
		是零集合.
	\end{enumerate}
\end{thm}

\begin{proof}
	若 \(A\subset E\) 且 \(E\) 是零集合, 则
	\[
	0\leq \chi_A\leq \chi_E,
	\]
	所以 \(A\) 是零集合.
	
	若 \(E_n\) 都是零集合, 则
	\[
	\chi_{\bigcup_{n=1}^{\infty}E_n}
	\leq
	\sum_{n=1}^{\infty}\chi_{E_n}.
	\]
	右端是零函数, 故 \(\bigcup_{n=1}^{\infty}E_n\) 是零集合.
	
	最后, 若 \(f\) 是零函数, 令
	\[
	E=\{x:f(x)\neq0\}.
	\]
	由定理 \ref{thm:null-function-supported-on-null-function}, 取 \(g=f\) 可知 \(\chi_E\) 是零函数, 所以 \(E\) 是零集合. 反过来, 若 \(E=\{x:f(x)\neq0\}\) 是零集合, 则
	\[
	f(x)\neq0\quad\Longrightarrow\quad \chi_E(x)\neq0.
	\]
	由定理 \ref{thm:null-function-supported-on-null-function}, \(f\) 是零函数.
\end{proof}

\begin{rem}
	每个点都是零集合. 因而, 由定理 \ref{thm:null-set-basic}, 任意可数集合都是零集合.
	
	若一个有界集合的 Jordan 外测度为零, 则它也是这里意义下的零集合. 但是反过来不一定成立: 有界可数集一定是零集合, 但它的 Jordan 外测度可以为正. 例如 \([0,1]\cap\mathbb Q\) 是可数集, 因而是零集合; 但它在任意非空子区间中都稠密, 其 Jordan 外测度等于 \(1\).
\end{rem}

\begin{exer}\label{exer:compact-null-set-jordan-measure-zero}
	设 \(K\subset\R^d\) 是紧零集. 证明 \(K\) 的 Jordan 外测度为零.
	
	提示: 若 \(g_n\in\C_0^+(\R^d)\), \(g_n\downarrow g\), 且 \(g\geq\chi_K\), 证明当 \(n\) 充分大时, 有
	\[
	g_n\geq (1-\varepsilon)\chi_K.
	\]
\end{exer}

由上面的练习可以得到一个熟悉的结论: 区间 \([0,1]\) 不可数. 否则 \([0,1]\) 作为可数集将是零集合; 又因为它是紧集, 按练习应有 Jordan 测度为零, 这与 \(m([0,1])=1\) 矛盾.

\begin{defn}[几乎处处]
	\label{def:almost-everywhere}
	设 \(P(x)\) 是一个依赖于 \(x\in\R^d\) 的命题. 若 \(P(x)\) 在除去某个零集合以外的所有点上成立, 则称 \(P(x)\) 几乎处处成立, 简记为 a.e.
\end{defn}

用这个术语, 定理 \ref{thm:null-set-basic} 的第三部分可以表述为: 一个非负函数是零函数, 当且仅当它几乎处处等于零.

\begin{thm}\label{thm:ae-equal-same-upper-integral}
	设 \(f,g\geq0\). 若 \(f=g\) 几乎处处成立, 则
	\[
	\int_{\R^d}^{*}f(x)\,\dd x
	=
	\int_{\R^d}^{*}g(x)\,\dd x.
	\]
\end{thm}

\begin{proof}
	令
	\[
	h=\min(f,g),
	\qquad
	H=\max(f,g).
	\]
	则 \(H-h\) 只可能在 \(f\neq g\) 的点上不为零, 而集合 \(\{x:f(x)\neq g(x)\}\) 是零集合. 因此 \(H-h\) 是零函数. 由上积分的次可加性,
	\[
	\int_{\R^d}^{*}H\,\dd x
	\leq
	\int_{\R^d}^{*}h\,\dd x
	+
	\int_{\R^d}^{*}(H-h)\,\dd x
	=
	\int_{\R^d}^{*}h\,\dd x.
	\]
	又因 \(h\leq H\), 单调性给出反向不等式, 故
	\[
	\int_{\R^d}^{*}H\,\dd x
	=
	\int_{\R^d}^{*}h\,\dd x.
	\]
	由于
	\[
	h\leq f,g\leq H,
	\]
	再次利用单调性, 得
	\[
	\int_{\R^d}^{*}f\,\dd x
	=
	\int_{\R^d}^{*}g\,\dd x.
	\]
\end{proof}

\begin{thm}\label{thm:finite-upper-integral-finite-ae}
	若 \(f\geq0\), 且
	\[
	\int_{\R^d}^{*}f(x)\,\dd x<+\infty,
	\]
	则 \(f\) 几乎处处有限.
\end{thm}

\begin{proof}
	令
	\[
	E=\{x:f(x)=+\infty\}.
	\]
	对任意正整数 \(n\), 有
	\[
	n\chi_E\leq f.
	\]
	由单调性和齐次性,
	\[
	n\int_{\R^d}^{*}\chi_E(x)\,\dd x
	=
	\int_{\R^d}^{*}n\chi_E(x)\,\dd x
	\leq
	\int_{\R^d}^{*}f(x)\,\dd x<+\infty.
	\]
	令 \(n\to\infty\), 只能有
	\[
	\int_{\R^d}^{*}\chi_E(x)\,\dd x=0.
	\]
	因此 \(E\) 是零集合. 即 \(f<+\infty\) 几乎处处成立.
\end{proof}

结合定理 \ref{thm:basic-properties-upper-integral-positive} 的次可加性, 可以得到下面的重要推论.

\begin{cor}\label{cor:summable-integrals-finite-ae}
	设 \(f_n\geq0\). 若
	\[
	\sum_{n=1}^{\infty}\int_{\R^d}^{*}f_n(x)\,\dd x<+\infty,
	\]
	则级数
	\[
	\sum_{n=1}^{\infty}f_n(x)
	\]
	在 \(\R^d\) 上几乎处处收敛到有限值.
\end{cor}

\begin{proof}
	由上积分的次可加性,
	\[
	\int_{\R^d}^{*}\left(\sum_{n=1}^{\infty}f_n(x)\right)\dd x
	\leq
	\sum_{n=1}^{\infty}\int_{\R^d}^{*}f_n(x)\,\dd x<+\infty.
	\]
	于是由定理 \ref{thm:finite-upper-integral-finite-ae}, 函数
	\[
	F(x)=\sum_{n=1}^{\infty}f_n(x)
	\]
	几乎处处有限. 这正说明该非负级数几乎处处收敛.
\end{proof}

\clearpage

\section{勒贝格积分的定义及其重要性质}
\label{sec:definition-properties-lebesgue-integral}

前面已经为任意非负函数定义了上积分. 现在我们开始定义真正的勒贝格可积函数. 思路与第一章中黎曼积分的定义相似: 若一个函数能够被紧支连续函数在上积分意义下任意逼近, 就把它称为可积函数.

\begin{defn}[勒贝格可积函数]
	\label{def:lebesgue-integrable-function}
	设 \(f\) 是取值于 \([-\infty,+\infty]\) 的函数. 若对任意 \(\varepsilon>0\), 存在 \(g\in\C_0(\R^d)\), 使得
	\begin{equation}
		\int_{\R^d}^{*}|f(x)-g(x)|\,\dd x<\varepsilon,
		\label{eq:lebesgue-integrable-definition}
	\end{equation}
	则称 \(f\) 在勒贝格意义下可积, 简称 \(f\) 可积.
\end{defn}

等价地, \(f\) 可积当且仅当对任意 \(\varepsilon>0\), 存在 \(g\in\C_0(\R^d)\) 与 \(h\in\mathcal I^+\), 使得
\[
|f-g|\leq h,
\qquad
\int_{\R^d}^{*}h(x)\,\dd x<\varepsilon.
\]
事实上, 若 \(\int^*|f-g|<\varepsilon\), 由上积分的定义可取 \(h\in\mathcal I^+\) 控制 \(|f-g|\), 并使 \(\int^*h\) 任意接近 \(\int^*|f-g|\); 反向则由单调性立即得到.

\begin{thm}[可积函数的稳定性准则]
	\label{thm:integrable-approx-by-integrable}
	设 \(f\) 为一个函数. 若对任意 \(\varepsilon>0\), 都存在一个可积函数 \(g\), 使得
	\[
	\int_{\R^d}^{*}|f(x)-g(x)|\,\dd x<\varepsilon,
	\]
	则 \(f\) 可积.
\end{thm}

\begin{proof}
	给定 \(\varepsilon>0\). 由假设可取可积函数 \(g\), 使得
	\[
	\int_{\R^d}^{*}|f-g|\,\dd x<\frac{\varepsilon}{2}.
	\]
	又因为 \(g\) 可积, 可取 \(u\in\C_0(\R^d)\), 使得
	\[
	\int_{\R^d}^{*}|g-u|\,\dd x<\frac{\varepsilon}{2}.
	\]
	由上积分的次可加性,
	\[
	\int_{\R^d}^{*}|f-u|\,\dd x
	\leq
	\int_{\R^d}^{*}|f-g|\,\dd x
	+
	\int_{\R^d}^{*}|g-u|\,\dd x
	<\varepsilon.
	\]
	因此 \(f\) 可积.
\end{proof}

\begin{thm}\label{thm:lsc-positive-integrable-iff-finite-upper-integral}
	设 \(f\in\mathcal I^+\). 则 \(f\) 可积的充分必要条件是
	\[
	\int_{\R^d}^{*}f(x)\,\dd x<+\infty.
	\]
\end{thm}

\begin{proof}
	若 \(f\) 可积, 则存在 \(g\in\C_0(\R^d)\), 使得 \(\int^*|f-g|<1\). 于是
	\[
	f\leq |f-g|+|g|,
	\]
	故由次可加性得
	\[
	\int^*f\leq \int^*|f-g|+\int |g|<+\infty.
	\]
	
	反过来, 若 \(f\in\mathcal I^+\) 且 \(\int^*f<+\infty\), 则由上一节下方半连续正函数的逼近定理, 可取 \(g_n\in\C_0^+(\R^d)\), 使得
	\[
	0\leq g_n\uparrow f,
	\qquad
	\int_{\R^d}g_n(x)\,\dd x\uparrow \int_{\R^d}^{*}f(x)\,\dd x.
	\]
	于是
	\[
	\int_{\R^d}^{*}(f-g_n)(x)\,\dd x
	=
	\int_{\R^d}^{*}f(x)\,\dd x-
	\int_{\R^d}g_n(x)\,\dd x\longrightarrow0.
	\]
	因此 \(f\) 可由 \(\C_0\) 中函数按定义 \ref{def:lebesgue-integrable-function} 逼近, 从而可积.
\end{proof}

现在设 \(f\) 可积. 由定义, 可取 \(g_n\in\C_0(\R^d)\), 使得
\begin{equation}
	\int_{\R^d}^{*}|f(x)-g_n(x)|\,\dd x\longrightarrow0.
	\label{eq:gn-approx-f-l1}
\end{equation}
则 \(\int g_n\) 是一个 Cauchy 数列. 事实上,
\[
\left|\int_{\R^d}g_n\,\dd x-
\int_{\R^d}g_m\,\dd x\right|
\leq
\int_{\R^d}^{*}|g_n-g_m|\,\dd x
\leq
\int_{\R^d}^{*}|g_n-f|\,\dd x+
\int_{\R^d}^{*}|f-g_m|\,\dd x.
\]
由 \eqref{eq:gn-approx-f-l1}, 右端趋于 \(0\). 因此 \(\int g_n\) 收敛.

如果 \(h_n\in\C_0(\R^d)\) 是另一列满足 \(\int^*|f-h_n|\to0\) 的函数, 则
\[
\left|\int_{\R^d}g_n\,\dd x-
\int_{\R^d}h_n\,\dd x\right|
\leq
\int_{\R^d}^{*}|g_n-f|\,\dd x+
\int_{\R^d}^{*}|f-h_n|\,\dd x\to0.
\]
所以极限值与逼近列的选取无关.

\begin{defn}[勒贝格积分]
	\label{def:lebesgue-integral}
	设 \(f\) 可积, 并取 \(g_n\in\C_0(\R^d)\), 使得
	\[
	\int_{\R^d}^{*}|f-g_n|\,\dd x\to0.
	\]
	定义
	\[
	\int_{\R^d}f(x)\,\dd x
	=
	\lim_{n\to\infty}\int_{\R^d}g_n(x)\,\dd x.
	\]
\end{defn}

由第二章第三节的结果, 在一个零集合上改变函数值, 不影响函数的可积性, 也不影响积分值. 因此, 若两个可积函数几乎处处相等, 便把它们看成同一个 \(L^1\) 元素. 所有这些等价类组成的集合记作 \(L^1(\R^d)\), 简记为 \(L^1\).

\begin{exer}\label{exer:l1-integral-equals-upper-for-nonnegative}
	设 \(f\in L^1\) 且 \(f\geq0\). 证明
	\[
	\int_{\R^d}f(x)\,\dd x
	=
	\int_{\R^d}^{*}f(x)\,\dd x.
	\]
\end{exer}

\begin{thm}[线性、格运算与单调性]
	\label{thm:basic-properties-l1}
	若 \(f,g\in L^1\), \(a,b\in\R\), 则 \(af+bg\in L^1\), 并且
	\begin{equation}
		\int_{\R^d}(af+bg)(x)\,\dd x
		=a\int_{\R^d}f(x)\,\dd x
		+b\int_{\R^d}g(x)\,\dd x.
		\label{eq:l1-linearity}
	\end{equation}
	此外,
	\[
	\max(f,g),\ \min(f,g),\ |f|\in L^1.
	\]
	若 \(f\leq g\) 几乎处处成立, 则
	\begin{equation}
		\int_{\R^d}f(x)\,\dd x
		\leq
		\int_{\R^d}g(x)\,\dd x.
		\label{eq:l1-monotonicity}
	\end{equation}
\end{thm}

\begin{proof}
	线性性由定义直接得到. 例如, 取 \(f_n,g_n\in\C_0(\R^d)\), 使得
	\[
	\int^*|f-f_n|\to0,
	\qquad
	\int^*|g-g_n|\to0.
	\]
	则 \(af_n+bg_n\in\C_0(\R^d)\), 且由次可加性,
	\[
	\int^*|(af+bg)-(af_n+bg_n)|
	\leq |a|\int^*|f-f_n|+|b|\int^*|g-g_n|\to0.
	\]
	于是 \(af+bg\in L^1\), 并且积分线性性由 \(\C_0\) 中函数积分的线性性传递而来.
	
	下面证明 \(\max(f,g)\in L^1\). 对 \(f_0,g_0\in\C_0(\R^d)\), 有
	\[
	\bigl|\max(f,g)-\max(f_0,g_0)\bigr|
	\leq |f-f_0|+|g-g_0|.
	\]
	因此若 \(f_0\) 与 \(g_0\) 分别逼近 \(f\) 与 \(g\), 则 \(\max(f_0,g_0)\in\C_0(\R^d)\) 逼近 \(\max(f,g)\). 所以 \(\max(f,g)\in L^1\). 类似可证 \(\min(f,g)\in L^1\). 特别地,
	\[
	|f|=\max(f,-f)\in L^1.
	\]
	
	若 \(f\leq g\) 几乎处处成立, 则 \(g-f\geq0\) 且 \(g-f\in L^1\). 由练习 \ref{exer:l1-integral-equals-upper-for-nonnegative},
	\[
	\int_{\R^d}(g-f)\,\dd x\geq0.
	\]
	结合线性性即得 \eqref{eq:l1-monotonicity}.
\end{proof}

\begin{exer}\label{exer:l1-translation-continuity}
	设 \(f\in L^1(\R^d)\). 证明
	\[
	h\longmapsto \int_{\R^d}|f(x+h)-f(x)|\,\dd x
	\]
	是 \(h\) 的连续函数, 并且当 \(h\to0\) 时, 上式趋于 \(0\).
\end{exer}

\begin{exer}\label{exer:linearity-characterizes-l1-positive}
	设 \(f\geq0\), 且 \(\int^* f<+\infty\). 证明等式
	\[
	\int_{\R^d}^{*}(f+g)(x)\,\dd x
	=
	\int_{\R^d}^{*}f(x)\,\dd x+
	\int_{\R^d}^{*}g(x)\,\dd x
	\]
	对任意 \(g\geq0\) 成立的充分必要条件是 \(f\in L^1\).
\end{exer}

\subsection{单调收敛定理}

\begin{thm}[Beppo Levi 单调收敛定理]
	\label{thm:beppo-levi-l1}
	设 \(f_n\in L^1\), 且
	\[
	f_1\leq f_2\leq\cdots,
	\qquad
	f_n(x)\uparrow f(x)
	\]
	几乎处处成立. 若
	\begin{equation}
		\lim_{n\to\infty}\int_{\R^d}f_n(x)\,\dd x<+\infty,
		\label{eq:beppo-levi-finite-limit}
	\end{equation}
	则 \(f\in L^1\), 且
	\begin{equation}
		\int_{\R^d}f(x)\,\dd x
		=
		\lim_{n\to\infty}\int_{\R^d}f_n(x)\,\dd x.
		\label{eq:beppo-levi-conclusion}
	\end{equation}
	反之, 若 \(f\in L^1\), 则右端极限有限, 并且 \eqref{eq:beppo-levi-conclusion} 成立.
\end{thm}

\begin{proof}
	由于 \(f_n\leq f_{n+1}\), 数列 \(\int f_n\) 单调递增. 下面先从 \eqref{eq:beppo-levi-finite-limit} 出发证明主要部分.
	
	令
	\[
	A=\lim_{n\to\infty}\int_{\R^d}f_n(x)\,\dd x<+\infty.
	\]
	固定 \(m\). 因为 \(f_n-f_m\uparrow f-f_m\), 且 \(f_n-f_m\geq0\), 上积分的单调收敛性质给出
	\[
	\int_{\R^d}^{*}(f-f_m)(x)\,\dd x
	=
	\lim_{n\to\infty}\int_{\R^d}(f_n-f_m)(x)\,\dd x
	=
	A-
	\int_{\R^d}f_m(x)\,\dd x.
	\]
	右端当 \(m\to\infty\) 时趋于 \(0\). 由定理 \ref{thm:integrable-approx-by-integrable}, 以 \(f_m\) 逼近 \(f\), 得 \(f\in L^1\). 又
	\[
	\int_{\R^d}|f-f_m|\,\dd x
	=
	\int_{\R^d}^{*}(f-f_m)\,\dd x\to0,
	\]
	所以根据积分定义,
	\[
	\int_{\R^d}f(x)\,\dd x
	=
	\lim_{m\to\infty}\int_{\R^d}f_m(x)\,\dd x.
	\]
	
	反过来, 若 \(f\in L^1\), 则 \(f_n\leq f\), 从而 \(\int f_n\leq \int f\), 所以 \(\int f_n\) 的极限有限. 令
	\[
	h_n=f-f_n\geq0.
	\]
	则 \(h_n\downarrow0\), 且 \(h_1\in L^1\). 又
	\[
	h_1-h_n=f_n-f_1\uparrow h_1.
	\]
	把刚才已经证明的部分应用于递增列 \(h_1-h_n\), 得
	\[
	\int_{\R^d}(h_1-h_n)\,\dd x\to \int_{\R^d}h_1\,\dd x.
	\]
	于是 \(\int h_n\,\dd x\to0\). 因此
	\[
	\int_{\R^d}f_n\,\dd x
	=
	\int_{\R^d}f\,\dd x-
	\int_{\R^d}h_n\,\dd x
	\longrightarrow
	\int_{\R^d}f\,\dd x.
	\]
\end{proof}

等价地, 单调收敛定理也可以写成正项级数的形式: 若 \(u_n\in L^1\) 且 \(u_n\geq0\), 则
\[
\sum_{n=1}^{\infty}u_n\in L^1
\]
的充分必要条件是
\[
\sum_{n=1}^{\infty}\int_{\R^d}u_n(x)\,\dd x<+\infty.
\]
此时
\begin{equation}
	\int_{\R^d}\left(\sum_{n=1}^{\infty}u_n(x)\right)\dd x
	=
	\sum_{n=1}^{\infty}\int_{\R^d}u_n(x)\,\dd x.
	\label{eq:beppo-levi-series-form}
\end{equation}
这正是第一章正项级数逐项积分定理在勒贝格积分框架下的自然改进.

\begin{thm}\label{thm:lsc-majorant-for-positive-finite-upper}
	设 \(f\geq0\), 且
	\[
	\int_{\R^d}^{*}f(x)\,\dd x<+\infty.
	\]
	则存在函数列 \(g_n\in\mathcal I^+\), 使得
	\[
	f\leq g_n\downarrow g,
	\qquad
	g\in L^1,
	\qquad
	g\geq f,
	\]
	并且
	\begin{equation}
		\int_{\R^d}^{*}f(x)\,\dd x
		=
		\int_{\R^d}g(x)\,\dd x.
		\label{eq:lsc-majorant-integral-equality}
	\end{equation}
	若 \(f\in L^1\), 则 \(f=g\) 几乎处处成立.
\end{thm}

\begin{proof}
	由上积分定义, 可取 \(h_n\in\mathcal I^+\), 使得
	\[
	f\leq h_n,
	\qquad
	\int^*h_n<\int^*f+\frac1n.
	\]
	令
	\[
	g_n=\inf_{1\leq k\leq n}h_k.
	\]
	由于有限个下方半连续函数的下确界仍下方半连续, 所以 \(g_n\in\mathcal I^+\), 且
	\[
	f\leq g_n\downarrow g:=\inf_{n}g_n.
	\]
	又有 \(g_n\leq h_n\), 因而
	\[
	\int^*g_n\leq \int^*h_n<\int^*f+\frac1n.
	\]
	因为 \(g_1-g_n\uparrow g_1-g\), 单调收敛定理给出 \(g_1-g\in L^1\), 并且可知 \(g\in L^1\). 同时
	\[
	\int g=\lim_{n\to\infty}\int^*g_n=\int^*f.
	\]
	若 \(f\in L^1\), 则 \(0\leq g-f\in L^1\), 且
	\[
	\int(g-f)\,\dd x=
	\int g\,\dd x-
	\int f\,\dd x=0.
	\]
	故 \(g=f\) 几乎处处成立.
\end{proof}

\subsection{Fatou 引理与控制收敛定理}

\begin{thm}[Fatou 引理]
	\label{thm:fatou-lemma}
	若 \(f_n\geq0\), 则
	\begin{equation}
		\int_{\R^d}^{*}\bigl(\liminf_{n\to\infty}f_n(x)\bigr)\,\dd x
		\leq
		\liminf_{n\to\infty}\int_{\R^d}^{*}f_n(x)\,\dd x.
		\label{eq:fatou-lemma}
	\end{equation}
	进一步, 若每个 \(f_n\in L^1\), 且右端有限, 则 \(\liminf_{n\to\infty}f_n\in L^1\).
\end{thm}

\begin{proof}
	对 \(m\geq n\), 令
	\[
	h_{n,m}(x)=\min\{f_n(x),f_{n+1}(x),\ldots,f_m(x)\}.
	\]
	记
	\[
	h_n(x)=\inf_{m\geq n}f_m(x).
	\]
	则对固定的 \(n\), 有 \(h_{n,m}\downarrow h_n\), 并且
	\[
	h_n\uparrow \liminf_{n\to\infty}f_n.
	\]
	又因为 \(h_n\leq h_{n,m}\leq f_m\) 对任意 \(m\geq n\) 成立, 由单调性得
	\[
	\int^*h_n\leq \int^*f_m,
	\qquad m\geq n.
	\]
	于是
	\[
	\int^*h_n\leq \inf_{m\geq n}\int^*f_m.
	\]
	令 \(n\to\infty\), 并利用上积分的单调收敛性质, 得
	\[
	\int^*\liminf_{n\to\infty}f_n
	=
	\lim_{n\to\infty}\int^*h_n
	\leq
	\liminf_{n\to\infty}\int^*f_n.
	\]
	
	若 \(f_n\in L^1\) 且右端有限, 则对固定的 \(n\), 有 \(0\leq h_{n,m}\leq f_n\), 且 \(h_{n,m}\downarrow h_n\). 由 Beppo Levi 定理的递减形式可知 \(h_n\in L^1\). 再由 \(h_n\uparrow \liminf f_n\) 以及上面得到的有限积分上界, 再次应用 Beppo Levi 定理, 得 \(\liminf f_n\in L^1\).
\end{proof}

\begin{exer}\label{exer:fatou-example-nexp}
	设一维函数
	\[
	f_n(x)=n^a|x|e^{-n|x|},
	\]
	其中 \(a\in\R\). 分别计算 Fatou 引理 \eqref{eq:fatou-lemma} 两端, 并讨论不同 \(a\) 下的情形.
\end{exer}

\begin{thm}[Lebesgue 控制收敛定理]
	\label{thm:dominated-convergence}
	设 \(f_n\in L^1\), 且
	\[
	f_n(x)\to f(x)
	\]
	几乎处处成立. 若存在非负函数 \(g\), 使得
	\[
	|f_n(x)|\leq g(x)\quad\text{对所有 }n\text{ 几乎处处成立},
	\qquad
	\int_{\R^d}^{*}g(x)\,\dd x<+\infty,
	\]
	则 \(f\in L^1\), 并且
	\begin{equation}
		\int_{\R^d}f(x)\,\dd x
		=
		\lim_{n\to\infty}\int_{\R^d}f_n(x)\,\dd x.
		\label{eq:dominated-convergence-integral-limit}
	\end{equation}
	事实上还有
	\[
	\int_{\R^d}|f_n(x)-f(x)|\,\dd x\to0.
	\]
\end{thm}

\begin{proof}
	由定理 \ref{thm:lsc-majorant-for-positive-finite-upper}, 可用一个 \(L^1\) 函数 \(G\geq g\) 几乎处处替代 \(g\). 因此不妨设 \(g\in L^1\). 由 \(|f_n|\leq g\) 和 \(f_n\to f\), 得 \(|f|\leq g\) 几乎处处, 从而 \(f\in L^1\).
	
	对非负函数列
	\[
	2g-|f_n-f|\geq0
	\]
	应用 Fatou 引理, 得
	\[
	\int 2g\,\dd x
	\leq
	\liminf_{n\to\infty}\int (2g-|f_n-f|)\,\dd x.
	\]
	因此
	\[
	\limsup_{n\to\infty}\int |f_n-f|\,\dd x\leq0.
	\]
	所以
	\[
	\int |f_n-f|\,\dd x\to0.
	\]
	最后,
	\[
	\left|\int f_n\,\dd x-
	\int f\,\dd x\right|
	\leq
	\int |f_n-f|\,\dd x\to0.
	\]
\end{proof}

\begin{exer}
	验证练习 \ref{exer:fatou-example-nexp} 中某些参数取值下, 公式
	\[
	\int f_n\,\dd x\to \int f\,\dd x
	\]
	不成立, 并说明控制函数条件不可缺少.
\end{exer}

\begin{exer}
	在练习 \ref{exer:fatou-example-nexp} 中, 直接验证 \(\sup_n f_n\) 的上积分为无穷大.
\end{exer}

\subsection{级数与完备性}

\begin{thm}\label{thm:absolute-summable-series-l1}
	设 \(u_n\in L^1\), 且
	\begin{equation}
		\sum_{n=1}^{\infty}\int_{\R^d}|u_n(x)|\,\dd x<+\infty.
		\label{eq:absolute-l1-summable-assumption}
	\end{equation}
	则级数
	\begin{equation}
		f(x)=\sum_{n=1}^{\infty}u_n(x)
		\label{eq:l1-series-pointwise}
	\end{equation}
	几乎处处绝对收敛, 且 \(f\in L^1\). 此外, 若
	\[
	S_N(x)=\sum_{n=1}^{N}u_n(x),
	\]
	则
	\begin{equation}
		\int_{\R^d}|f(x)-S_N(x)|\,\dd x\to0.
		\label{eq:l1-series-tail}
	\end{equation}
\end{thm}

\begin{proof}
	由 Beppo Levi 定理,
	\[
	F(x)=\sum_{n=1}^{\infty}|u_n(x)|
	\]
	属于 \(L^1\). 因此 \(F(x)<+\infty\) 几乎处处, 这说明 \(\sum u_n(x)\) 几乎处处绝对收敛. 记其和为 \(f\). 因 \(|S_N|\leq F\), 且 \(S_N\to f\) 几乎处处, 由控制收敛定理知 \(f\in L^1\), 且
	\[
	\int |S_N-f|\,\dd x\to0.
	\]
\end{proof}

\begin{exer}
	利用分解
	\[
	u_n^+=\max(u_n,0),
	\qquad
	u_n^-=\max(-u_n,0)
	\]
	以及 Beppo Levi 定理, 重新证明定理 \ref{thm:absolute-summable-series-l1}.
\end{exer}

\begin{exer}\label{exer:l1-decomposition-i-plus}
	设 \(f\in L^1\). 证明存在 \(f_1,f_2\in\mathcal I^+\), 使得
	\[
	f=f_1-f_2
	\]
	几乎处处成立. 提示: 先把 \(f\) 几乎处处表示成 \(\C_0\) 中函数的绝对收敛级数, 再利用正负部分分解.
\end{exer}

\begin{exer}
	设 \(f\in L^1(\R)\). 证明
	\[
	\sum_{n=1}^{\infty}f(x+n)
	\]
	几乎处处收敛.
\end{exer}

\begin{thm}[Riesz--Fischer 定理]
	\label{thm:riesz-fischer-l1}
	设 \(f_n\in L^1\), 且
	\begin{equation}
		\int_{\R^d}|f_n(x)-f_m(x)|\,\dd x\to0,
		\qquad n,m\to\infty.
		\label{eq:l1-cauchy-condition}
	\end{equation}
	则存在 \(f\in L^1\), 使得
	\begin{equation}
		\int_{\R^d}|f_n(x)-f(x)|\,\dd x\to0.
		\label{eq:l1-converges-to-f}
	\end{equation}
	反过来, \eqref{eq:l1-converges-to-f} 当然推出 \eqref{eq:l1-cauchy-condition}.
\end{thm}

\begin{proof}
	反向结论由三角不等式立即得到. 下面证明正向结论.
	
	由 \eqref{eq:l1-cauchy-condition}, 可取严格递增的整数列 \(n_k\), 使得
	\[
	\int_{\R^d}|f_{n_{k+1}}-f_{n_k}|\,\dd x<2^{-k},
	\qquad k=1,2,\ldots.
	\]
	于是
	\[
	\sum_{k=1}^{\infty}\int_{\R^d}|f_{n_{k+1}}-f_{n_k}|\,\dd x<+\infty.
	\]
	由定理 \ref{thm:absolute-summable-series-l1}, 级数
	\[
	f_{n_1}+\sum_{k=1}^{\infty}(f_{n_{k+1}}-f_{n_k})
	\]
	在 \(L^1\) 中收敛到某个函数 \(f\). 也就是说,
	\[
	\int_{\R^d}|f_{n_k}-f|\,\dd x\to0.
	\]
	最后, 对任意 \(n\), 有
	\[
	\int |f_n-f|\,\dd x
	\leq
	\int |f_n-f_{n_k}|\,\dd x+
	\int |f_{n_k}-f|\,\dd x.
	\]
	先取 \(k\) 充分大, 再令 \(n\) 充分大, 由 Cauchy 条件得到 \eqref{eq:l1-converges-to-f}.
\end{proof}

\begin{rem}
	Riesz--Fischer 定理说明: 若用
	\[
	\|f\|_1=\int_{\R^d}|f(x)|\,\dd x
	\]
	定义距离, 那么 \(L^1\) 是完备空间. 这与实数的 Cauchy 收敛准则完全平行; 证明的关键也是``绝对收敛级数必收敛''.
\end{rem}

\begin{exer}
	在一维情形中, 设 \(n\) 为正整数, 并令 \(m\) 是满足
	\[
	m^2\leq n<(m+1)^2
	\]
	的整数. 定义
	\[
	f_n(x)=
	\begin{cases}
		1, & 0<n-m^2-mx<1,\\
		0, & \text{其他情形}.
	\end{cases}
	\]
	证明 \(\int |f_n|\,\dd x\to0\), 但 \(\lim_{n\to\infty}f_n(x)\) 对某些 \(x\) 不存在.
\end{exer}

\clearpage

\section{可测和局部可积函数}
\label{sec:measurable-and-locally-integrable-functions}

为了刻画可积性, 除了整体的积分有限性之外, 还需要把函数本身的正则性分离出来. 这就引出可测函数的概念. 粗略地说, 可测函数是可以被可积函数局部截断逼近的函数.

若 \(g\in\C_0^+(\R^d)\), 定义截断函数
\[
T_g f(x)=\max\{-g(x),\min(g(x),f(x))\}.
\]
也就是说, \(T_g f\) 把 \(f\) 的值限制在区间 \([-g(x),g(x)]\) 内. 当 \(f(x)=+\infty\) 时, 约定 \(T_g f(x)=g(x)\); 当 \(f(x)=-\infty\) 时, 约定 \(T_g f(x)=-g(x)\).

\begin{defn}[可测函数]
	\label{def:measurable-function}
	设 \(f\) 是 \(\R^d\) 上几乎处处定义并取值于 \([-\infty,+\infty]\) 的函数. 若对任意 \(g\in\C_0^+(\R^d)\), 都有
	\[
	T_g f\in L^1(\R^d),
	\]
	则称 \(f\) 为可测函数.
\end{defn}

显然, 在一个零集合上改变 \(f\) 的取值不会影响可测性. 此外, 连续函数是可测的, 可积函数也是可测的. 后者是因为若 \(f\in L^1\), 则由 \(L^1\) 对 \(\max\) 与 \(\min\) 运算封闭可知 \(T_g f\in L^1\).

\begin{thm}\label{thm:l1-iff-measurable-finite-upper-abs}
	函数 \(f\) 可积的充分必要条件是 \(f\) 可测, 并且
	\[
	\int_{\R^d}^{*}|f(x)|\,\dd x<+\infty.
	\]
\end{thm}

\begin{proof}
	若 \(f\in L^1\), 则上面已经说明 \(f\) 可测. 又因为 \(|f|\in L^1\), 故
	\[
	\int_{\R^d}^{*}|f(x)|\,\dd x
	=
	\int_{\R^d}|f(x)|\,\dd x<+\infty.
	\]
	
	反过来, 设 \(f\) 可测且 \(\int^*|f|<+\infty\). 由前一节的结论, 可取 \(G\in L^1\), 使得
	\[
	|f|\leq G
	\]
	几乎处处成立. 取函数列 \(g_n\in\C_0^+(\R^d)\), 使得
	\[
	0\leq g_n(x)\to +\infty
	\]
	对每个 \(x\in\R^d\) 成立. 例如可以取 \(g_n\) 在 \(\|x\|\leq n\) 上等于 \(n\), 在 \(\|x\|\geq n+1\) 上等于 \(0\), 并在中间连续过渡.
	
	由可测性的定义,
	\[
	T_{g_n}f\in L^1.
	\]
	同时 \(T_{g_n}f\to f\) 几乎处处, 且
	\[
	|T_{g_n}f|\leq |f|\leq G.
	\]
	由 Lebesgue 控制收敛定理,
	\[
	\int_{\R^d}|T_{g_n}f-f|\,\dd x\to0.
	\]
	因此 \(f\) 是 \(L^1\) 函数列 \(T_{g_n}f\) 的 \(L^1\) 极限, 从而 \(f\in L^1\).
\end{proof}

\begin{thm}\label{thm:ae-limit-measurable}
	设 \(f_n\) 是一列可测函数. 若
	\[
	f_n(x)\to f(x)
	\]
	几乎处处成立, 则 \(f\) 可测.
\end{thm}

\begin{proof}
	任取 \(g\in\C_0^+(\R^d)\). 因 \(f_n\) 可测, 有
	\[
	T_g f_n\in L^1.
	\]
	并且
	\[
	T_g f_n(x)\to T_g f(x)
	\]
	几乎处处成立, 同时
	\[
	|T_g f_n(x)|\leq g(x).
	\]
	由于 \(g\in L^1\), 由控制收敛定理得 \(T_g f\in L^1\). 因 \(g\) 任意, \(f\) 可测.
\end{proof}

根据定理 \ref{thm:l1-iff-measurable-finite-upper-abs} 的证明, 一个函数 \(f\) 可测的充分必要条件是: 存在一列有界的 \(L^1\) 函数 \(f_n\), 使得
\[
f_n(x)\to f(x)
\]
几乎处处成立. 事实上, 若 \(f\) 可测, 可以取 \(f_n=T_{g_n}f\), 其中 \(g_n\in\C_0^+\) 逐点趋于 \(+\infty\); 反过来则由定理 \ref{thm:ae-limit-measurable} 得到.

\begin{thm}\label{thm:measurable-closed-lattice-algebra}
	设 \(f,g\) 可测. 则
	\[
	\max(f,g),
	\qquad
	\min(f,g)
	\]
	都是可测函数. 若 \(f,g\) 几乎处处有限, 则
	\[
	f+g,
	\qquad
	fg
	\]
	也是可测函数.
\end{thm}

\begin{proof}
	由上面的等价刻画, 可取有界 \(L^1\) 函数列 \(f_n,g_n\), 使得
	\[
	f_n\to f,
	\qquad
	g_n\to g
	\]
	几乎处处成立. 因为 \(L^1\) 对 \(\max\) 与 \(\min\) 运算封闭, 所以
	\[
	\max(f_n,g_n),
	\qquad
	\min(f_n,g_n)
	\]
	都是 \(L^1\) 函数, 且分别几乎处处收敛到 \(\max(f,g)\) 与 \(\min(f,g)\). 由定理 \ref{thm:ae-limit-measurable}, 二者可测.
	
	若 \(f,g\) 几乎处处有限, 则 \(f_n+g_n\to f+g\) 几乎处处, 且每个 \(f_n+g_n\in L^1\), 因而 \(f+g\) 可测.
	
	最后证明乘积. 由恒等式
	\[
	fg=\frac{(f+g)^2-(f-g)^2}{4},
	\]
	只需证明可测函数的平方仍可测. 若 \(f_n\) 是有界 \(L^1\) 函数且 \(f_n\to f\) 几乎处处, 则 \(f_n^2\in L^1\), 因为 \(f_n\) 有界且可积. 又 \(f_n^2\to f^2\) 几乎处处, 所以 \(f^2\) 可测. 因此 \(fg\) 可测.
\end{proof}

由定理 \ref{thm:ae-limit-measurable} 和定理 \ref{thm:measurable-closed-lattice-algebra} 可知, 若 \(\{f_n\}\) 是可测函数列, 则
\[
\sup_n f_n,
\qquad
\inf_n f_n,
\qquad
\limsup_{n\to\infty}f_n,
\qquad
\liminf_{n\to\infty}f_n
\]
只要表达式有意义, 都是可测函数.

\begin{cor}\label{cor:l1-times-bounded-measurable}
	若 \(f\in L^1\), \(g\) 可测且有界, 则 \(fg\in L^1\).
\end{cor}

\begin{proof}
	由定理 \ref{thm:measurable-closed-lattice-algebra}, \(fg\) 可测. 若 \(|g|\leq M\), 则
	\[
	|fg|\leq M|f|.
	\]
	所以
	\[
	\int^* |fg|\,\dd x\leq M\int |f|\,\dd x<+\infty.
	\]
	由定理 \ref{thm:l1-iff-measurable-finite-upper-abs}, 得 \(fg\in L^1\).
\end{proof}

\subsection{局部可积函数}

可积性不仅包含函数的正则性, 还包含无穷远处的积分有限性. 如果只保留局部积分有限性, 就得到局部可积函数.

\begin{defn}[局部可积函数]
	\label{def:locally-integrable-function}
	设 \(f\) 是 \(\R^d\) 上几乎处处定义并取值于 \([-\infty,+\infty]\) 的函数. 若对任意 \(\varphi\in\C_0(\R^d)\), 都有
	\[
	\varphi f\in L^1(\R^d),
	\]
	则称 \(f\) 为局部可积函数. 所有局部可积函数的集合记作 \(L^1_{\operatorname{loc}}(\R^d)\), 简记为 \(L^1_{\operatorname{loc}}\).
\end{defn}

\begin{exer}\label{exer:locally-integrable-finite-ae}
	证明: 若 \(f\in L^1_{\operatorname{loc}}\), 则 \(f\) 几乎处处有限.
\end{exer}

\begin{exer}\label{exer:locally-integrable-basic}
	证明:
	\begin{enumerate}
		\item 若 \(f\in L^1_{\operatorname{loc}}\), 则 \(f\) 可测;
		\item 若 \(f\) 可测且局部有界, 即 \(f\) 在每个紧集上有界, 则 \(f\in L^1_{\operatorname{loc}}\).
	\end{enumerate}
\end{exer}

下面给出局部可积函数的一个逼近刻画. 注意这里的连续函数不要求具有紧致支集.

\begin{thm}\label{thm:locally-integrable-continuous-approx}
	函数 \(f\in L^1_{\operatorname{loc}}\) 的充分必要条件是: 对任意 \(\varepsilon>0\), 存在一个连续函数 \(g\), 使得
	\[
	\int_{\R^d}^{*}|f(x)-g(x)|\,\dd x<\varepsilon.
	\]
\end{thm}

\begin{proof}
	先设 \(f\in L^1_{\operatorname{loc}}\). 取函数列 \(h_n\in\C_0^+(\R^d)\), 满足
	\[
	0\leq h_n\leq 1,
	\qquad
	h_n(x)=1 \quad (\|x\|\leq n),
	\qquad
	h_n(x)=0 \quad (\|x\|\geq n+1),
	\]
	并令 \(h_0=h_{-1}=0\). 对每个 \(n\geq1\), 函数
	\[
	f(h_n-h_{n-1})
	\]
	属于 \(L^1\). 因而可以取 \(u_n\in\C_0(\R^d)\), 使得
	\[
	\int^*\left|f(h_n-h_{n-1})-u_n\right|\,\dd x
	<\frac{\varepsilon}{2^n}.
	\]
	由于 \(h_{n+1}-h_{n-2}=1\) 于 \(h_n-h_{n-1}\) 的支集上, 且
	\[
	0\leq h_{n+1}-h_{n-2}\leq1,
	\]
	令
	\[
	v_n=u_n(h_{n+1}-h_{n-2}),
	\]
	则 \(v_n\in\C_0(\R^d)\), 且
	\[
	\int^*\left|f(h_n-h_{n-1})-v_n\right|\,\dd x
	\leq
	\int^*\left|f(h_n-h_{n-1})-u_n\right|\,\dd x
	<\frac{\varepsilon}{2^n}.
	\]
	
	定义
	\[
	g(x)=\sum_{n=1}^{\infty}v_n(x).
	\]
	这个和在每个紧集上只有有限多项非零, 因而 \(g\) 是连续函数. 又因为
	\[
	\sum_{n=1}^{\infty}(h_n-h_{n-1})=1,
	\]
	所以由上积分的次可加性,
	\begin{align*}
		\int^*|f-g|\,\dd x
		&\leq
		\sum_{n=1}^{\infty}
		\int^*\left|f(h_n-h_{n-1})-v_n\right|\,\dd x\\
		&<\sum_{n=1}^{\infty}\frac{\varepsilon}{2^n}
		=\varepsilon.
	\end{align*}
	
	反过来, 假设对任意 \(\varepsilon>0\) 都存在连续函数 \(g\), 使 \(\int^*|f-g|<\varepsilon\). 任取 \(\varphi\in\C_0(\R^d)\). 由于 \(g\) 连续, \(\varphi g\in\C_0(\R^d)\subset L^1\). 又
	\[
	|\varphi f-\varphi g|
	\leq \|\varphi\|_\infty |f-g|.
	\]
	因此 \(\varphi f\) 可以在 \(L^1\) 意义下由 \(\varphi g\) 任意逼近. 由定理 \ref{thm:integrable-approx-by-integrable}, 得 \(\varphi f\in L^1\). 因 \(\varphi\) 任意, \(f\in L^1_{\operatorname{loc}}\).
\end{proof}

\begin{exer}\label{exer:locally-integrable-dominated}
	若 \(f\) 可测, \(|f|\leq g\), 且 \(g\in L^1_{\operatorname{loc}}\), 证明 \(f\in L^1_{\operatorname{loc}}\).
\end{exer}

\begin{exer}\label{exer:locally-integrable-lattice}
	若 \(f,g\in L^1_{\operatorname{loc}}\), 证明
	\[
	|f|,
	\quad af\ (a\in\R),
	\quad f+g,
	\quad \max(f,g),
	\quad \min(f,g)
	\]
	都属于 \(L^1_{\operatorname{loc}}\).
\end{exer}

\begin{rem}
	两个局部可积函数的乘积未必局部可积. 例如在一维情形下, 函数 \(|x|^{-1/2}\) 在原点附近局部可积, 但它的平方 \(|x|^{-1}\) 在原点附近不可积.
\end{rem}

\begin{exer}\label{exer:locally-integrable-times-locally-bounded}
	若 \(f\in L^1_{\operatorname{loc}}\), 且 \(g\) 可测并局部有界, 证明
	\[
	fg\in L^1_{\operatorname{loc}}.
	\]
\end{exer}

\clearpage

\section{可测和可积集合}
\label{sec:measurable-and-integrable-sets}

借助前几节建立的积分理论, 我们现在可以定义 \(\R^d\) 中集合的测度. 基本思想很直接: 用集合的特征函数来度量集合的大小.

\begin{defn}[外测度、可测集与可积集]
	\label{def:outer-measure-measurable-integrable-set}
	设 \(E\subset\R^d\), 记 \(\chi_E\) 为 \(E\) 的特征函数. 定义 \(E\) 的外测度为
	\begin{equation}
		m^*(E)=\int_{\R^d}^{*}\chi_E(x)\,\dd x.
		\label{eq:outer-measure-definition}
	\end{equation}
	若 \(\chi_E\) 是可测函数, 则称 \(E\) 为可测集. 若 \(\chi_E\in L^1\), 则称 \(E\) 为可积集, 并定义它的测度为
	\begin{equation}
		m(E)=\int_{\R^d}\chi_E(x)\,\dd x.
		\label{eq:measure-integrable-set-definition}
	\end{equation}
\end{defn}

若 \(E\) 可测且 \(m^*(E)<+\infty\), 则由定理 \ref{thm:l1-iff-measurable-finite-upper-abs} 可知 \(\chi_E\in L^1\), 因而 \(E\) 是可积集. 此时
\[
m(E)=m^*(E).
\]

\subsection{外测度的基本性质}

由上积分的性质立刻得到外测度的基本性质.

\begin{thm}\label{thm:outer-measure-basic-properties}
	设 \(E,E_n\subset\R^d\). 则:
	\begin{enumerate}
		\item 若 \(E_1\subset E_2\), 则
		\[
		m^*(E_1)\leq m^*(E_2).
		\]
		\item 若
		\[
		E\subset \bigcup_{n=1}^{\infty}E_n,
		\]
		则
		\begin{equation}
			m^*(E)\leq \sum_{n=1}^{\infty}m^*(E_n).
			\label{eq:outer-measure-subadditivity}
		\end{equation}
		\item 若 \(E_n\uparrow E\), 即
		\[
		E_1\subset E_2\subset\cdots,
		\qquad
		E=\bigcup_{n=1}^{\infty}E_n,
		\]
		则
		\[
		m^*(E_n)\uparrow m^*(E).
		\]
	\end{enumerate}
\end{thm}

\begin{proof}
	若 \(E_1\subset E_2\), 则 \(\chi_{E_1}\leq\chi_{E_2}\), 由上积分单调性得到第一条.
	
	若 \(E\subset\bigcup_nE_n\), 则
	\[
	\chi_E\leq \sum_{n=1}^{\infty}\chi_{E_n}.
	\]
	由上积分的次可加性得到 \eqref{eq:outer-measure-subadditivity}.
	
	若 \(E_n\uparrow E\), 则 \(\chi_{E_n}\uparrow\chi_E\). 由上积分的单调收敛性质得到
	\[
	\int^*\chi_{E_n}\,\dd x\uparrow\int^*\chi_E\,\dd x,
	\]
	即 \(m^*(E_n)\uparrow m^*(E)\).
\end{proof}

容易看出, \(\chi_E\in\mathcal I^+\) 的充分必要条件是 \(E\) 为开集. 因此, 开集在测度理论中所起的作用, 类似于 \(\mathcal I^+\) 中函数在积分理论中所起的作用.

\begin{thm}[开集外逼近]
	\label{thm:outer-measure-open-approximation}
	对任意集合 \(E\subset\R^d\), 有
	\begin{equation}
		m^*(E)=\inf\{m^*(O): O\supset E,\ O\text{为开集}\}.
		\label{eq:outer-measure-open-approximation}
	\end{equation}
\end{thm}

\begin{proof}
	若 \(O\supset E\), 则由单调性有 \(m^*(E)\leq m^*(O)\). 因而
	\[
	m^*(E)\leq \inf_{O\supset E}m^*(O).
	\]
	下面证明反向不等式. 若 \(m^*(E)=+\infty\), 无需证明. 设 \(m^*(E)<+\infty\). 给定 \(\varepsilon>0\), 由上积分定义, 可取 \(f\in\mathcal I^+\), 使得
	\[
	\chi_E\leq f,
	\qquad
	\int^* f\,\dd x<m^*(E)+\varepsilon.
	\]
	令
	\[
	O=\{x:f(x)>1-\varepsilon\}.
	\]
	因为 \(f\) 下方半连续, 所以 \(O\) 是开集. 又 \(\chi_E\leq f\), 若 \(x\in E\), 则 \(f(x)\geq1\), 因而 \(E\subset O\). 当 \(0<\varepsilon<1\) 时,
	\[
	\chi_O\leq \frac{1}{1-\varepsilon}f.
	\]
	于是
	\[
	m^*(O)\leq \frac{m^*(E)+\varepsilon}{1-\varepsilon}.
	\]
	令 \(\varepsilon\to0\), 得到反向不等式.
\end{proof}

\subsection{可测集与可积集}

从可测函数的运算性质可以直接得到可测集的运算性质.

\begin{thm}\label{thm:measurable-set-basic-properties}
	可测集具有如下性质:
	\begin{enumerate}
		\item 可测集的余集仍为可测集;
		\item 可数多个可测集的并与交仍为可测集;
		\item 若 \(\{E_n\}\) 是可测集列, 则
		\[
		\limsup_{n\to\infty}E_n,
		\qquad
		\liminf_{n\to\infty}E_n
		\]
		都是可测集;
		\item 可测集的对称差是可测集;
		\item 任意开集和闭集都是可测集.
	\end{enumerate}
\end{thm}

\begin{proof}
	设 \(E\) 可测. 因
	\[
	\chi_{E^c}=1-\chi_E,
	\]
	且常数函数局部可测, 由可测函数的线性运算性质可知 \(E^c\) 可测.
	
	若 \(E_n\) 可测, 则 \(\chi_{E_n}\) 可测. 注意
	\[
	\chi_{\bigcup_{n=1}^{\infty}E_n}
	=
	\sup_n \chi_{E_n}.
	\]
	可测函数列的上确界可测, 故 \(\bigcup_nE_n\) 可测. 交集由余集和并集得到.
	
	由集合极限的表示式
	\[
	\limsup_{n\to\infty}E_n=\bigcap_{N=1}^{\infty}\bigcup_{n\geq N}E_n,
	\qquad
	\liminf_{n\to\infty}E_n=\bigcup_{N=1}^{\infty}\bigcap_{n\geq N}E_n,
	\]
	立刻得到第三条. 对称差满足
	\[
	\chi_{E\triangle F}=|\chi_E-\chi_F|,
	\]
	故可测. 最后, 开集 \(O\) 的特征函数 \(\chi_O\) 下方半连续, 因而可测; 闭集是开集的余集, 也可测.
\end{proof}

\begin{thm}\label{thm:integrable-set-basic-properties}
	可积集具有如下性质:
	\begin{enumerate}
		\item 可数多个可积集的交仍为可积集;
		\item 开集或闭集为可积集的充分必要条件是它的外测度有限. 特别地, 任意紧集都是可积集;
		\item 若 \(E_n\) 是可积集列, 且
		\[
		E=\bigcup_{n=1}^{\infty}E_n,
		\]
		并且 \(\sum_{n=1}^{\infty}m(E_n)<+\infty\), 则 \(E\) 是可积集, 且
		\[
		m(E)\leq \sum_{n=1}^{\infty}m(E_n).
		\]
		若进一步 \(E_n\) 两两不交, 则
		\begin{equation}
			m(E)=\sum_{n=1}^{\infty}m(E_n).
			\label{eq:countable-additivity-measure}
		\end{equation}
	\end{enumerate}
\end{thm}

\begin{proof}
	若 \(E_n\) 都可积, 则它们都可测. 因此 \(E=\bigcap_nE_n\) 可测. 又 \(E\subset E_1\), 所以
	\[
	m^*(E)\leq m(E_1)<+\infty.
	\]
	因此 \(E\) 可积.
	
	开集和闭集都可测, 故它们可积当且仅当外测度有限. 紧集包含在某个有限区间中, 因而外测度有限, 所以紧集可积.
	
	若 \(E=\bigcup_nE_n\), 则 \(E\) 可测. 又由外测度次可加性,
	\[
	m^*(E)\leq\sum_{n=1}^{\infty}m(E_n)<+\infty,
	\]
	所以 \(E\) 可积, 且得到不等式.
	
	若 \(E_n\) 两两不交, 则对每个 \(N\),
	\[
	\chi_{\bigcup_{n=1}^NE_n}=\sum_{n=1}^N\chi_{E_n}.
	\]
	令 \(N\to\infty\), 由 Beppo Levi 单调收敛定理得到
	\[
	m(E)=\int\chi_E\,\dd x
	=\sum_{n=1}^{\infty}\int\chi_{E_n}\,\dd x
	=\sum_{n=1}^{\infty}m(E_n).
	\]
\end{proof}

\begin{rem}
	对于不是可积集的可测集, 也常常定义它的测度为
	\[
	m(E)=m^*(E)=+\infty.
	\]
	在这个约定下, 公式 \eqref{eq:countable-additivity-measure} 对所有两两不交的可测集列都成立, 允许等式两端同时为 \(+\infty\).
\end{rem}

\begin{exer}\label{exer:measurable-by-compact-intersections}
	证明: 集合 \(E\subset\R^d\) 可测的充分必要条件是对任意紧集 \(K\), 集合 \(E\cap K\) 可测. 等价地, 也可以要求 \(E\cap K\) 可积.
\end{exer}

\begin{exer}\label{exer:outer-measure-by-interval-cover}
	证明
	\[
	m^*(E)=\inf\sum_{n=1}^{\infty}m(I_n),
	\]
	其中下确界取遍所有覆盖 \(E\) 的区间列 \(\{I_n\}_{n=1}^{\infty}\), 即
	\[
	E\subset\bigcup_{n=1}^{\infty}I_n.
	\]
	提示: 先证明每个开集都可以写成可数多个区间的并, 这些区间的顶点可取为有理点. 然后对有限个区间作加细, 使它们没有公共内点.
\end{exer}

\begin{rem}
	由练习 \ref{exer:outer-measure-by-interval-cover} 可得零集的通常刻画: 集合 \(E\) 是零集的充分必要条件是对任意 \(\varepsilon>0\), 存在区间列 \(\{I_n\}\), 使得
	\[
	E\subset\bigcup_{n=1}^{\infty}I_n,
	\qquad
	\sum_{n=1}^{\infty}m(I_n)<\varepsilon.
	\]
\end{rem}

\begin{exer}\label{exer:integrable-set-finite-interval-approx}
	证明: \(E\) 为可积集的充分必要条件是对任意 \(\varepsilon>0\), 存在有限多个区间 \(I_1,\ldots,I_N\), 使得
	\[
	m^*\left(E\triangle\bigcup_{j=1}^N I_j\right)<\varepsilon.
	\]
\end{exer}

\subsection{定义在可测集上的积分}

至此, 我们已经在整个 \(\R^d\) 上建立了积分. 若 \(E\subset\R^d\) 是可测集, \(f\) 是定义在 \(E\) 上的函数, 令它的零延拓为
\[
\widetilde f(x)=
\begin{cases}
	f(x), & x\in E,\\
	0, & x\notin E.
\end{cases}
\]
若 \(\widetilde f\) 可测或可积, 则称 \(f\) 在 \(E\) 上可测或可积, 并定义
\[
\int_E f(x)\,\dd x
=
\int_{\R^d}\widetilde f(x)\,\dd x.
\]

若 \(E'\subset E\) 也是可测集, 且 \(f\in L^1(E)\), 则 \(f\in L^1(E')\). 这是因为
\[
\widetilde{f|_{E'}}=\widetilde f\,\chi_{E'},
\]
而 \(\chi_{E'}\) 是有界可测函数.

若 \(E\) 与 \(F\) 为不交可测集, 且 \(f\) 在 \(E\cup F\) 上可积, 则
\[
\int_{E\cup F}f(x)\,\dd x
=
\int_E f(x)\,\dd x+
\int_F f(x)\,\dd x.
\]
这些性质都是由整个空间上的相应结论直接得到的.

\begin{exer}\label{exer:lloc-iff-integrable-on-compact}
	证明: \(f\in L^1_{\operatorname{loc}}(\R^d)\) 的充分必要条件是 \(f\) 在每个紧集上可积.
\end{exer}

\subsection{可测集的刻画}

下面的定理给出可测集的外逼近和内逼近刻画. 在从测度论出发的教科书中, 这些条件往往被用作可测性的定义.

\begin{thm}\label{thm:measurable-set-characterizations}
	对集合 \(E\subset\R^d\), 下列条件等价:
	\begin{enumerate}
		\item \(E\) 是可测集;
		\item 对任意 \(\varepsilon>0\), 存在开集 \(O\supset E\), 使得
		\[
		m^*(O\setminus E)<\varepsilon;
		\]
		\item 对任意 \(\varepsilon>0\), 存在闭集 \(F\subset E\), 使得
		\[
		m^*(E\setminus F)<\varepsilon;
		\]
		\item 存在下降的开集列 \(O_n\), 使得 \(E\subset O_n\), 且
		\[
		\left(\bigcap_{n=1}^{\infty}O_n\right)\setminus E
		\]
		是零集合;
		\item 存在上升的闭集列 \(F_n\), 使得 \(F_n\subset E\), 且
		\[
		E\setminus\left(\bigcup_{n=1}^{\infty}F_n\right)
		\]
		是零集合.
	\end{enumerate}
\end{thm}

\begin{proof}
	由于余集把开集变成闭集, 并把外逼近变成内逼近, 只需证明 \((1)\)、\((2)\)、\((4)\) 等价.
	
	先设 \(E\) 可测. 则 \(\chi_E\in L^1_{\operatorname{loc}}\). 由定理 \ref{thm:locally-integrable-continuous-approx}, 对给定 \(\varepsilon>0\), 可取连续函数 \(g\), 使得
	\[
	\int^*|\chi_E-g|\,\dd x<\frac{\varepsilon}{4}.
	\]
	于是存在 \(h\in\mathcal I^+\), 使得
	\[
	|\chi_E-g|\leq h,
	\qquad
	\int^*h\,\dd x<\frac{\varepsilon}{4}.
	\]
	令
	\[
	O=\{x:g(x)+h(x)>1/2\}.
	\]
	因为 \(g+h\) 下方半连续, \(O\) 是开集. 若 \(x\in E\), 则 \(\chi_E(x)=1\), 且 \(|\chi_E(x)-g(x)|\leq h(x)\), 从而
	\[
	g(x)+h(x)\geq1,
	\]
	所以 \(E\subset O\).
	若 \(x\in O\setminus E\), 则 \(\chi_E(x)=0\). 此时 \(|g(x)|\leq h(x)\), 而 \(g(x)+h(x)>1/2\), 故
	\[
	2h(x)>1/2,
	\qquad
	\text{即}\qquad
	\chi_{O\setminus E}(x)\leq4h(x).
	\]
	因此
	\[
	m^*(O\setminus E)
	\leq4\int^*h\,\dd x<\varepsilon.
	\]
	这证明 \((1)\Rightarrow(2)\).
	
	再设 \((2)\) 成立. 对每个 \(n\), 取开集 \(O_n'\supset E\), 使得
	\[
	m^*(O_n'\setminus E)<\frac1n.
	\]
	令
	\[
	O_n=O_1'\cap\cdots\cap O_n'.
	\]
	则 \(O_n\) 是下降的开集列, 且 \(E\subset O_n\). 又
	\[
	\bigcap_{n=1}^{\infty}O_n\setminus E
	\subset O_n'\setminus E
	\]
	对每个 \(n\) 成立, 所以
	\[
	m^*\left(\bigcap_{n=1}^{\infty}O_n\setminus E\right)=0.
	\]
	于是 \((4)\) 成立.
	
	最后, 若 \((4)\) 成立, 设
	\[
	N=\left(\bigcap_{n=1}^{\infty}O_n\right)\setminus E.
	\]
	则 \(N\) 是零集, 从而可测. 又 \(\bigcap_nO_n\) 可测, 因此
	\[
	E=\left(\bigcap_{n=1}^{\infty}O_n\right)\setminus N
	\]
	可测. 这证明 \((4)\Rightarrow(1)\).
\end{proof}

\begin{exer}\label{exer:integrable-set-open-compact-approx}
	证明: \(E\) 为可积集的充分必要条件是对任意 \(\varepsilon>0\), 存在开集 \(O\) 和紧集 \(K\), 使得
	\[
	K\subset E\subset O,
	\qquad
	m(O\setminus K)=m(O)-m(K)<\varepsilon.
	\]
\end{exer}

\begin{exer}\label{exer:integrable-measure-inner-regular}
	假定 \(E\) 可积. 证明
	\[
	m(E)=\sup\{m(K):K\subset E,\ K\text{为紧集}\}.
	\]
\end{exer}

\subsection{一个不可测集合}

下面给出一个不可测集合的例子. 在 \(\R\) 上定义等价关系:
\[
x\sim y
\quad\Longleftrightarrow\quad
x-y\in\mathbb Q.
\]
于是实数被分成若干等价类. 从每个等价类中选取恰好一个属于 \([0,1)\) 的代表元, 由此得到集合 \(E\). 这个选择需要选择公理.

下面证明 \(E\) 不可测. 对 \(r\in\mathbb Q\), 令
\[
E_r=r+E=\{r+x:x\in E\}.
\]
因为每个实数都与某个 \([0,1)\) 中的代表元相差一个有理数, 所以
\begin{equation}
	\bigcup_{r\in\mathbb Q}E_r=\R.
	\label{eq:vitali-cover-real-line}
\end{equation}
另一方面, 若 \(0<r<1\), 则
\begin{equation}
	E_r\subset(0,2).
	\label{eq:vitali-translates-contained}
\end{equation}
并且这些集合 \(E_r\) 两两不交. 因若 \(r+x=s+y\), 其中 \(x,y\in E\), \(r,s\in\mathbb Q\), 则 \(x-y=s-r\in\mathbb Q\), 所以 \(x\sim y\). 由于 \(E\) 每个等价类只取一个点, 得 \(x=y\), 从而 \(r=s\).

若 \(E\) 可测, 则每个 \(E_r\) 可测且测度相同. 由 \eqref{eq:vitali-translates-contained} 和可数可加性,
\[
\sum_{\substack{r\in\mathbb Q\\0<r<1}}m(E_r)
=m\left(\bigcup_{\substack{r\in\mathbb Q\\0<r<1}}E_r\right)
\leq2.
\]
由于左端有可数无穷多个相等的非负项, 只能有 \(m(E_r)=0\). 于是由 \eqref{eq:vitali-cover-real-line}, \(\R\) 是可数个零集的并, 从而应是零集, 矛盾. 所以 \(E\) 不可测.

\subsection{可测函数的集合刻画}

现在证明这里关于可测性的定义与通常定义等价.

\begin{thm}\label{thm:measurable-function-level-set-characterization}
	函数 \(f\) 可测的充分必要条件是: 对任意实数 \(a\), 集合
	\[
	\{x:f(x)>a\}
	\]
	都是可测集.
\end{thm}

\begin{proof}
	先证必要性. 若 \(f\) 可测, 则 \(f-a\) 可测, 故只需考虑 \(a=0\). 令
	\[
	E=\{x:f(x)>0\}.
	\]
	因为 \(f^+=\max(f,0)\) 可测, 且
	\[
	\chi_E=\lim_{n\to\infty}\min(nf^+,1),
	\]
	所以 \(\chi_E\) 可测, 即 \(E\) 可测.
	
	反过来, 假设对任意实数 \(a\), 集合 \(\{x:f(x)>a\}\) 可测. 对每个正整数 \(n\), 定义阶梯函数
	\[
	f_n(x)=\frac1n\lfloor nf(x)\rfloor,
	\]
	其中 \(\lfloor t\rfloor\) 表示不超过 \(t\) 的最大整数. 对每个 \(k\in\mathbb Z\), 集合
	\[
	E_{n,k}=\left\{x:\frac{k}{n}<f(x)\leq \frac{k+1}{n}\right\}
	\]
	可测, 因而 \(f_n\) 可测. 又
	\[
	0\leq f(x)-f_n(x)\leq \frac1n
	\]
	在 \(f\) 有限处成立. 故 \(f_n\to f\) 几乎处处, 从而 \(f\) 可测.
\end{proof}

\begin{exer}\label{exer:composition-measurable}
	假定 \(f\) 可测, 且 \(\phi\) 在 \(\R\) 上连续或单调递增. 若 \(f\) 几乎处处有限, 证明 \(\phi\circ f\) 可测.
\end{exer}

\subsection{Lusin 定理与 Riemann 可积性的判别}

下面的定理说明, 可测函数在去掉一个任意小测度的集合后, 可以看成连续函数.

\begin{thm}[Lusin 定理]
	\label{thm:lusin}
	设 \(f\) 是几乎处处有限的可测函数. 则对任意 \(\varepsilon>0\), 存在闭集 \(F\subset\R^d\), 使得 \(f\) 在 \(F\) 上的限制连续, 且
	\[
	m(F^c)<\varepsilon.
	\]
	反过来, 若函数 \(f\) 具有上述性质, 则 \(f\) 可测.
\end{thm}

\begin{proof}
	先证必要性. 令
	\[
	\psi(t)=\frac{t}{1+|t|},
	\qquad
	\psi(+\infty)=1,
	\quad
	\psi(-\infty)=-1.
	\]
	函数 \(f\) 的可测性与 \(\psi\circ f\) 的可测性等价, 而连续性也可通过 \(\psi\) 还原. 因此不妨假定 \(f\in L^1_{\operatorname{loc}}\) 且 \(|f|\leq1\).
	
	由定理 \ref{thm:locally-integrable-continuous-approx}, 可取连续函数列 \(g_n\), 使得
	\[
	\int^*|f-g_n|\,\dd x<2^{-2n}\varepsilon.
	\]
	取 \(h_n\in\mathcal I^+\), 使得
	\[
	|f-g_n|\leq h_n,
	\qquad
	\int^*h_n\,\dd x<2^{-2n}\varepsilon.
	\]
	令
	\[
	E_n=\{x:h_n(x)>2^{-n}\}.
	\]
	则 \(E_n\) 为开集, 且
	\[
	2^{-n}\chi_{E_n}\leq h_n,
	\qquad
	m^*(E_n)\leq 2^n\int^*h_n\,\dd x<2^{-n}\varepsilon.
	\]
	取 \(N\) 足够大, 使得
	\[
	\sum_{n=N}^{\infty}m^*(E_n)<\varepsilon.
	\]
	令
	\[
	F=\left(\bigcup_{n=N}^{\infty}E_n\right)^c.
	\]
	则 \(F\) 为闭集且 \(m(F^c)<\varepsilon\). 对 \(x\in F\), 有
	\[
	|f(x)-g_n(x)|\leq h_n(x)\leq 2^{-n},
	\qquad n\geq N.
	\]
	因此 \(g_n\) 在 \(F\) 上一致收敛于 \(f\). 由于每个 \(g_n\) 连续, 所以 \(f|_F\) 连续.
	
	再证充分性. 由假设, 可取闭集 \(F_n\), 使得 \(f|_{F_n}\) 连续, 且
	\[
	m(F_n^c)<2^{-n}.
	\]
	定义
	\[
	f_n(x)=
	\begin{cases}
		f(x), & x\in F_n,\\
		+\infty, & x\notin F_n.
	\end{cases}
	\]
	因为 \(f|_{F_n}\) 连续且 \(F_n\) 闭, 可验证 \(f_n\) 是下方半连续函数, 从而可测. 又
	\[
	m\left(\bigcap_{N=1}^{\infty}\bigcup_{n\geq N}F_n^c\right)=0
	\]
	由 Borel--Cantelli 型估计得到, 即除一个零集外, 每个点属于除有限多个以外的所有 \(F_n\). 因此 \(f_n\to f\) 几乎处处. 由可测函数列的极限可测性, 得 \(f\) 可测.
\end{proof}

\begin{exer}\label{exer:lusin-for-characteristic}
	证明: 当 \(f=\chi_E\) 是可测集的特征函数时, Lusin 定理可以由定理 \ref{thm:measurable-set-characterizations} 中的开集外逼近和闭集内逼近直接推出.
\end{exer}

\begin{rem}
	Lusin 定理并不意味着可测函数一定在某些点连续. 例如 Dirichlet 函数
	\[
	f(x)=
	\begin{cases}
		1, & x\in\mathbb Q,\\
		0, & x\notin\mathbb Q,
	\end{cases}
	\]
	是可测函数, 但它在每一点都不连续.
\end{rem}

\begin{thm}[Lebesgue 的 Riemann 可积性判别]
	\label{thm:riemann-integrable-iff-discontinuities-null}
	设 \(f\) 是具有紧致支集的有界函数. 则 \(f\) Riemann 可积的充分必要条件是它的不连续点集合为零集合.
\end{thm}

\begin{proof}
	对每个 \(x\), 定义
	\[
	f_1(x)=\liminf_{y\to x}f(y),
	\qquad
	f_2(x)=\limsup_{y\to x}f(y).
	\]
	则 \(f_1\) 下方半连续, \(f_2\) 上方半连续, 且
	\[
	f_1\leq f\leq f_2.
	\]
	函数 \(f\) 在 \(x\) 连续当且仅当 \(f_1(x)=f_2(x)\). 因而不连续点集合为
	\[
	D=\{x:f_1(x)<f_2(x)\}.
	\]
	
	按照上下 Riemann 积分的定义, 可以证明
	\[
	\underline{\int} f\,\dd x=\int f_1\,\dd x,
	\qquad
	\overline{\int} f\,\dd x=\int f_2\,\dd x,
	\]
	其中右端是 Lebesgue 积分. 因此 \(f\) Riemann 可积的充分必要条件是
	\[
	\int (f_2-f_1)\,\dd x=0.
	\]
	这又等价于 \(f_2=f_1\) 几乎处处, 即不连续点集合 \(D\) 为零集合.
\end{proof}

\subsection{Egorov 定理}

\begin{thm}[Egorov 定理]
	\label{thm:egorov}
	设 \(f_n\) 是几乎处处取有限值的可测函数列, 且
	\[
	f_n(x)\to f(x)
	\]
	几乎处处成立, 其中 \(f\) 也几乎处处有限. 则对任意紧集 \(K\subset\R^d\) 和任意 \(\varepsilon>0\), 存在紧集 \(K_1\subset K\), 使得
	\[
	m(K\setminus K_1)<\varepsilon,
	\]
	并且 \(f_n\) 在 \(K_1\) 上一致收敛于 \(f\).
\end{thm}

\begin{proof}
	对 \(n=1,2,\ldots\), 定义
	\[
	h_n(x)=\sup_{m\geq n}|f_m(x)-f(x)|.
	\]
	若某个函数值为无限, 则约定 \(h_n(x)=+\infty\). 由假设, 除一个零集外, \(h_n(x)\downarrow0\). 对 \(\delta>0\), 令
	\[
	E_{n,\delta}=\{x\in K:h_n(x)>\delta\}.
	\]
	则 \(E_{n,\delta}\) 是可测集, 且随 \(n\) 下降. 又
	\[
	\bigcap_{n=1}^{\infty}E_{n,\delta}
	\]
	是零集. 因为 \(K\) 可积, 由控制收敛定理或测度的连续性可得
	\[
	m(E_{n,\delta})\to0,
	\qquad n\to\infty.
	\]
	
	对每个 \(j\), 取 \(N_j\) 足够大, 使得
	\[
	m(E_{N_j,1/j})<\frac{\varepsilon}{2^j}.
	\]
	令
	\[
	E=\bigcup_{j=1}^{\infty}E_{N_j,1/j}.
	\]
	则
	\[
	m(E)<\varepsilon.
	\]
	若 \(x\in K\setminus E\), 则对每个 \(j\), 当 \(n\geq N_j\) 时,
	\[
	|f_n(x)-f(x)|
	\leq h_{N_j}(x)
	\leq \frac1j.
	\]
	这说明 \(f_n\to f\) 在 \(K\setminus E\) 上一致成立.
	
	最后, 由可积集的内正则性, 可取紧集 \(K_1\subset K\setminus E\), 使得
	\[
	m(K\setminus K_1)<\varepsilon.
	\]
	于是 \(f_n\) 在 \(K_1\) 上一致收敛到 \(f\).
\end{proof}

\begin{exer}\label{exer:dct-lusin-from-egorov}
	证明 Lebesgue 控制收敛定理和 Lusin 定理都可以由 Egorov 定理推出.
\end{exer}

\clearpage

\section{变量替换}
\label{sec:change-of-variables-lebesgue}

首先说明一个记号上的约定. 若 \(\Omega\subset\R^d\) 是开集, 则可以完全仿照前面在 \(\R^d\) 上的作法, 对只定义在 \(\Omega\) 上的函数定义上积分、可测性和积分. 具体地说, 从
\[
\C_0(\Omega)=\{f\in \C_0(\R^d):\operatorname{supp}f\subset\Omega\}
\]
出发, 先定义 \(\mathcal I^+(\Omega)\), 再定义非负函数的上积分, 最后定义 \(L^1(\Omega)\).

等价地, 若 \(f\) 是定义在 \(\Omega\) 上的函数, 令其零延拓为
\[
\widetilde f(x)=
\begin{cases}
	f(x), & x\in\Omega,\\
	0, & x\notin\Omega.
\end{cases}
\]
则 \(f\in L^1(\Omega)\) 当且仅当 \(\widetilde f\in L^1(\R^d)\). 在这种情况下,
\[
\int_\Omega f(x)\,\dd x
=
\int_{\R^d}\widetilde f(x)\,\dd x.
\]
同样的约定也适用于可测函数与上积分.

\begin{exer}\label{exer:l1-on-open-set-zero-extension}
	证明: 若 \(\Omega\subset\R^d\) 为开集, 则 \(f\in L^1(\Omega)\) 的充分必要条件是其零延拓 \(\widetilde f\in L^1(\R^d)\). 并证明相应积分相等.
\end{exer}

现在设 \(\Omega_1,\Omega_2\subset\R^d\) 是开集, 并设
\[
\phi:\Omega_1\longrightarrow\Omega_2
\]
是一个一一映射, 且 \(\phi\) 与 \(\phi^{-1}\) 都连续可微. 若 \(f\) 是定义在 \(\Omega_2\) 上的函数, 仍记
\[
\phi^*f=f\circ\phi,
\qquad
(\phi^*f)(x)=f(\phi(x)).
\]

\begin{thm}[Lebesgue 积分的变量替换公式]
	\label{thm:change-of-variables-lebesgue}
	在上述假设下, \(f\in L^1(\Omega_2)\) 的充分必要条件是
	\[
	(\phi^*f)|\det\phi'|\in L^1(\Omega_1).
	\]
	此时有
	\begin{equation}
		\int_{\Omega_2}f(y)\,\dd y
		=
		\int_{\Omega_1}f(\phi(x))|\det\phi'(x)|\,\dd x.
		\label{eq:change-of-variables-lebesgue}
	\end{equation}
	这里
	\[
	\phi'(x)=\left(\frac{\partial \phi_j}{\partial x_k}(x)\right)_{j,k=1}^{d}.
	\]
\end{thm}

\begin{proof}
	由第一章的变量替换公式可知, 当 \(f\in\C_0(\Omega_2)\) 时,
	\begin{equation}
		\int_{\Omega_2}f(y)\,\dd y
		=
		\int_{\Omega_1}f(\phi(x))|\det\phi'(x)|\,\dd x.
		\label{eq:change-of-variables-c0}
	\end{equation}
	
	先把公式推广到 \(\mathcal I^+(\Omega_2)\) 中的函数. 设 \(f\in\mathcal I^+(\Omega_2)\). 由下方半连续正函数的表示定理, 可取 \(f_n\in\C_0^+(\Omega_2)\), 使得
	\[
	0\leq f_n\uparrow f.
	\]
	由 \eqref{eq:change-of-variables-c0}, 对每个 \(n\) 有
	\[
	\int_{\Omega_2}f_n(y)\,\dd y
	=
	\int_{\Omega_1}f_n(\phi(x))|\det\phi'(x)|\,\dd x.
	\]
	令 \(n\to\infty\), 并利用上积分对单调极限的性质, 得
	\begin{equation}
		\int_{\Omega_2}^{*}f(y)\,\dd y
		=
		\int_{\Omega_1}^{*}(\phi^*f)(x)|\det\phi'(x)|\,\dd x.
		\label{eq:change-of-variables-i-plus}
	\end{equation}
	
	由于 \(\phi^{-1}\) 连续可微, 有
	\[
	(\phi^{-1})'(\phi(x))\phi'(x)=I_d,
	\]
	所以 \(\det\phi'(x)\neq0\). 因而乘以正连续函数 \(|\det\phi'|\) 不破坏下方半连续性. 更具体地说,
	\[
	f\in\mathcal I^+(\Omega_2)
	\quad\Longleftrightarrow\quad
	(\phi^*f)|\det\phi'|\in\mathcal I^+(\Omega_1),
	\]
	而且这个对应是双射. 因此, 对任意非负函数 \(f\), 由上方 \(\mathcal I^+\) 函数控制并取下确界, 可从 \eqref{eq:change-of-variables-i-plus} 推出
	\begin{equation}
		\int_{\Omega_2}^{*}f(y)\,\dd y
		=
		\int_{\Omega_1}^{*}(\phi^*f)(x)|\det\phi'(x)|\,\dd x.
		\label{eq:change-of-variables-upper-integral}
	\end{equation}
	
	现在证明可积性的等价. 若 \(f\in L^1(\Omega_2)\), 则对任意 \(\varepsilon>0\), 存在 \(g\in\C_0(\Omega_2)\), 使得
	\[
	\int_{\Omega_2}^{*}|f(y)-g(y)|\,\dd y<\varepsilon.
	\]
	由 \eqref{eq:change-of-variables-upper-integral}, 得
	\[
	\int_{\Omega_1}^{*}|f(\phi(x))-g(\phi(x))|\,|\det\phi'(x)|\,\dd x<\varepsilon.
	\]
	又因为 \(g\in\C_0(\Omega_2)\), 所以
	\[
	(\phi^*g)|\det\phi'|\in\C_0(\Omega_1).
	\]
	于是 \((\phi^*f)|\det\phi'|\) 可以在 \(L^1\) 意义下由 \(\C_0(\Omega_1)\) 中函数任意逼近, 从而
	\[
	(\phi^*f)|\det\phi'|\in L^1(\Omega_1).
	\]
	反过来, 对 \(\phi^{-1}\) 应用同样论证, 即可由 \((\phi^*f)|\det\phi'|\in L^1(\Omega_1)\) 推出 \(f\in L^1(\Omega_2)\).
	
	最后, 公式 \eqref{eq:change-of-variables-lebesgue} 先对 \(\C_0(\Omega_2)\) 中函数成立, 再由刚才证明的可积逼近和上积分公式 \eqref{eq:change-of-variables-upper-integral} 传递到所有 \(L^1(\Omega_2)\) 函数. 定理得证.
\end{proof}

\begin{rem}
	与 Riemann 积分情形相比, Lebesgue 积分的变量替换公式更稳定. 关键原因是上积分满足单调收敛性质, 因而可以从 \(\C_0\) 函数自然推广到 \(\mathcal I^+\), 再推广到任意非负函数和一般可积函数.
\end{rem}

\clearpage

\section{累次积分, Lebesgue--Fubini 定理}
\label{sec:iterated-integrals-lebesgue-fubini}

本节讨论 \(\R^{d+e}\) 上函数的积分与累次积分之间的关系. 我们把 \(\R^{d+e}\) 中的点写成 \((x,y)\), 其中
\[
x\in\R^d,\qquad y\in\R^e.
\]
若 \(f\in\C_0(\R^{d+e})\), 则由第一章连续函数的累次积分公式可知
\begin{equation}
	\int_{\R^{d+e}}f(x,y)\,\dd x\dd y
	=
	\int_{\R^d}\left(\int_{\R^e}f(x,y)\,\dd y\right)\dd x.
	\label{eq:c0-fubini-start}
\end{equation}
这里左端表示把 \(f\) 看作 \(\R^{d+e}\) 上函数的积分, 右端表示先固定 \(x\), 对 \(y\) 积分, 再把所得函数对 \(x\) 积分.

我们现在要把公式 \eqref{eq:c0-fubini-start} 推广到 Lebesgue 积分的情形. 仍然先从 \(\mathcal I^+\) 中的函数开始.

\begin{thm}\label{thm:i-plus-fubini}
	设 \(f\in\mathcal I^+(\R^{d+e})\). 则对每个固定的 \(x\in\R^d\), 函数
	\[
	y\longmapsto f(x,y)
	\]
	属于 \(\mathcal I^+(\R^e)\). 并且函数
	\[
	x\longmapsto \int_{\R^e}^{*}f(x,y)\,\dd y
	\]
	属于 \(\mathcal I^+(\R^d)\), 且
	\begin{equation}
		\int_{\R^{d+e}}^{*}f(x,y)\,\dd x\dd y
		=
		\int_{\R^d}^{*}\left(\int_{\R^e}^{*}f(x,y)\,\dd y\right)\dd x.
		\label{eq:i-plus-fubini}
	\end{equation}
\end{thm}

\begin{proof}
	由 \(\mathcal I^+\) 的表示定理, 可取 \(f_n\in\C_0^+(\R^{d+e})\), 使得
	\[
	0\leq f_n\uparrow f.
	\]
	于是对固定的 \(x\), 有
	\[
	0\leq f_n(x,\cdot)\uparrow f(x,\cdot).
	\]
	由于每个 \(f_n(x,\cdot)\in\C_0^+(\R^e)\), 得
	\[
	f(x,\cdot)\in\mathcal I^+(\R^e).
	\]
	
	令
	\[
	g_n(x)=\int_{\R^e}f_n(x,y)\,\dd y.
	\]
	由 \(f_n\in\C_0(\R^{d+e})\) 和连续函数的累次积分性质可知 \(g_n\in\C_0^+(\R^d)\). 又由上积分的单调收敛性质,
	\[
	g_n(x)\uparrow \int_{\R^e}^{*}f(x,y)\,\dd y.
	\]
	因此
	\[
	x\longmapsto \int_{\R^e}^{*}f(x,y)\,\dd y
	\]
	属于 \(\mathcal I^+(\R^d)\). 最后, 由 \eqref{eq:c0-fubini-start} 和单调收敛性质,
	\begin{align*}
		\int_{\R^{d+e}}^{*}f(x,y)\,\dd x\dd y
		&=\lim_{n\to\infty}\int_{\R^{d+e}}f_n(x,y)\,\dd x\dd y\\
		&=\lim_{n\to\infty}\int_{\R^d}\left(\int_{\R^e}f_n(x,y)\,\dd y\right)\dd x\\
		&=\int_{\R^d}^{*}\left(\int_{\R^e}^{*}f(x,y)\,\dd y\right)\dd x.
	\end{align*}
\end{proof}

对于任意非负函数, 一般只能先得到一个方向的不等式.

\begin{thm}\label{thm:upper-integral-iterated-inequality}
	设 \(f\geq0\) 是 \(\R^{d+e}\) 上的任意函数. 则
	\begin{equation}
		\int_{\R^d}^{*}\left(\int_{\R^e}^{*}f(x,y)\,\dd y\right)\dd x
		\leq
		\int_{\R^{d+e}}^{*}f(x,y)\,\dd x\dd y.
		\label{eq:upper-integral-iterated-inequality}
	\end{equation}
\end{thm}

\begin{proof}
	若 \(f\leq g\in\mathcal I^+(\R^{d+e})\), 则对每个 \(x\), 有
	\[
	\int_{\R^e}^{*}f(x,y)\,\dd y
	\leq
	\int_{\R^e}^{*}g(x,y)\,\dd y.
	\]
	于是由定理 \ref{thm:i-plus-fubini},
	\[
	\int_{\R^d}^{*}\left(\int_{\R^e}^{*}f(x,y)\,\dd y\right)\dd x
	\leq
	\int_{\R^{d+e}}^{*}g(x,y)\,\dd x\dd y.
	\]
	再对所有满足 \(f\leq g\in\mathcal I^+(\R^{d+e})\) 的 \(g\) 取下确界, 即得 \eqref{eq:upper-integral-iterated-inequality}.
\end{proof}

\begin{exer}\label{exer:product-equality-upper-integral}
	证明: 若
	\[
	f(x,y)=g(x)h(y),
	\]
	其中 \(g,h\geq0\), 则 \eqref{eq:upper-integral-iterated-inequality} 中的等号成立.
\end{exer}

\begin{exer}\label{exer:upper-integral-inequality-can-be-strict}
	设 \(g_1,g_2\) 是 \(\R^d\) 上的非负函数. 证明
	\[
	\int_{\R^d}^{*}(g_1+g_2)(x)\,\dd x
	\leq
	\int_{\R^d}^{*}g_1(x)\,\dd x+
	\int_{\R^d}^{*}g_2(x)\,\dd x.
	\]
	再取 \(\R^e\) 中两个互不相交、测度皆为 \(1\) 的开集的特征函数 \(h_1,h_2\), 并令
	\[
	f(x,y)=g_1(x)h_1(y)+g_2(x)h_2(y).
	\]
	证明 \eqref{eq:upper-integral-iterated-inequality} 的两端分别等于上式中的对应两端. 因此, \eqref{eq:upper-integral-iterated-inequality} 中的等号不一定成立.
\end{exer}

\begin{cor}\label{cor:null-function-sections}
	若 \(f\geq0\) 是 \(\R^{d+e}\) 上的零函数, 则对几乎所有的 \(x\in\R^d\), 函数
	\[
	y\longmapsto f(x,y)
	\]
	是 \(\R^e\) 上的零函数.
\end{cor}

\begin{proof}
	由定理 \ref{thm:upper-integral-iterated-inequality},
	\[
	\int_{\R^d}^{*}\left(\int_{\R^e}^{*}f(x,y)\,\dd y\right)\dd x
	\leq
	\int_{\R^{d+e}}^{*}f(x,y)\,\dd x\dd y=0.
	\]
	于是非负函数
	\[
	F(x)=\int_{\R^e}^{*}f(x,y)\,\dd y
	\]
	是零函数. 因而 \(F(x)=0\) 对几乎所有 \(x\) 成立, 即对几乎所有 \(x\), 截面函数 \(f(x,\cdot)\) 是零函数.
\end{proof}

\subsection{Lebesgue--Fubini 定理}

现在处理一般可积函数. 下面的结果是 Fubini 定理的基本形式.

\begin{thm}[Lebesgue--Fubini 定理]
	\label{thm:lebesgue-fubini}
	若 \(f\in L^1(\R^{d+e})\), 则对几乎所有的 \(x\in\R^d\), 函数
	\[
	y\longmapsto f(x,y)
	\]
	属于 \(L^1(\R^e)\). 同时, 函数
	\[
	F(x)=\int_{\R^e}f(x,y)\,\dd y
	\]
	在几乎处处意义下有定义, 且 \(F\in L^1(\R^d)\). 并且
	\begin{equation}
		\int_{\R^{d+e}}f(x,y)\,\dd x\dd y
		=
		\int_{\R^d}\left(\int_{\R^e}f(x,y)\,\dd y\right)\dd x.
		\label{eq:lebesgue-fubini}
	\end{equation}
\end{thm}

\begin{proof}
	因为 \(f\in L^1(\R^{d+e})\), 可以取 \(g_n\in\C_0(\R^{d+e})\), 使得
	\begin{equation}
		\sum_{n=1}^{\infty}
		\int_{\R^{d+e}}|f(x,y)-g_n(x,y)|\,\dd x\dd y<+\infty.
		\label{eq:fubini-approx-summable}
	\end{equation}
	由定理 \ref{thm:upper-integral-iterated-inequality},
	\[
	\int_{\R^d}^{*}\left(
	\int_{\R^e}^{*}|f(x,y)-g_n(x,y)|\,\dd y
	\right)\dd x
	\leq
	\int_{\R^{d+e}}|f-g_n|\,\dd x\dd y.
	\]
	将上式对 \(n\) 求和, 并利用 \eqref{eq:fubini-approx-summable}, 得
	\[
	\int_{\R^d}^{*}\left(
	\sum_{n=1}^{\infty}
	\int_{\R^e}^{*}|f(x,y)-g_n(x,y)|\,\dd y
	\right)\dd x<+\infty.
	\]
	因此对几乎所有 \(x\), 级数
	\[
	\sum_{n=1}^{\infty}
	\int_{\R^e}^{*}|f(x,y)-g_n(x,y)|\,\dd y
	\]
	收敛. 特别地, 对几乎所有 \(x\), 有
	\[
	\int_{\R^e}^{*}|f(x,y)-g_n(x,y)|\,\dd y\to0.
	\]
	由于 \(g_n(x,\cdot)\in\C_0(\R^e)\), 这说明对几乎所有 \(x\), 函数 \(f(x,\cdot)\) 属于 \(L^1(\R^e)\).
	
	对这些 \(x\) 定义
	\[
	F(x)=\int_{\R^e}f(x,y)\,\dd y.
	\]
	同时令
	\[
	G_n(x)=\int_{\R^e}g_n(x,y)\,\dd y.
	\]
	则 \(G_n\in\C_0(\R^d)\). 并且
	\begin{align*}
		\int_{\R^d}^{*}|F(x)-G_n(x)|\,\dd x
		&\leq
		\int_{\R^d}^{*}\left(
		\int_{\R^e}^{*}|f(x,y)-g_n(x,y)|\,\dd y
		\right)\dd x\\
		&\leq
		\int_{\R^{d+e}}|f(x,y)-g_n(x,y)|\,\dd x\dd y\to0.
	\end{align*}
	因此 \(F\in L^1(\R^d)\). 由积分的定义和 \eqref{eq:c0-fubini-start},
	\begin{align*}
		\int_{\R^d}F(x)\,\dd x
		&=\lim_{n\to\infty}\int_{\R^d}G_n(x)\,\dd x\\
		&=\lim_{n\to\infty}\int_{\R^{d+e}}g_n(x,y)\,\dd x\dd y\\
		&=\int_{\R^{d+e}}f(x,y)\,\dd x\dd y.
	\end{align*}
\end{proof}

\begin{rem}
	定理 \ref{thm:lebesgue-fubini} 中的截面函数 \(y\mapsto f(x,y)\) 一般不必对所有 \(x\) 可积. 事实上, 在固定的某个 \(x\) 处任意改变 \(f(x,\cdot)\), 并不会影响 \(f\) 作为 \(\R^{d+e}\) 上函数的可积性.
\end{rem}

\begin{exer}\label{exer:product-measurable-integrable}
	设
	\[
	f(x,y)=g(x)h(y).
	\]
	若 \(g\) 和 \(h\) 可积, 证明 \(f\) 可积. 若 \(g\) 和 \(h\) 可测, 证明 \(f\) 可测.
\end{exer}

\begin{exer}\label{exer:convolution-l1}
	设 \(f,g\in L^1(\R^d)\). 证明对几乎所有的 \(x\in\R^d\), 积分
	\[
	h(x)=\int_{\R^d}f(x-y)g(y)\,\dd y
	\]
	存在. 并证明 \(h\in L^1(\R^d)\), 且
	\[
	\int_{\R^d}|h(x)|\,\dd x
	\leq
	\left(\int_{\R^d}|f(x)|\,\dd x\right)
	\left(\int_{\R^d}|g(x)|\,\dd x\right).
	\]
	进一步, 证明
	\[
	\int_{\R^d}h(x)\,\dd x
	=
	\left(\int_{\R^d}f(x)\,\dd x\right)
	\left(\int_{\R^d}g(x)\,\dd x\right).
	\]
\end{exer}

\begin{exer}\label{exer:sections-of-measurable-functions}
	若 \(f(x,y)\) 是 \(\R^{d+e}\) 上的可测函数, 证明
	\[
	x\longmapsto \int_{\R^e}^{*}|f(x,y)|\,\dd y
	\]
	是可测函数, 并且对几乎所有的 \(x\), 截面函数
	\[
	y\longmapsto f(x,y)
	\]
	是 \(\R^e\) 上的可测函数.
\end{exer}

\begin{rem}
	即使 \eqref{eq:lebesgue-fubini} 的右端作为一个累次积分存在, 也不能单凭这一点断定左端积分存在. 前面的练习 \ref{exer:upper-integral-inequality-can-be-strict} 可以给出相应的例子: 特别地取 \(g_1\) 和 \(g_2\), 使得 \(g_1+g_2\in L^1\), 就能看到仅有某一方向的累次积分存在是不够的.
\end{rem}

\begin{thm}\label{thm:fubini-integrability-criterion}
	设 \(f\) 是 \(\R^{d+e}\) 上的可测函数. 则 \(f\in L^1(\R^{d+e})\) 的充分必要条件是
	\begin{equation}
		\int_{\R^d}^{*}\left(
		\int_{\R^e}^{*}|f(x,y)|\,\dd y
		\right)\dd x<+\infty.
		\label{eq:fubini-integrability-criterion}
	\end{equation}
\end{thm}

\begin{proof}
	若 \(f\in L^1(\R^{d+e})\), 则 \(|f|\in L^1(\R^{d+e})\). 由 Fubini 定理应用于 \(|f|\), 得 \eqref{eq:fubini-integrability-criterion}.
	
	反过来, 假设 \eqref{eq:fubini-integrability-criterion} 成立. 由定理 \ref{thm:l1-iff-measurable-finite-upper-abs}, 只需证明
	\[
	\int_{\R^{d+e}}^{*}|f(x,y)|\,\dd x\dd y<+\infty.
	\]
	因此不妨假定 \(f\geq0\). 取 \(g_n\in\C_0^+(\R^{d+e})\), 使得
	\[
	g_n(x,y)\uparrow +\infty
	\]
	对每个 \((x,y)\) 成立. 因 \(f\) 可测, 截断函数
	\[
	T_{g_n}f=\min(g_n,f)
	\]
	属于 \(L^1(\R^{d+e})\), 且 \(T_{g_n}f\uparrow f\). 由 Fubini 定理,
	\begin{align*}
		\int_{\R^{d+e}}T_{g_n}f\,\dd x\dd y
		&=\int_{\R^d}\left(\int_{\R^e}T_{g_n}f(x,y)\,\dd y\right)\dd x\\
		&\leq
		\int_{\R^d}^{*}\left(
		\int_{\R^e}^{*}f(x,y)\,\dd y
		\right)\dd x<+\infty.
	\end{align*}
	由 Beppo Levi 单调收敛定理, \(f\in L^1(\R^{d+e})\). 定理得证.
\end{proof}

由定理 \ref{thm:lebesgue-fubini} 和定理 \ref{thm:fubini-integrability-criterion} 可知: 若可测函数 \(f\) 满足 \eqref{eq:fubini-integrability-criterion}, 则
\begin{equation}
	\int_{\R^d}\left(\int_{\R^e}f(x,y)\,\dd y\right)\dd x
	=
	\int_{\R^e}\left(\int_{\R^d}f(x,y)\,\dd x\right)\dd y.
	\label{eq:fubini-change-order}
\end{equation}
特别地, 当 \(f\geq0\) 且 \(f\) 可测时, 只要两个累次积分都有意义, 就可以交换积分次序. 这里 \(f\geq0\) 的假定是关键的.

\begin{exer}\label{exer:two-iterated-integrals-unequal}
	证明下面两个累次积分不相等:
	\[
	\int_0^1\left(\int_0^{+\infty}\bigl(e^{-xy}-2e^{-2xy}\bigr)\,\dd y\right)\dd x,
	\]
	和
	\[
	\int_0^{+\infty}\left(\int_0^1\bigl(e^{-xy}-2e^{-2xy}\bigr)\,\dd x\right)\dd y.
	\]
\end{exer}

\begin{cor}\label{cor:measurable-set-null-sections}
	设 \(E\subset\R^{d+e}\) 是可测集. 对 \(x\in\R^d\), 记
	\[
	E_x=\{y\in\R^e:(x,y)\in E\}.
	\]
	则 \(E\) 是零集的充分必要条件是: 对几乎所有的 \(x\), \(E_x\) 是 \(\R^e\) 中的零集.
\end{cor}

\begin{proof}
	若 \(E\) 是零集, 则 \(\chi_E\) 是零函数. 由推论 \ref{cor:null-function-sections}, 对几乎所有的 \(x\), 函数 \(\chi_E(x,\cdot)=\chi_{E_x}\) 是零函数, 即 \(E_x\) 是零集.
	
	反过来, 若对几乎所有的 \(x\), \(E_x\) 是零集, 则
	\[
	\int_{\R^e}^{*}\chi_E(x,y)\,\dd y=0
	\]
	对几乎所有的 \(x\) 成立. 因为 \(E\) 可测, \(\chi_E\) 是可测函数. 由定理 \ref{thm:fubini-integrability-criterion}, \(\chi_E\in L^1\), 且
	\[
	m(E)=\int_{\R^{d+e}}\chi_E(x,y)\,\dd x\dd y=0.
	\]
	故 \(E\) 是零集.
\end{proof}

\begin{rem}
	推论 \ref{cor:measurable-set-null-sections} 中, \(E\) 可测这一条件不能省略. 即使每个截面 \(E_x\) 都只是一个点, 也不能保证 \(E\) 是零集, 甚至不能保证 \(E\) 可测.
\end{rem}

\clearpage

\section{\texorpdfstring{\(L^p\)}{Lp} 空间}
\label{sec:lp-spaces}

设 \(1\leq p<+\infty\). 我们现在定义 \(L^p\) 空间. 这里采用一个基于定理 \ref{thm:l1-iff-measurable-finite-upper-abs} 的定义; 当 \(p=1\) 时, 它与前面已经定义的 \(L^1\) 完全一致.

\begin{defn}[\(L^p\) 空间]
	\label{def:lp-space}
	设 \(1\leq p<+\infty\). 若 \(f\) 是可测函数, 且
	\begin{equation}
		\|f\|_p
		=
		\left(\int_{\R^d}^{*}|f(x)|^p\,\dd x\right)^{1/p}
		<+\infty,
		\label{eq:lp-norm-definition}
	\end{equation}
	则称 \(f\in L^p(\R^d)\), 简记为 \(f\in L^p\).
\end{defn}

由于 \(|f|^p\) 是可测非负函数, 条件 \eqref{eq:lp-norm-definition} 等价于 \(|f|^p\in L^1\). 因而在 \(f\in L^p\) 时, 上式中的上积分可以写成通常积分:
\[
\|f\|_p
=
\left(\int_{\R^d}|f(x)|^p\,\dd x\right)^{1/p}.
\]
和 \(L^1\) 一样, 我们总是把几乎处处相等的函数视为同一个 \(L^p\) 元素.

\begin{thm}[H\"older 不等式]
	\label{thm:holder-inequality}
	设 \(1<p<+\infty\), 且 \(q\) 满足
	\[
	\frac1p+\frac1q=1.
	\]
	若 \(f\in L^p\), \(g\in L^q\), 则 \(fg\in L^1\), 并且
	\begin{equation}
		\left|\int_{\R^d}f(x)g(x)\,\dd x\right|
		\leq
		\|f\|_p\|g\|_q.
		\label{eq:holder-inequality}
	\end{equation}
	反过来, 若 \(f\in L^p\), 且存在常数 \(C\), 使得对任意 \(g\in L^q\) 都有
	\begin{equation}
		\left|\int_{\R^d}f(x)g(x)\,\dd x\right|
		\leq C\|g\|_q,
		\label{eq:holder-converse-assumption}
	\end{equation}
	则
	\[
	\|f\|_p\leq C.
	\]
\end{thm}

\begin{proof}
	令
	\[
	\alpha=\frac1p,
	\qquad
	\beta=\frac1q.
	\]
	则 \(\alpha+\beta=1\). 对任意非负数 \(a,b\), 有 Young 不等式
	\[
	a^\alpha b^\beta\leq \alpha a+\beta b.
	\]
	这是指数函数凸性的直接结果. 令
	\[
	a=|f(x)|^p,
	\qquad
	b=|g(x)|^q,
	\]
	便得
	\[
	|f(x)g(x)|
	\leq
	\frac1p|f(x)|^p+\frac1q|g(x)|^q.
	\]
	这说明 \(fg\in L^1\) 当 \(f\in L^p\), \(g\in L^q\) 时成立.
	
	若 \(\|f\|_p=\|g\|_q=1\), 积分上式得到
	\[
	\int_{\R^d}|f(x)g(x)|\,\dd x
	\leq
	\frac1p+\frac1q=1.
	\]
	一般情形对 \(f/\|f\|_p\) 和 \(g/\|g\|_q\) 应用上面的结论即可, 零函数情形显然. 于是得到 \eqref{eq:holder-inequality}.
	
	下面证明反向论断. 若 \(f=0\) 几乎处处, 结论显然. 否则定义
	\[
	g(x)=
	\begin{cases}
		\dfrac{|f(x)|^p}{f(x)}, & f(x)\neq0,\\[0.8em]
		0, & f(x)=0.
	\end{cases}
	\]
	也就是说, 在实值函数情形下
	\[
	g(x)=\operatorname{sgn}(f(x))|f(x)|^{p-1}.
	\]
	则
	\[
	|g(x)|^q=|f(x)|^p,
	\]
	所以 \(g\in L^q\), 且
	\[
	\|g\|_q=\|f\|_p^{p-1}.
	\]
	把这个 \(g\) 代入 \eqref{eq:holder-converse-assumption}, 得
	\[
	\|f\|_p^p
	=
	\int_{\R^d}|f(x)|^p\,\dd x
	=
	\int_{\R^d}f(x)g(x)\,\dd x
	\leq
	C\|f\|_p^{p-1}.
	\]
	故 \(\|f\|_p\leq C\).
\end{proof}

\begin{exer}
	\label{exer:holder-upper-integral}
	证明: 若把 H\"older 不等式中的积分都换成上积分, 则对任意非负函数 \(f,g\) 仍有
	\[
	\int_{\R^d}^{*}f(x)g(x)\,\dd x
	\leq
	\left(\int_{\R^d}^{*}f(x)^p\,\dd x\right)^{1/p}
	\left(\int_{\R^d}^{*}g(x)^q\,\dd x\right)^{1/q}.
	\]
\end{exer}

\begin{thm}[Minkowski 不等式]
	\label{thm:minkowski-inequality}
	若 \(f,g\in L^p\), \(1\leq p<+\infty\), 则 \(f+g\in L^p\), 并且
	\begin{equation}
		\|f+g\|_p\leq \|f\|_p+\|g\|_p.
		\label{eq:minkowski-inequality}
	\end{equation}
\end{thm}

\begin{proof}
	当 \(p=1\) 时, 这就是 \(L^1\) 中的三角不等式. 下面设 \(p>1\), 并令 \(q\) 为共轭指数.
	
	由可测函数的运算性质, \(f+g\) 可测. 又有
	\[
	|f+g|\leq 2\max(|f|,|g|),
	\]
	从而
	\[
	|f+g|^p\leq 2^p\bigl(|f|^p+|g|^p\bigr).
	\]
	故 \(f+g\in L^p\).
	
	任取 \(h\in L^q\). 由 H\"older 不等式,
	\begin{align*}
		\left|\int_{\R^d}(f+g)h\,\dd x\right|
		&\leq
		\left|\int_{\R^d}fh\,\dd x\right|
		+
		\left|\int_{\R^d}gh\,\dd x\right|\\
		&\leq
		\bigl(\|f\|_p+\|g\|_p\bigr)\|h\|_q.
	\end{align*}
	由定理 \ref{thm:holder-inequality} 的反向部分, 得
	\[
	\|f+g\|_p\leq \|f\|_p+\|g\|_p.
	\]
\end{proof}

\begin{exer}
	\label{exer:minkowski-upper-integral}
	把不等式 \eqref{eq:minkowski-inequality} 推广到 \(f,g\) 不一定可测的情形. 也就是说, 用上积分定义
	\[
	\|f\|_p^{*}
	=
	\left(\int_{\R^d}^{*}|f(x)|^p\,\dd x\right)^{1/p},
	\]
	证明
	\[
	\|f+g\|_p^{*}
	\leq
	\|f\|_p^{*}+\|g\|_p^{*}.
	\]
\end{exer}

下面的定理说明, \(L^p\) 函数仍可以用紧支连续函数按 \(L^p\) 范数逼近. 这与 \(L^1\) 的定义完全平行.

\begin{thm}[\(\C_0\) 在 \(L^p\) 中的稠密性]
	\label{thm:c0-dense-in-lp}
	设 \(1\leq p<+\infty\). 函数 \(f\in L^p\) 的充分必要条件是: 对任意 \(\varepsilon>0\), 存在 \(g\in\C_0(\R^d)\), 使得
	\begin{equation}
		\int_{\R^d}^{*}|f(x)-g(x)|^p\,\dd x<\varepsilon.
		\label{eq:c0-density-lp}
	\end{equation}
\end{thm}

\begin{proof}
	当 \(p=1\) 时, 这就是 \(L^1\) 的定义. 下面设 \(p>1\).
	
	先证必要性. 设 \(f\in L^p\). 定义截断函数
	\[
	f_n(x)=
	\begin{cases}
		f(x), & \dfrac1n<|f(x)|<n,\\[0.8em]
		0, & \text{其他情形}.
	\end{cases}
	\]
	则 \(|f-f_n|^p\leq |f|^p\), 且 \(|f-f_n|^p\to0\) 几乎处处. 由控制收敛定理,
	\[
	\int_{\R^d}|f-f_n|^p\,\dd x\to0.
	\]
	固定一个足够大的 \(n\). 因 \(f_n\) 有界且 \(|f_n|\geq 1/n\) 于其非零处成立, 有
	\[
	|f_n|\leq n^{p-1}|f_n|^p.
	\]
	所以 \(f_n\in L^1\). 由 \(L^1\) 中 \(\C_0\) 的逼近性质, 可取 \(g\in\C_0(\R^d)\), 使得
	\[
	\int_{\R^d}|f_n(x)-g(x)|\,\dd x
	<
	\frac{\varepsilon}{(2n)^{p-1}}.
	\]
	把 \(g\) 截断为
	\[
	\max(-n,\min(n,g))
	\]
	不会增大左端, 因而可假定 \(|g|\leq n\). 此时
	\[
	|f_n-g|^p\leq (2n)^{p-1}|f_n-g|,
	\]
	故
	\[
	\int_{\R^d}|f_n-g|^p\,\dd x<\varepsilon.
	\]
	再结合 \(f_n\to f\) 于 \(L^p\) 中, 即得必要性.
	
	反过来, 假设 \eqref{eq:c0-density-lp} 对任意 \(\varepsilon>0\) 成立. 首先由
	\[
	|f|^p\leq 2^{p-1}\bigl(|f-g|^p+|g|^p\bigr)
	\]
	可知 \(\int^*|f|^p<+\infty\). 还需证明 \(f\) 可测.
	
	取一个处处正的连续函数 \(h\in L^q\), 其中 \(1/p+1/q=1\). 例如可以取
	\[
	h(x)=e^{-\|x\|^2}.
	\]
	由练习 \ref{exer:holder-upper-integral}, 对满足 \eqref{eq:c0-density-lp} 的 \(g\in\C_0\), 有
	\[
	\int_{\R^d}^{*}|h(x)f(x)-h(x)g(x)|\,\dd x
	\leq
	\|h\|_q
	\left(\int_{\R^d}^{*}|f(x)-g(x)|^p\,\dd x\right)^{1/p}.
	\]
	因此 \(hf\) 可以在 \(L^1\) 意义下由 \(hg\in\C_0(\R^d)\) 逼近, 从而 \(hf\in L^1\). 于是 \(hf\) 可测. 因为 \(h>0\) 且 \(1/h\) 连续, 得
	\[
	f=\frac{hf}{h}
	\]
	可测. 故 \(f\in L^p\).
\end{proof}

现在推广 Riesz--Fischer 定理到 \(L^p\) 空间.

\begin{thm}[Riesz--Fischer 定理]
	\label{thm:riesz-fischer-lp}
	设 \(1\leq p<+\infty\), \(f_n\in L^p\), 且
	\begin{equation}
		\|f_n-f_m\|_p\to0,
		\qquad n,m\to\infty.
		\label{eq:lp-cauchy-condition}
	\end{equation}
	则存在 \(f\in L^p\), 使得
	\begin{equation}
		\|f_n-f\|_p\to0.
		\label{eq:lp-convergence}
	\end{equation}
	反过来, \eqref{eq:lp-convergence} 显然推出 \eqref{eq:lp-cauchy-condition}.
\end{thm}

\begin{proof}
	当 \(p=1\) 时, 这已经在前一节证明过. 下面设 \(p>1\), 并令 \(q\) 为共轭指数.
	
	由 Cauchy 条件, 可取严格递增的整数列 \(n_k\), 使得
	\[
	\sum_{k=1}^{\infty}\|f_{n_{k+1}}-f_{n_k}\|_p<+\infty.
	\]
	取一个处处正的连续函数 \(h\in L^q\). 由 H\"older 不等式,
	\[
	\sum_{k=1}^{\infty}
	\|h(f_{n_{k+1}}-f_{n_k})\|_1
	\leq
	\|h\|_q
	\sum_{k=1}^{\infty}\|f_{n_{k+1}}-f_{n_k}\|_p
	<+\infty.
	\]
	由 \(L^1\) 情形的级数定理, 级数
	\[
	h f_{n_1}
	+
	\sum_{k=1}^{\infty}h(f_{n_{k+1}}-f_{n_k})
	\]
	几乎处处绝对收敛. 因 \(h(x)>0\), 可知 \(f_{n_k}(x)\) 对几乎所有 \(x\) 收敛. 记其极限为 \(f(x)\). 于是 \(f\) 可测.
	
	给定 \(\varepsilon>0\). 由 \eqref{eq:lp-cauchy-condition}, 可取 \(N\), 使得当 \(m,n\geq N\) 时
	\[
	\|f_m-f_n\|_p<\varepsilon.
	\]
	固定 \(m\geq N\), 令 \(k\to\infty\), 并使用 Fatou 引理, 得
	\[
	\int_{\R^d}|f_m(x)-f(x)|^p\,\dd x
	\leq
	\liminf_{k\to\infty}
	\int_{\R^d}|f_m(x)-f_{n_k}(x)|^p\,\dd x
	\leq
	\varepsilon^p.
	\]
	因此 \(f_m-f\in L^p\), 且 \(\|f_m-f\|_p\leq\varepsilon\) 对所有 \(m\geq N\) 成立. 由于 \(f_m\in L^p\), 得 \(f\in L^p\), 并且 \(f_m\to f\) 于 \(L^p\) 中.
\end{proof}

\begin{rem}
	定理 \ref{thm:riesz-fischer-lp} 表明, \(L^p\) 在范数 \(\|\cdot\|_p\) 下是完备的. 因而 \(1\leq p<+\infty\) 时, \(L^p\) 是 Banach 空间.
\end{rem}

最后引入 \(L^\infty\) 空间.

\begin{defn}[\(L^\infty\) 空间]
	\label{def:l-infinity}
	若 \(f\) 是可测函数, 且存在常数 \(C\geq0\), 使得
	\[
	|f(x)|\leq C
	\]
	几乎处处成立, 则称 \(f\in L^\infty(\R^d)\). 定义
	\[
	\|f\|_\infty
	=
	\inf\{C\geq0: |f(x)|\leq C\text{ 几乎处处成立}\}.
	\]
	这个数称为 \(f\) 的本质上确界范数.
\end{defn}

实际上, 上述下确界可以达到. 也就是说,
\[
|f(x)|\leq \|f\|_\infty
\]
几乎处处成立. 这是因为可取 \(C_n\downarrow \|f\|_\infty\), 且 \(|f|\leq C_n\) 几乎处处成立; 去掉这些例外零集的可数并后, 令 \(n\to\infty\) 即可.

\begin{rem}
	H\"older 不等式在 \(p=1,q=\infty\) 的情形仍然成立: 若 \(f\in L^1\), \(g\in L^\infty\), 则 \(fg\in L^1\), 且
	\[
	\int_{\R^d}|f(x)g(x)|\,\dd x
	\leq
	\|g\|_\infty\int_{\R^d}|f(x)|\,\dd x.
	\]
\end{rem}

\begin{exer}
	\label{exer:lp-limit-to-linfty}
	假定 \(f\in L^\infty\), 并且 \(f\in L^p\) 对任意 \(1\leq p<+\infty\) 成立. 证明
	\[
	\|f\|_\infty
	=
	\lim_{p\to+\infty}\|f\|_p.
	\]
\end{exer}

\clearpage

\section{勒贝格积分的微商}
\label{sec:derivatives-of-lebesgue-integrals}

本节证明勒贝格积分的一个基本定理: 局部可积函数几乎处处可以由它在小球上的平均值恢复. 这个结果是积分理论在三角级数、调和函数以及许多分析问题中的基础.

若 \(S\subset\R^d\) 是球, 记 \(m(S)\) 为它的 Lebesgue 测度. 对局部可积函数 \(f\), 我们关心如下平均值:
\[
\frac{1}{m(S)}\int_S f(y)\,\dd y,
\]
其中 \(S\) 是包含点 \(x\) 的球, 并且 \(m(S)\to0\).

\begin{thm}[Lebesgue 微分定理]
	\label{thm:lebesgue-differentiation-balls}
	若 \(f\in L^1_{\operatorname{loc}}(\R^d)\), 则对几乎所有的 \(x\in\R^d\), 有
	\begin{equation}
		f(x)=\lim_{\substack{m(S)\to0\\ x\in S}}
		\frac{1}{m(S)}\int_S f(y)\,\dd y,
		\label{eq:lebesgue-differentiation-average}
	\end{equation}
	其中 \(S\) 取遍包含 \(x\) 的球. 进一步, 对几乎所有的 \(x\), 还有
	\begin{equation}
		\lim_{\substack{m(S)\to0\\ x\in S}}
		\frac{1}{m(S)}\int_S |f(y)-f(x)|\,\dd y=0.
		\label{eq:lebesgue-point-property}
	\end{equation}
\end{thm}

为了证明这个定理, 我们引入 Hardy--Littlewood 最大函数. 对 \(f\in L^1_{\operatorname{loc}}(\R^d)\), 定义
\begin{equation}
	Mf(x)=\sup_{S\ni x}\frac{1}{m(S)}\int_S |f(y)|\,\dd y,
	\label{eq:hardy-littlewood-maximal-function}
\end{equation}
其中上确界取遍所有包含 \(x\) 的球 \(S\).

\begin{exer}\label{exer:maximal-function-lsc}
	证明 \(Mf\) 是下方半连续函数, 即 \(Mf\in\mathcal I^+\).
\end{exer}

\begin{thm}[Hardy 最大定理]
	\label{thm:hardy-maximal-weak-l1}
	若 \(f\in L^1(\R^d)\), 则对任意 \(\alpha>0\), 有
	\begin{equation}
		m\{x:Mf(x)>\alpha\}
		\leq
		\frac{4^d}{\alpha}\int_{\R^d}|f(y)|\,\dd y.
		\label{eq:hardy-maximal-weak-l1}
	\end{equation}
\end{thm}

\begin{proof}
	记
	\[
	E_\alpha=\{x:Mf(x)>\alpha\}.
	\]
	由定义, \(E_\alpha\) 是一族球的并. 这些球满足
	\[
	\frac{1}{m(S)}\int_S |f(y)|\,\dd y>\alpha.
	\]
	记这族球为 \(\mathcal F\). 这些球的半径有上界, 因为若 \(S\in\mathcal F\), 则
	\[
	\alpha m(S)<\int_S |f(y)|\,\dd y\leq \int_{\R^d}|f(y)|\,\dd y.
	\]
	
	我们从 \(\mathcal F\) 中选出一列互不相交的球 \(S_k\), 使得 \(E_\alpha\) 被这些球的四倍同心放大球覆盖. 设 \(R_1\) 是 \(\mathcal F\) 中球半径的上确界, 取 \(S_1\in\mathcal F\), 使其半径 \(r_1>2R_1/3\). 若已经取出互不相交的 \(S_1,\ldots,S_{k-1}\), 令 \(R_k\) 为 \(\mathcal F\) 中所有与 \(S_1,\ldots,S_{k-1}\) 都不相交的球的半径上确界. 若 \(R_k=0\), 选取过程停止; 否则取 \(S_k\in\mathcal F\), 使其半径 \(r_k>2R_k/3\).
	
	若这个过程无限进行, 则由于 \(S_k\) 互不相交, 有
	\[
	\sum_k m(S_k)
	<\frac{1}{\alpha}\sum_k\int_{S_k}|f(y)|\,\dd y
	\leq \frac{1}{\alpha}\int_{\R^d}|f(y)|\,\dd y<+\infty.
	\]
	因此 \(m(S_k)\to0\), 也就是 \(r_k\to0\), 从而 \(R_k\to0\).
	
	现在取 \(x\in E_\alpha\). 则存在 \(S\in\mathcal F\) 使 \(x\in S\). 设 \(S\) 的半径为 \(r\). 若 \(S\) 与所有 \(S_k\) 都不相交, 则 \(R_k\geq r\) 对所有 \(k\) 成立, 与 \(R_k\to0\) 矛盾. 因此 \(S\) 必与某个 \(S_k\) 相交. 取第一个与 \(S\) 相交的 \(S_k\). 此时 \(S\) 与 \(S_1,\ldots,S_{k-1}\) 不相交, 所以 \(r\leq R_k<3r_k/2\). 若 \(a_k\) 为 \(S_k\) 的中心, 则
	\[
	\|x-a_k\|\leq 2r+r_k<4r_k.
	\]
	因此 \(x\) 属于 \(S_k\) 的四倍同心放大球. 于是
	\[
	E_\alpha\subset \bigcup_k 4S_k.
	\]
	从而
	\[
	m(E_\alpha)
	\leq \sum_k m(4S_k)
	=4^d\sum_k m(S_k)
	\leq \frac{4^d}{\alpha}\int_{\R^d}|f(y)|\,\dd y.
	\]
\end{proof}

\begin{proof}[定理 \ref{thm:lebesgue-differentiation-balls} 的证明]
	结论是局部的, 因此可以先证明 \(f\in L^1(\R^d)\) 的情形. 一般的局部可积情形, 只需在任意紧集附近乘以一个等于 \(1\) 的紧支连续截断函数, 再把结论局部化即可.
	
	对 \(\varepsilon>0\), 令
	\[
	E_\varepsilon(f)=\left\{x:
	\limsup_{\substack{m(S)\to0\\x\in S}}
	\frac{1}{m(S)}\int_S |f(y)-f(x)|\,\dd y>\varepsilon
	\right\}.
	\]
	我们证明每个 \(E_\varepsilon(f)\) 都是零集. 任取 \(g\in\C_0(\R^d)\). 因为 \(g\) 连续, 所以
	\[
	\frac{1}{m(S)}\int_S |g(y)-g(x)|\,\dd y\longrightarrow0
	\]
	当 \(m(S)\to0\) 且 \(x\in S\). 由三角不等式,
	\begin{align*}
		\frac{1}{m(S)}\int_S |f(y)-f(x)|\,\dd y
		&\leq
		\frac{1}{m(S)}\int_S |f(y)-g(y)|\,\dd y\nobreak\\
		&\quad+
		\frac{1}{m(S)}\int_S |g(y)-g(x)|\,\dd y
		+|g(x)-f(x)|.
	\end{align*}
	因此
	\[
	E_{3\varepsilon}(f)
	\subset
	\{x:M(f-g)(x)>\varepsilon\}
	\cup
	\{x:|f(x)-g(x)|>\varepsilon\}.
	\]
	由 Hardy 最大定理和 Chebyshev 不等式, 得
	\[
	m^*(E_{3\varepsilon}(f))
	\leq
	\frac{4^d+1}{\varepsilon}
	\int_{\R^d}|f(x)-g(x)|\,\dd x.
	\]
	由于 \(\C_0(\R^d)\) 在 \(L^1\) 中稠密, 可令右端任意小. 所以
	\[
	m^*(E_{3\varepsilon}(f))=0.
	\]
	令 \(\varepsilon\) 取遍正有理数, 即得 \eqref{eq:lebesgue-point-property} 对几乎所有 \(x\) 成立. 由
	\[
	\left|\frac{1}{m(S)}\int_S f(y)\,\dd y-f(x)\right|
	\leq
	\frac{1}{m(S)}\int_S |f(y)-f(x)|\,\dd y,
	\]
	立即得到 \eqref{eq:lebesgue-differentiation-average}.
\end{proof}

\begin{exer}\label{exer:differentiation-general-bases}
	证明定理 \ref{thm:hardy-maximal-weak-l1} 和定理 \ref{thm:lebesgue-differentiation-balls} 在下述情形仍成立: 不只考虑包含 \(x\) 的球 \(S\), 而考虑所有满足以下条件的可测集合 \(E\): 对某个固定常数 \(K\), 存在一个球 \(S\supset E\cup\{x\}\), 使得
	\[
	m(S)\leq K m(E).
	\]
\end{exer}

\begin{exer}\label{exer:approximate-identity-l1}
	设 \(f\in L^1(\R^d)\), 并设 \(g\in L^1\) 有界且关于 \(\|x\|\) 单调下降. 令
	\[
	g_\varepsilon(x)=\varepsilon^{-d}g(x/\varepsilon),
	\qquad \varepsilon>0.
	\]
	证明
	\[
	\lim_{\varepsilon\to0}
	\int_{\R^d}\left|f(x-y)g_\varepsilon(y)-f(x)g_\varepsilon(y)\right|\,\dd y=0
	\]
	对几乎所有的 \(x\) 成立. 并证明左端作为 \(x\) 的函数是连续的. 提示: 先考虑 \(g\) 具有紧致支集的情形, 并利用定理 \ref{thm:hardy-maximal-weak-l1}.
\end{exer}

对于 \(L^p\) 函数, 可以把 Hardy 最大定理的弱型估计改进为强型估计.

\begin{thm}[最大函数的 \(L^p\) 估计]
	\label{thm:maximal-lp-estimate}
	若 \(f\in L^p(\R^d)\), \(1<p<+\infty\), 则
	\begin{equation}
		\int_{\R^d}(Mf(x))^p\,\dd x
		\leq
		C_p\int_{\R^d}|f(x)|^p\,\dd x,
		\label{eq:maximal-lp-estimate}
	\end{equation}
	其中可以取
	\[
	C_p=4^d2^p\frac{p}{p-1}.
	\]
\end{thm}

\begin{proof}
	设
	\[
	E_s=\{x:Mf(x)>s\},
	\qquad
	\Phi(s)=m(E_s).
	\]
	由分布函数公式,
	\[
	\int_{\R^d}(Mf(x))^p\,\dd x
	=
	\int_0^{+\infty}\Phi(s)\,\dd(s^p).
	\]
	对每个 \(s>0\), 将 \(f\) 分解为
	\[
	f=g_s+h_s,
	\]
	其中
	\[
	g_s(x)=
	\begin{cases}
		f(x), & |f(x)|\leq s/2,\\
		0, & |f(x)|>s/2,
	\end{cases}
	\qquad
	h_s=f-g_s.
	\]
	则 \(Mg_s\leq s/2\). 因而在 \(E_s\) 上必有 \(Mh_s>s/2\). 由 Hardy 最大定理,
	\[
	\Phi(s)
	\leq
	m\{x:Mh_s(x)>s/2\}
	\leq
	\frac{2\cdot4^d}{s}\int_{\{|f|>s/2\}}|f(x)|\,\dd x.
	\]
	于是
	\begin{align*}
		\int_{\R^d}(Mf)^p\,\dd x
		&\leq
		2\cdot4^d\int_0^{+\infty}s^{-1}
		\left(\int_{\{|f|>s/2\}}|f(x)|\,\dd x\right)\dd(s^p)\\
		&=2\cdot4^d p
		\int_{\R^d}|f(x)|
		\left(\int_0^{2|f(x)|}s^{p-2}\,\dd s\right)\dd x\\
		&=4^d2^p\frac{p}{p-1}
		\int_{\R^d}|f(x)|^p\,\dd x.
	\end{align*}
	这就是所需估计.
\end{proof}

\begin{rem}
	上面的常数并不是最佳的. 但是, 通过取适当的非负函数, 可以看出当 \(p\downarrow1\) 时, 常数关于 \((p-1)^{-1}\) 的阶不能改进. 特别地, 当 \(p=1\) 时不存在形如 \eqref{eq:maximal-lp-estimate} 的强型 \(L^1\) 不等式; 此时只能得到定理 \ref{thm:hardy-maximal-weak-l1} 的弱型估计.
\end{rem}

\begin{thm}[区间基下的微分定理]
	\label{thm:lebesgue-differentiation-rectangles-lp}
	设 \(f\in L^p(\R^d)\), \(1<p<+\infty\). 则对几乎所有 \(x\in\R^d\), 有
	\begin{equation}
		f(x)=\lim_{\substack{\operatorname{diam} I\to0\\ x\in I}}
		\frac{1}{m(I)}\int_I f(y)\,\dd y,
		\label{eq:rectangle-differentiation-lp}
	\end{equation}
	其中 \(I\) 取遍包含 \(x\) 的区间.
\end{thm}

\begin{proof}
	定义区间最大函数
	\[
	M_{\Box}f(x)=\sup_{I\ni x}
	\frac{1}{m(I)}\int_I |f(y)|\,\dd y,
	\]
	其中 \(I\) 取遍包含 \(x\) 的区间. 当 \(d=1\) 时, 它与前面的球最大函数只差无关紧要的常数. 一般维数下, 可以按坐标方向逐次应用一维最大算子, 得到
	\[
	\|M_{\Box}f\|_p\leq K_{p,d}\|f\|_p.
	\]
	于是有
	\[
	m\{x:M_{\Box}f(x)>s\}
	\leq
	K_{p,d}^p s^{-p}\|f\|_p^p.
	\]
	现在重复定理 \ref{thm:lebesgue-differentiation-balls} 的证明, 把球换成区间, 把 Hardy 最大定理换成上面的 \(L^p\) 控制, 即得结论.
\end{proof}

\begin{rem}
	定理 \ref{thm:lebesgue-differentiation-rectangles-lp} 中对区间各边长的比例没有任何限制. 这点与只用球作平均时不同. 当 \(p=1\) 且 \(d>1\) 时, 这个结论一般不成立.
\end{rem}

\begin{exer}\label{exer:convex-neighborhood-differentiation}
	设 \(\mathcal F\) 是原点的一族闭的、有界的、凸的、对称邻域, 满足
	\[
	\bigcap_{S\in\mathcal F}S=\{0\},
	\]
	并且对任意 \(S,S'\in\mathcal F\), 有 \(S\subset S'\) 或 \(S'\subset S\). 证明: 若 \(f\in L^1(\R^d)\), 则
	\[
	\lim_{\substack{S\in\mathcal F\\ m(S)\to0}}
	\frac{1}{m(S)}\int_S |f(x+y)-f(x)|\,\dd y=0
	\]
	对几乎所有的 \(x\) 成立.
	
	提示: 先把族 \(\mathcal F\) 扩大, 使得若 \(S_n\in\mathcal F\), \(S_n\downarrow S\), 则 \(S\in\mathcal F\). 注意若 \(S,S'\in\mathcal F\), \(S'\subset S\), 且 \((x+S)\cap(x'+S')\neq\varnothing\), 则
	\[
	x'+S'\subset x+3S.
	\]
	然后把定理 \ref{thm:hardy-maximal-weak-l1} 推广到最大函数
	\[
	M_{\mathcal F}f(x)=\sup_{S\in\mathcal F}
	\frac{1}{m(S)}\int_S |f(x+y)|\,\dd y.
	\]
\end{exer}

现在利用定理 \ref{thm:lebesgue-differentiation-balls} 和定理 \ref{thm:hardy-maximal-weak-l1}, 证明实直线上的单调函数几乎处处可微. 虽然这个结果也可以由第三章的一般理论更自然地推出, 但直接证明本身也很有意义.

设 \(f\) 是 \(\R\) 上的增函数. 对区间 \(I=[a,b]\), 记
\[
f(I)=f(b)-f(a),
\qquad
m(I)=b-a.
\]
定义 \(f\) 的下微商和上微商为
\[
f_*(x)=\liminf_{\substack{m(I)\to0\\x\in I}}\frac{f(I)}{m(I)},
\qquad
f^*(x)=\limsup_{\substack{m(I)\to0\\x\in I}}\frac{f(I)}{m(I)}.
\]
容易看出 \(f_*\in L^1_{\operatorname{loc}}(\R)\). 进一步, 函数
\[
g(x)=f(x)-\int_0^x f_*(t)\,\dd t
\]
仍是增函数. 当 \(x<y\) 时, 可由 Fatou 引理得到
\[
\int_x^y f_*(t)\,\dd t\leq f(y)-f(x),
\]
这正说明 \(g(y)-g(x)\geq0\).

由定理 \ref{thm:lebesgue-differentiation-balls}, 函数 \(x\mapsto\int_0^x f_*(t)\,\dd t\) 的微商几乎处处等于 \(f_*(x)\). 因而 \(g_*(x)=0\) 几乎处处. 我们还要证明 \(g^*(x)=0\) 几乎处处.

为此, 对任意 \(\varepsilon>0\), 考虑集合
\[
E=\{x:g_*(x)=0\text{ 且 }g^*(x)>\varepsilon\}.
\]
下面证明 \(E\) 是零集. 若 \(\delta>0\) 且 \(x\in E\), 则可以找到包含 \(x\) 的区间 \(I=[a,b]\), 使得
\[
g(I)<\delta m(I).
\]
定义一个增函数 \(G\): 在 \(I\) 内令 \(G=g\), 在 \(a\) 左侧令 \(G(t)=g(a)\), 在 \(b\) 右侧令 \(G(t)=g(b)\). 对 \(G\) 定义相应的最大微商
\[
H(x)=\sup_{J\ni x}\frac{G(J)}{m(J)}.
\]
把 Hardy 最大定理的证明稍作修改, 可得对任意增函数 \(G\), 有
\[
m\{x:H(x)>\varepsilon\}\leq \frac{4}{\varepsilon}\bigl(G(+\infty)-G(-\infty)\bigr).
\]
这里对上面这个 \(G\), 右端等于 \(4g(I)/\varepsilon\), 从而
\[
m(I\cap E)
\leq \frac{4g(I)}{\varepsilon}
\leq \frac{4\delta}{\varepsilon}m(I).
\]
因为 \(\delta>0\) 任意, 可知 \(E\) 的特征函数在 \(E\) 内几乎处处具有微商零, 从而 \(E\) 必为零集. 于是 \(g^*(x)=0\) 几乎处处.

由此得到如下定理.

\begin{thm}\label{thm:monotone-function-ae-differentiable}
	若 \(f\) 是 \(\R\) 上的增函数, 则 \(f'(x)\) 几乎处处存在, 并且 \(f'\in L^1_{\operatorname{loc}}(\R)\). 此外, 可以写成
	\begin{equation}
		f(y)-f(x)=\int_x^y f'(t)\,\dd t+g(y)-g(x),
		\label{eq:monotone-decomposition}
	\end{equation}
	其中 \(g\) 是增函数, 且 \(g'(x)=0\) 几乎处处.
\end{thm}

显然, \(g\) 不必是常数. 例如阶梯函数就是增函数, 且它的导数几乎处处为零. 更有意思的是, 存在连续的增函数, 它的导数几乎处处等于零. 下面给出经典例子.

\begin{exam}[Cantor 函数]
	\label{exam:cantor-function}
	我们在 \([0,1]\) 上构造一个连续的增函数 \(g\), 使得
	\[
	g'(x)=0
	\]
	在 Cantor 集 \(C\) 的余集上成立. 由于 \(C\) 是零集, 这就给出一个连续增函数, 其导数几乎处处为零, 但函数本身并非常数.
	
	设 \(E_n\) 是构造 Cantor 集时第 \(n\) 步剩下的闭集, 并令
	\[
	O_n=[0,1]\setminus E_n.
	\]
	则 \(O_n\) 由 \(2^n-1\) 个开区间组成, 它们正是前 \(n\) 步删去的那些中间三分之一开区间. 在 \(O_n\) 上定义函数 \(g_n\): 若 \(x\) 属于 \(O_n\) 中从左到右数的第 \(k\) 个区间, 则令
	\[
	g_n(x)=\frac{k}{2^n}.
	\]
	由于 \(O_n\subset O_{n+1}\), 并且 \(g_{n+1}\) 在 \(O_n\) 上与 \(g_n\) 一致, 可以在
	\[
	O=[0,1]\setminus C=\bigcup_{n=1}^{\infty}O_n
	\]
	上定义函数 \(g\), 即
	\[
	g(x)=g_n(x),\qquad x\in O_n.
	\]
	由构造可知
	\[
	|g(x)-g(y)|\leq 2^{-n},
	\qquad |x-y|<3^{-n}.
	\]
	因此 \(g\) 在 \(O\) 上一致连续, 从而可以唯一地连续延拓到 \([0,1]\) 上. 这个延拓仍记为 \(g\). 容易看出 \(g\) 是不减函数. 若 \(x\in O\), 则对充分大的 \(n\), 点 \(x\) 落在 \(O_n\) 的某个开区间内, 而在该区间上 \(g=g_n\) 为常数, 因此
	\[
	g'(x)=0.
	\]
	因为 \([0,1]\setminus O=C\) 是零集, 所以 \(g'(x)=0\) 几乎处处.
\end{exam}
	
	\chapter{拉东--斯蒂尔杰斯积分}

\section{正测度的定义}
\label{sec:positive-radon-measures}

在第二章中, 我们已经看到: 从黎曼积分推广到 Lebesgue 积分时, 真正使用的并不是黎曼积分的全部性质, 而主要是 \(\C_0(\R^d)\) 上积分的两个基本性质:
\begin{align}
	\int_{\R^d}(af+bg)\,\dd x
	&=a\int_{\R^d}f\,\dd x+b\int_{\R^d}g\,\dd x,\label{eq:lebesgue-basic-linear}\\
	\int_{\R^d}f\,\dd x&\geq0,\qquad f\in\C_0^+(\R^d).\label{eq:lebesgue-basic-positive}
\end{align}
因此, 可以把 Lebesgue 积分的构造抽象出来: 不一定从通常的体积积分出发, 而是从 \(\C_0(\R^d)\) 上的一个正线性泛函出发. 这样得到的积分就是 Radon 积分的基本形式.

设
\[
\mu:\C_0(\R^d)\longrightarrow\R
\]
是一个映射. 若它满足
\begin{equation}
	\mu(af+bg)=a\mu(f)+b\mu(g),
	\qquad a,b\in\R,\quad f,g\in\C_0(\R^d),
	\label{eq:positive-measure-linear}
\end{equation}
以及
\begin{equation}
	\mu(f)\geq0,
	\qquad f\in\C_0^+(\R^d),
	\label{eq:positive-measure-positive}
\end{equation}
则 \(\mu\) 称为 \(\C_0(\R^d)\) 上的一个正线性泛函.

由 \eqref{eq:positive-measure-linear} 和 \eqref{eq:positive-measure-positive} 立刻得到单调性: 若 \(f,g\in\C_0(\R^d)\) 且 \(f\leq g\), 则
\[
\mu(f)\leq \mu(g).
\]
此外还有下面的连续性性质.

\begin{lem}\label{lem:positive-functional-continuity-from-above}
	设 \(\mu\) 满足 \eqref{eq:positive-measure-linear} 和 \eqref{eq:positive-measure-positive}. 若 \(f_n\in\C_0^+(\R^d)\), 且
	\[
	f_1\geq f_2\geq\cdots\geq0,
	\qquad
	f_n(x)\downarrow0
	\]
	对每个 \(x\in\R^d\) 成立, 并且所有 \(f_n\) 的支集都包含在同一个紧集内, 则
	\begin{equation}
		\mu(f_n)\longrightarrow0.
		\label{eq:positive-functional-continuity-from-above}
	\end{equation}
\end{lem}

\begin{proof}
	因为所有 \(f_n\) 的支集都包含在同一个紧集 \(K\) 中, 可取 \(\varphi\in\C_0^+(\R^d)\), 使得
	\[
	\varphi=1\quad\text{于 }K\text{ 上}.
	\]
	令
	\[
	M_n=\max_{x\in K} f_n(x).
	\]
	由 Dini 引理, \(f_n\downarrow0\) 在紧集 \(K\) 上一致成立, 所以 \(M_n\to0\). 又有
	\[
	0\leq f_n\leq M_n\varphi.
	\]
	由单调性和齐次性,
	\[
	0\leq\mu(f_n)\leq M_n\mu(\varphi)\to0.
	\]
\end{proof}

现在可以仿照第二章的构造, 从 \(\mu\) 出发定义积分. 对 \(f\in\mathcal I^+\), 定义
\begin{equation}
	\mu^*(f)=\sup\{\mu(g):g\in\C_0(\R^d),\ 0\leq g\leq f\}.
	\label{eq:mu-star-on-i-plus}
\end{equation}
若 \(f\geq0\) 是任意非负函数, 再定义
\begin{equation}
	\mu^*(f)=\inf\{\mu^*(h):h\in\mathcal I^+,\ f\leq h\}.
	\label{eq:mu-star-positive-function}
\end{equation}
这里允许取值为 \(+\infty\). 与第二章完全相同, 可以证明 \(\mu^*\) 对非负函数具有单调性、正齐次性、次可加性和单调收敛性质.

\begin{defn}[关于 \(\mu\) 的可积函数]
	\label{def:l1-mu}
	设 \(f\) 是取值于 \([-\infty,+\infty]\) 的函数. 若对任意 \(\varepsilon>0\), 存在 \(g\in\C_0(\R^d)\), 使得
	\[
	\mu^*(|f-g|)<\varepsilon,
	\]
	则称 \(f\) 关于 \(\mu\) 可积, 并记 \(f\in L^1_\mu\).
\end{defn}

若 \(f\in L^1_\mu\), 取 \(g_n\in\C_0(\R^d)\), 使得
\[
\mu^*(|f-g_n|)\to0.
\]
则 \(\mu(g_n)\) 是 Cauchy 数列. 事实上,
\[
|\mu(g_n)-\mu(g_m)|
\leq \mu^*(|g_n-g_m|)
\leq \mu^*(|g_n-f|)+\mu^*(|f-g_m|)\to0.
\]
因此 \(\mu(g_n)\) 有极限, 且与逼近列 \(g_n\) 的选取无关.

\begin{defn}[关于 \(\mu\) 的积分]
	\label{def:integral-with-respect-to-mu}
	设 \(f\in L^1_\mu\). 取 \(g_n\in\C_0(\R^d)\), 使 \(\mu^*(|f-g_n|)\to0\). 定义
	\[
	\mu(f)=\lim_{n\to\infty}\mu(g_n).
	\]
	在不致混淆时, 也写作
	\[
	\int f\,\dd\mu.
	\]
\end{defn}

用这种方法, 第二章关于 Lebesgue 积分的许多结论都可以逐字推广到 \(\mu\)-积分. 例如, \(L^1_\mu\) 是线性空间; 若 \(f,g\in L^1_\mu\), \(a,b\in\R\), 则
\[
\mu(af+bg)=a\mu(f)+b\mu(g).
\]
若 \(f\geq0\) 且 \(f\in L^1_\mu\), 则
\[
\mu(f)\geq0.
\]
单调收敛定理、Fatou 引理和控制收敛定理也都有相应版本. 证明与第二章中的证明没有本质差别, 只要把通常积分 \(\int f\,\dd x\) 换成 \(\mu(f)\) 即可.

\begin{rem}
	需要注意的是, 零集合的概念现在依赖于 \(\mu\). 也就是说, 一个集合 \(E\) 是否满足
	\[
	\mu^*(\chi_E)=0
	\]
	取决于所选的正线性泛函 \(\mu\). 因此, 在涉及``几乎处处''的表述时, 必须说明是相对于哪个测度而言.
\end{rem}

\begin{defn}[Radon 正测度]
	\label{def:positive-radon-measure}
	一个满足 \eqref{eq:positive-measure-linear} 和 \eqref{eq:positive-measure-positive} 的正线性泛函
	\[
	\mu:
	\C_0(\R^d)\to\R
	\]
	称为 \(\R^d\) 上的一个 Radon 正测度.
\end{defn}

特别地, 通常的 Lebesgue 测度就是由
\[
\mu(f)=\int_{\R^d}f(x)\,\dd x,
\qquad f\in\C_0(\R^d),
\]
所给出的 Radon 正测度.

更一般地, 若 \(\Omega\subset\R^d\) 是开集, 用 \(\C_0(\Omega)\) 代替 \(\C_0(\R^d)\), 则 \(\C_0(\Omega)\) 上的正线性泛函称为 \(\Omega\) 上的 Radon 正测度. 为了避免记号过于繁复, 后面主要讨论 \(\R^d\) 上的情形; 对开集 \(\Omega\) 的情形, 相应结论可以逐字改写.

\clearpage

\section{一维情形: Stieltjes 积分}
\label{sec:one-dimensional-stieltjes-integral}

设 \(\Omega=(a,b)\subset\R\) 是一个有限或无限开区间. 本节的目标是描述 \(\Omega\) 上所有 Radon 正测度. 结果表明, 一维 Radon 正测度正好就是 Stieltjes 积分.

对 \(x\in\Omega\), 定义函数
\[
H_x(y)=
\begin{cases}
	1, & y>x,\\
	0, & y\leq x.
\end{cases}
\]
则 \(H_x\in\mathcal I^+(\Omega)\). 若 \(x_1<x_2\), 则 \(H_{x_1}-H_{x_2}\) 是区间 \((x_1,x_2]\) 的特征函数, 因而关于任意 Radon 正测度 \(\mu\) 都是可积的.

设 \(\mu\) 是 \(\Omega\) 上的 Radon 正测度. 我们要构造一个增函数 \(\varphi\), 使得
\begin{equation}
	\varphi(x_2)-\varphi(x_1)
	=
	\mu(H_{x_1}-H_{x_2}),
	\qquad x_1<x_2.
	\label{eq:stieltjes-distribution-function}
\end{equation}
取定一点 \(x_0\in\Omega\) 和常数 \(C\in\R\), 定义
\begin{equation}
	\varphi(x)=\mu(H_{x_0}-H_x)+C.
	\label{eq:distribution-function-from-measure}
\end{equation}
这里若 \(x<x_0\), 则 \(H_{x_0}-H_x\) 应理解为 \(- (H_x-H_{x_0})\). 由积分的可加性可知, \eqref{eq:distribution-function-from-measure} 正好满足 \eqref{eq:stieltjes-distribution-function}. 反过来, 满足 \eqref{eq:stieltjes-distribution-function} 的函数只能相差一个常数.

由于 \(H_{x_1}-H_{x_2}\geq0\) 当 \(x_1<x_2\), 由 \(\mu\) 的正性可知 \(\varphi\) 是增函数. 进一步, \(\varphi\) 是右连续的. 事实上, 当 \(h\downarrow0\) 时,
\[
H_x-H_{x+h}\downarrow0
\]
在紧支意义下成立, 从而
\[
\varphi(x+h)-\varphi(x)
=
\mu(H_x-H_{x+h})\longrightarrow0.
\]

下面说明, \(\mu\) 可以由 \(\varphi\) 重新构造出来. 设 \(f\in\C_0(\Omega)\), 取一个包含 \(\operatorname{supp}f\) 的紧区间 \([\alpha,\beta]\subset\Omega\), 并作剖分
\[
\alpha=x_0<x_1<\cdots <x_N=\beta.
\]
记 \(M_k\) 和 \(m_k\) 分别为 \(f\) 在 \([x_{k-1},x_k]\) 上的最大值和最小值. 由 \eqref{eq:stieltjes-distribution-function} 和单调性, 有
\begin{equation}
	\sum_{k=1}^N m_k\bigl(\varphi(x_k)-\varphi(x_{k-1})\bigr)
	\leq
	\mu(f)
	\leq
	\sum_{k=1}^N M_k\bigl(\varphi(x_k)-\varphi(x_{k-1})\bigr).
	\label{eq:stieltjes-upper-lower-bound}
\end{equation}

反过来, 若 \(\varphi\) 是 \(\Omega\) 上的一个增函数, 则可以仿照 Riemann 积分定义 Riemann--Stieltjes 上、下积分:
\[
\overline{\int_\Omega} f\,\dd\varphi
=
\inf_P\sum_{k=1}^N M_k\bigl(\varphi(x_k)-\varphi(x_{k-1})\bigr),
\]
\[
\underline{\int_\Omega} f\,\dd\varphi
=
\sup_P\sum_{k=1}^N m_k\bigl(\varphi(x_k)-\varphi(x_{k-1})\bigr),
\]
其中 \(P\) 取遍包含 \(\operatorname{supp}f\) 的紧区间的所有剖分. 当 \(f\in\C_0(\Omega)\) 时, 由于 \(f\) 一致连续, 上、下 Stieltjes 积分相等. 这个共同值记作
\[
\int_\Omega f\,\dd\varphi.
\]
于是由 \eqref{eq:stieltjes-upper-lower-bound} 得
\[
\mu(f)=\int_\Omega f\,\dd\varphi,
\qquad f\in\C_0(\Omega).
\]
另一方面, 对任意增函数 \(\varphi\), 映射
\[
f\longmapsto \int_\Omega f\,\dd\varphi,
\qquad f\in\C_0(\Omega),
\]
显然是正线性的, 从而定义 \(\Omega\) 上的一个 Radon 正测度.

现在讨论唯一性. 设两个增函数 \(\varphi\) 与 \(\psi\) 给出同一个正测度, 即
\[
\int_\Omega f\,\dd\varphi
=
\int_\Omega f\,\dd\psi,
\qquad f\in\C_0(\Omega).
\]
取 \(x_1<x_2\). 令 \(f_n\in\C_0(\Omega)\) 满足 \(0\leq f_n\leq1\), 在 \([x_1,x_2]\) 上等于 \(1\), 在 \((x_1-1/n,x_2+1/n)^c\) 上等于 \(0\), 并在中间连续过渡. 令 \(n\to\infty\), 得
\[
\varphi(x_2+0)-\varphi(x_1-0)
=
\psi(x_2+0)-\psi(x_1-0).
\]
因此
\[
\varphi(x+0)-\psi(x+0)
=
\varphi(x-0)-\psi(x-0)=C
\]
对所有 \(x\in\Omega\) 成立, 其中 \(C\) 为常数. 特别地, 在所有连续点上, \(\varphi\) 与 \(\psi\) 只相差同一个常数. 若进一步要求二者右连续, 则
\[
\varphi(x)-\psi(x)=C,
\qquad x\in\Omega.
\]

于是得到如下定理.

\begin{thm}[Riesz 表示的一维形式]
	\label{thm:one-dimensional-positive-measure-stieltjes}
	设 \(\Omega=(a,b)\subset\R\). \(\Omega\) 上任意 Radon 正测度 \(\mu\) 都可以表示为关于某个增函数 \(\varphi\) 的 Stieltjes 积分:
	\[
	\mu(f)=\int_\Omega f\,\dd\varphi,
	\qquad f\in\C_0(\Omega).
	\]
	若要求 \(\varphi\) 右连续, 并指定它在某一点的值, 则 \(\varphi\) 由 \(\mu\) 唯一确定.
\end{thm}

\begin{exer}
	证明: 一个增函数至多有可数多个不连续点. 提示: 先证明对任意 \(\varepsilon>0\), 在任一紧区间上, 跃度大于 \(\varepsilon\) 的不连续点只有有限多个.
\end{exer}

由这个练习可知, 若只在可数多个点上改变一个增函数的取值, 则相应的 Stieltjes 积分不变. 因为紧支连续函数的 Stieltjes 积分只感受到增函数的左右极限所决定的跳跃量, 而不取决于函数在跳跃点本身如何赋值.

\begin{rem}
	一般的 Radon 测度可以看作 Stieltjes 积分的高维推广. 因此, 即使在多变量情形中, 我们也常把测度作用写成积分形式
	\[
	\mu(f)=\int f\,\dd\mu.
	\]
	这种记号强调的是: 测度本身充当了积分的“微分元”.
\end{rem}

\clearpage

\section{一般的 Radon 测度及其正部与负部的分解}
\label{sec:general-radon-measures-jordan-decomposition}

若 \(\mu\) 为 Radon 正测度, 则它具有如下连续性: 若 \(f_n\in\C_0(\R^d)\) 都在同一个紧集外等于零, 并且 \(f_n\) 一致收敛到零, 则
\[
\mu(f_n)\longrightarrow0.
\]
事实上, 取 \(\varphi\in\C_0^+(\R^d)\), 使得 \(\varphi=1\) 于这个公共紧集上. 若
\[
M_n=\sup_{x\in\R^d}|f_n(x)|,
\]
则 \(M_n\to0\), 且
\[
-M_n\varphi\leq f_n\leq M_n\varphi.
\]
由正性和单调性便得
\[
|\mu(f_n)|\leq M_n\mu(\varphi)\to0.
\]
基于这个性质, 我们可以把正测度推广为一般的 Radon 测度.

\begin{defn}[Radon 测度]
	\label{def:radon-measure-general}
	一个定义在 \(\C_0(\R^d)\) 上的实值泛函 \(\mu\), 若满足:
	\begin{enumerate}
		\item 线性性:
		\[
		\mu(af+bg)=a\mu(f)+b\mu(g),
		\qquad a,b\in\R,\quad f,g\in\C_0(\R^d);
		\]
		\item 连续性: 若 \(f_n\in\C_0(\R^d)\) 都在同一个紧集外等于零, 且 \(f_n\) 一致收敛到零, 则
		\[
		\mu(f_n)\to0,
		\]
	\end{enumerate}
	则称 \(\mu\) 为 \(\R^d\) 上的一个 Radon 测度.
\end{defn}

Radon 正测度显然是这里意义下的 Radon 测度. 引入一般 Radon 测度之后, 测度的线性组合也自然有意义: 若 \(\mu,\nu\) 为 Radon 测度, \(a,b\in\R\), 则
\[
(a\mu+b\nu)(f)=a\mu(f)+b\nu(f),
\qquad f\in\C_0(\R^d),
\]
定义了新的 Radon 测度.

\begin{defn}[测度的序]
	\label{def:order-of-radon-measures}
	设 \(\mu,\nu\) 为 Radon 测度. 若 \(\nu-\mu\) 是 Radon 正测度, 即
	\[
	\mu(f)\leq \nu(f),
	\qquad f\in\C_0^+(\R^d),
	\]
	则记作
	\[
	\mu\leq \nu.
	\]
\end{defn}

下面的定理是测度的 Jordan 分解. 它是函数分解
\[
f=f^+-f^-,
\qquad
f^+=\max(f,0),
\quad
f^-=\max(-f,0)
\]
在测度层面的对应物.

\begin{thm}[Jordan 分解]
	\label{thm:jordan-decomposition-radon-measure}
	任意 Radon 测度 \(\mu\) 都可以唯一地写成
	\[
	\mu=\mu^+-\mu^-,
	\]
	其中 \(\mu^+\) 和 \(\mu^-\) 是 Radon 正测度. 这个分解具有如下最小性: 若
	\[
	\mu=\mu_1-\mu_2,
	\qquad \mu_1,\mu_2\geq0,
	\]
	则
	\[
	\mu^+\leq \mu_1,
	\qquad
	\mu^-\leq \mu_2.
	\]
\end{thm}

\begin{proof}
	先构造 \(\mu^+\). 对 \(f\in\C_0^+(\R^d)\), 定义
	\begin{equation}
		\mu^+(f)=\sup\{\mu(g):g\in\C_0(\R^d),\ 0\leq g\leq f\}.
		\label{eq:positive-part-measure-definition}
	\end{equation}
	显然 \(\mu^+(f)\geq0\), 且 \(\mu^+(cf)=c\mu^+(f)\) 对 \(c\geq0\) 成立.
	
	下面证明可加性. 设 \(f_1,f_2\in\C_0^+(\R^d)\). 一方面, 若 \(0\leq g_i\leq f_i\), 则
	\[
	0\leq g_1+g_2\leq f_1+f_2,
	\]
	所以
	\[
	\mu(g_1)+\mu(g_2)=\mu(g_1+g_2)\leq \mu^+(f_1+f_2).
	\]
	取上确界得
	\[
	\mu^+(f_1)+\mu^+(f_2)\leq \mu^+(f_1+f_2).
	\]
	另一方面, 任取 \(g\in\C_0(\R^d)\) 满足 \(0\leq g\leq f_1+f_2\). 令
	\[
	g_1=\min(g,f_1),
	\qquad
	g_2=g-g_1.
	\]
	则 \(g_1,g_2\in\C_0^+(\R^d)\), 且 \(g_1\leq f_1\). 又
	\[
	g_2=g-\min(g,f_1)=\max(0,g-f_1)\leq f_2.
	\]
	因此
	\[
	\mu(g)=\mu(g_1)+\mu(g_2)\leq \mu^+(f_1)+\mu^+(f_2).
	\]
	再对 \(g\) 取上确界, 得反向不等式. 所以 \(\mu^+\) 在 \(\C_0^+\) 上可加.
	
	由齐次性和可加性, \(\mu^+\) 可以唯一地扩张为 \(\C_0(\R^d)\) 上的正线性泛函. 由正线性泛函的连续性可知, \(\mu^+\) 是 Radon 正测度.
	
	定义
	\[
	\mu^-=\mu^+-\mu.
	\]
	若 \(f\in\C_0^+(\R^d)\), 则由 \eqref{eq:positive-part-measure-definition} 中取 \(g=f\) 可知
	\[
	\mu^+(f)\geq \mu(f),
	\]
	从而 \(\mu^-(f)\geq0\). 因此 \(\mu^-\) 也是 Radon 正测度, 且
	\[
	\mu=\mu^+-\mu^-.
	\]
	
	现在证明最小性. 若 \(\mu=\mu_1-\mu_2\), 其中 \(\mu_1,\mu_2\geq0\), 则对任意 \(0\leq g\leq f\) 有
	\[
	\mu(g)=\mu_1(g)-\mu_2(g)\leq \mu_1(g)\leq \mu_1(f).
	\]
	取上确界得
	\[
	\mu^+(f)\leq \mu_1(f),
	\qquad f\in\C_0^+(\R^d).
	\]
	于是 \(\mu^+\leq\mu_1\). 又
	\[
	\mu^-=\mu^+-\mu\leq \mu_1-\mu=\mu_2.
	\]
	最小性得证.
	
	唯一性也由最小性得到. 若 \(\mu=\alpha-\beta\) 是另一个具有相同最小性的分解, 则由最小性有 \(\mu^+\leq\alpha\) 和 \(\alpha\leq\mu^+\), 故 \(\alpha=\mu^+\); 同理 \(\beta=\mu^-\).
\end{proof}

\begin{defn}[全变差测度]
	\label{def:total-variation-measure}
	设 \(\mu\) 为 Radon 测度, 并设
	\[
	\mu=\mu^+-\mu^-
	\]
	是它的 Jordan 分解. 定义 \(\mu\) 的绝对值或全变差测度为
	\[
	|\mu|=\mu^++\mu^-.
	\]
\end{defn}

\begin{exer}
	\label{exer:total-variation-characterization}
	设 \(f\in\C_0^+(\R^d)\). 证明
	\[
	|\mu|(f)=\sup\{\mu(g):g\in\C_0(\R^d),\ |g|\leq f\}.
	\]
	并证明三角不等式
	\[
	|\mu+\nu|\leq |\mu|+|\nu|.
	\]
\end{exer}

\subsection{测度的最大值与最小值}

类似于两个函数的最大值和最小值, 可以定义两个测度的最大值和最小值.

\begin{defn}[测度的最大值与最小值]
	\label{def:max-min-radon-measures}
	设 \(\mu,\nu\) 为 Radon 测度. 定义
	\[
	\max(\mu,\nu)=\frac{\mu+\nu+|\mu-\nu|}{2},
	\]
	以及
	\[
	\min(\mu,\nu)=\frac{\mu+\nu-|\mu-\nu|}{2}.
	\]
\end{defn}

\begin{thm}
	\label{thm:max-min-measures-order-property}
	\(\max(\mu,\nu)\) 是所有大于等于 \(\mu\) 和 \(\nu\) 的 Radon 测度中最小的一个; \(\min(\mu,\nu)\) 是所有小于等于 \(\mu\) 和 \(\nu\) 的 Radon 测度中最大的一个.
\end{thm}

\begin{proof}
	只证明最大值的断言, 最小值的断言类似. 首先
	\[
	\max(\mu,\nu)-\mu
	=
	\frac{|\mu-\nu|-(\mu-\nu)}{2}
	=(\nu-\mu)^+\geq0,
	\]
	所以 \(\max(\mu,\nu)\geq\mu\). 同理 \(\max(\mu,\nu)\geq\nu\).
	
	若 \(\sigma\geq\mu\) 且 \(\sigma\geq\nu\), 则 \(\sigma-\mu\geq0\), \(\sigma-\nu\geq0\). 又因为
	\[
	|\mu-\nu|\leq (\sigma-\mu)+(\sigma-\nu),
	\]
	故
	\[
	\max(\mu,\nu)
	=
	\frac{\mu+\nu+|\mu-\nu|}{2}
	\leq
	\frac{\mu+\nu+\sigma-\mu+\sigma-\nu}{2}
	=\sigma.
	\]
	因此 \(\max(\mu,\nu)\) 是最小上界.
\end{proof}

特别地,
\[
\mu^+=\max(\mu,0),
\qquad
\mu^-=\max(-\mu,0).
\]

\subsection{Hahn 分解与奇异性}

下面证明测度的 Hahn 分解. 它说明一个有符号测度可以把空间分成两部分, 在其中一部分上只看见正部, 在另一部分上只看见负部.

\begin{thm}[Hahn 分解]
	\label{thm:hahn-decomposition-radon-measure}
	设 \(\mu\) 为 Radon 测度. 则存在互为余集的集合 \(E^+\) 和 \(E^-\), 使得 \(E^+\) 关于 \(\mu^-\) 是零集, 而 \(E^-\) 关于 \(\mu^+\) 是零集. 并且可以选择这样的 \(E^+\) 和 \(E^-\), 使它们关于任意 Radon 正测度都是可测的.
\end{thm}

\begin{proof}
	取一个连续函数 \(F>0\), 使得 \(F\in L^1_{\mu^+}\). 例如可用一列逐渐扩大的紧集和快速下降的系数构造这样的函数. 由 \(\mu^+\) 的定义, 可以取 \(g_n\in\C_0(\R^d)\), \(0\leq g_n\leq F\), 使得
	\[
	\mu(g_n)>\mu^+(F)-2^{-n}.
	\]
	因为
	\[
	\mu(g_n)=\mu^+(g_n)-\mu^-(g_n),
	\]
	上式等价于
	\[
	\bigl(\mu^+(F)-\mu^+(g_n)\bigr)+\mu^-(g_n)<2^{-n}.
	\]
	于是特别有
	\[
	\mu^-(g_n)<2^{-n},
	\qquad
	\mu^+(F-g_n)<2^{-n}.
	\]
	令
	\[
	g(x)=\limsup_{n\to\infty}g_n(x).
	\]
	因为每个 \(g_n\) 连续, 所以 \(g\) 关于任意 Radon 正测度都是可测的, 且 \(0\leq g\leq F\).
	
	对任意 \(N\), 有
	\[
	g\leq \sum_{n=N}^{\infty}g_n.
	\]
	因此
	\[
	\mu^-(g)
	\leq
	\sum_{n=N}^{\infty}\mu^-(g_n)
	\leq
	2^{1-N}.
	\]
	令 \(N\to\infty\), 得 \(\mu^-(g)=0\). 又由 Fatou 引理,
	\[
	\mu^+(F-g)
	=
	\mu^+\bigl(\liminf_{n\to\infty}(F-g_n)\bigr)
	\leq
	\liminf_{n\to\infty}\mu^+(F-g_n)=0.
	\]
	
	定义
	\[
	E^+=\{x:g(x)>0\},
	\qquad
	E^-=\{x:g(x)=0\}.
	\]
	显然二者互为余集, 且由于 \(g\) 可测, 它们关于任意 Radon 正测度都是可测的. 由 \(\mu^-(g)=0\) 且 \(g>0\) 于 \(E^+\) 上可知, \(E^+\) 关于 \(\mu^-\) 是零集. 又因为 \(F-g=F>0\) 于 \(E^-\) 上, 且 \(\mu^+(F-g)=0\), 故 \(E^-\) 关于 \(\mu^+\) 是零集.
\end{proof}

\begin{thm}
	\label{thm:singularity-equivalent-characterizations}
	设 \(\mu,\nu\) 为两个 Radon 测度. 下列条件等价:
	\begin{enumerate}
		\item
		\[
		\min(|\mu|,|\nu|)=0;
		\]
		\item 存在互为余集的集合 \(E_\mu\) 和 \(E_\nu\), 使得 \(E_\mu\) 关于 \(|\nu|\) 是零集, 而 \(E_\nu\) 关于 \(|\mu|\) 是零集.
	\end{enumerate}
\end{thm}

\begin{proof}
	只需考虑 \(\mu\) 和 \(\nu\) 都为正测度的情形, 因为一般情形可对 \(|\mu|\) 和 \(|\nu|\) 应用同样论证.
	
	先设 \(\min(\mu,\nu)=0\). 令
	\[
	\rho=\mu-\nu.
	\]
	由 \(\min(\mu,\nu)=0\) 可知 \(\rho^+=\mu\), \(\rho^-=\nu\). 对 \(\rho\) 应用 Hahn 分解, 得到互余集合 \(E_\mu,E_\nu\), 使得 \(E_\mu\) 关于 \(\nu\) 是零集, \(E_\nu\) 关于 \(\mu\) 是零集.
	
	反过来, 假设存在这样的互余集合. 令
	\[
	\sigma=\min(\mu,\nu).
	\]
	则 \(\sigma\leq\mu\) 且 \(\sigma\leq\nu\). 因为 \(E_\nu\) 关于 \(\mu\) 是零集, 所以 \(E_\nu\) 关于 \(\sigma\) 也是零集; 因为 \(E_\mu\) 关于 \(\nu\) 是零集, 所以 \(E_\mu\) 关于 \(\sigma\) 也是零集. 二者互为余集, 因而整个空间关于 \(\sigma\) 是零集, 即 \(\sigma=0\). 因此 \(\min(\mu,\nu)=0\).
\end{proof}

\begin{defn}[互为奇异]
	\label{def:mutually-singular-measures}
	若定理 \ref{thm:singularity-equivalent-characterizations} 中的等价条件成立, 则称 \(\mu\) 与 \(\nu\) 互为奇异, 记作
	\[
	\mu\perp\nu.
	\]
\end{defn}

\begin{rem}
	古典术语中常说``\(\mu\) 关于 \(\nu\) 是奇异的''. 但由定义可见, 奇异性是对称关系, 因而说两个测度互为奇异更准确.
\end{rem}

\clearpage

\section{一维情形}
\label{sec:one-dimensional-general-radon-measures}

为了用 Stieltjes 积分描述一维的一般 Radon 测度, 我们需要引入有界变差函数的概念.

\begin{defn}[局部有界变差函数]
	\label{def:locally-bounded-variation-function}
	设 \(\Omega=(a,b)\subset\R\) 是一个开区间. 若函数 \(\varphi:\Omega\to\R\) 满足: 对任意紧子区间 \([\alpha,\beta]\subset\Omega\), 存在常数 \(V\), 使得对 \([\alpha,\beta]\) 的任意剖分
	\[
	\alpha=x_0<x_1<\cdots<x_N=\beta,
	\]
	都有
	\[
	\sum_{k=0}^{N-1}\left|\varphi(x_{k+1})-\varphi(x_k)\right|\leq V,
	\]
	则称 \(\varphi\) 在 \(\Omega\) 上是局部有界变差函数. 使上式成立的最小常数称为 \(\varphi\) 在 \([\alpha,\beta]\) 上的全变差.
\end{defn}

\begin{thm}\label{thm:bv-decomposition-difference-increasing}
	任一有界变差函数 \(\varphi\) 都可以写成两个增函数之差:
	\[
	\varphi=\varphi^+-\varphi^-.
	\]
	反过来, 两个增函数之差必为有界变差函数. 若 \(\varphi\) 右连续, 则可以选择 \(\varphi^+\) 和 \(\varphi^-\) 都右连续.
\end{thm}

\begin{proof}
	只需在一个紧区间 \([\alpha,\beta]\) 上证明. 为了简化记号, 先减去常数, 假设 \(\varphi(\alpha)=0\). 对 \(x\in[\alpha,\beta]\), 定义
	\[
	\varphi^+(x)=
	\sup\sum_{k=0}^{N-1}\bigl(\varphi(x_{k+1})-\varphi(x_k)\bigr)^+,
	\]
	\[
	\varphi^-(x)=
	\sup\sum_{k=0}^{N-1}\bigl(\varphi(x_{k+1})-\varphi(x_k)\bigr)^-,
	\]
	以及
	\[
	V_\varphi(x)=
	\sup\sum_{k=0}^{N-1}\left|\varphi(x_{k+1})-\varphi(x_k)\right|.
	\]
	这里上确界均取遍 \([\alpha,x]\) 的所有剖分
	\[
	\alpha=x_0<x_1<\cdots<x_N=x,
	\]
	并记
	\[
	r^+=\max(r,0),
	\qquad
	r^- =\max(-r,0).
	\]
	由有界变差性, 上述三个量都是有限的. 它们显然都是关于 \(x\) 的增函数.
	
	下面证明
	\[
	\varphi=\varphi^+-\varphi^-.
	\]
	对任一剖分, 有
	\[
	\varphi(x)-\varphi(\alpha)
	=
	\sum_{k=0}^{N-1}\bigl(\varphi(x_{k+1})-\varphi(x_k)\bigr)
	=
	\sum_{k=0}^{N-1}\bigl(\varphi(x_{k+1})-\varphi(x_k)\bigr)^+
	-
	\sum_{k=0}^{N-1}\bigl(\varphi(x_{k+1})-\varphi(x_k)\bigr)^-.
	\]
	由此得到
	\[
	\varphi(x)+\varphi^-(x)\leq \varphi^+(x).
	\]
	反过来, 由同一个等式也得到
	\[
	\varphi^+(x)\leq \varphi(x)+\varphi^-(x).
	\]
	故
	\[
	\varphi(x)=\varphi^+(x)-\varphi^-(x).
	\]
	
	类似地可得
	\[
	V_\varphi(x)=\varphi^+(x)+\varphi^-(x).
	\]
	这说明分解中的 \(\varphi^+\) 和 \(\varphi^-\) 是最小的: 若
	\[
	\varphi=u-v,
	\]
	其中 \(u,v\) 为增函数且 \(u(\alpha)=v(\alpha)=0\), 则必有
	\[
	u\geq \varphi^+,
	\qquad
	v\geq \varphi^-.
	\]
	
	反过来, 若 \(\varphi=u-v\), 其中 \(u,v\) 为增函数, 则对任意剖分有
	\[
	\sum_{k=0}^{N-1}|\varphi(x_{k+1})-\varphi(x_k)|
	\leq
	\sum_{k=0}^{N-1}\bigl(u(x_{k+1})-u(x_k)\bigr)
	+
	\sum_{k=0}^{N-1}\bigl(v(x_{k+1})-v(x_k)\bigr),
	\]
	右端等于
	\[
	u(\beta)-u(\alpha)+v(\beta)-v(\alpha).
	\]
	所以 \(\varphi\) 有界变差.
	
	最后, 若 \(\varphi\) 右连续, 则 \(\varphi^+\) 和 \(\varphi^-\) 可以取为右连续. 事实上, 最小分解中 \(\varphi^+\) 与 \(\varphi^-\) 不可能在同一点同时有正的右跳跃; 否则可从二者同时减去较小的那个跳跃, 从而得到更小的分解. 又因为
	\[
	\varphi=\varphi^+-\varphi^-
	\]
	右连续, 二者的右间断必须相同, 于是它们都右连续.
\end{proof}

设 \(\varphi\) 是 \(\Omega\) 上的局部有界变差函数. 对 \(f\in\C_0(\Omega)\), 取包含 \(\operatorname{supp} f\) 的紧区间 \([\alpha,\beta]\subset\Omega\). 若
\[
\alpha=x_0<x_1<\cdots<x_N=\beta
\]
是剖分, 并在每个 \([x_k,x_{k+1}]\) 中任取一点 \(\xi_k\), 定义 Riemann--Stieltjes 和
\[
\sum_{k=0}^{N-1} f(\xi_k)\bigl(\varphi(x_{k+1})-\varphi(x_k)\bigr).
\]
当剖分细度趋于零时, 这个和有极限, 记作
\[
\int_\Omega f\,\dd\varphi.
\]

这个极限存在的理由如下: 由定理 \ref{thm:bv-decomposition-difference-increasing}, 可写
\[
\varphi=\varphi^+-\varphi^-,
\]
其中 \(\varphi^+\) 和 \(\varphi^-\) 为增函数. 对增函数情形, 上一节已经定义了 Stieltjes 积分; 因此可令
\[
\int_\Omega f\,\dd\varphi
=
\int_\Omega f\,\dd\varphi^+
-
\int_\Omega f\,\dd\varphi^-.
\]
这也说明, 由有界变差函数定义的 Stieltjes 积分, 正是两个关于增函数的 Stieltjes 积分之差.

于是映射
\[
f\longmapsto \int_\Omega f\,\dd\varphi,
\qquad f\in\C_0(\Omega),
\]
是一个 Radon 测度. 反过来, 由上一节正测度的一维表示以及 Jordan 分解可知, 一维的一般 Radon 测度都可以用这种方式表示.

\begin{thm}\label{thm:one-dimensional-radon-measure-bv-stieltjes}
	设 \(\Omega=(a,b)\subset\R\). \(\Omega\) 上的一般 Radon 测度恰好是关于局部有界变差函数的 Stieltjes 积分. 更准确地说, 对任意 Radon 测度 \(\mu\), 存在局部有界变差函数 \(\varphi\), 使得
	\[
	\mu(f)=\int_\Omega f\,\dd\varphi,
	\qquad f\in\C_0(\Omega).
	\]
	若要求 \(\varphi\) 右连续并指定它在某一点的值, 则 \(\varphi\) 由 \(\mu\) 唯一确定.
\end{thm}

\begin{proof}
	若 \(\varphi\) 局部有界变差, 写成
	\[
	\varphi=\varphi^+-\varphi^-,
	\]
	则
	\[
	f\longmapsto \int f\,\dd\varphi
	=
	\int f\,\dd\varphi^+
	-
	\int f\,\dd\varphi^-
	\]
	是两个 Radon 正测度之差, 因而是 Radon 测度.
	
	反过来, 设 \(\mu\) 是 Radon 测度. 由 Jordan 分解,
	\[
	\mu=\mu^+-\mu^-.
	\]
	根据上一节的一维正测度表示, 存在右连续增函数 \(\varphi^+\) 与 \(\varphi^-\), 使得
	\[
	\mu^+(f)=\int f\,\dd\varphi^+,
	\qquad
	\mu^-(f)=\int f\,\dd\varphi^-.
	\]
	令
	\[
	\varphi=\varphi^+-\varphi^-.
	\]
	则 \(\varphi\) 局部有界变差, 且
	\[
	\mu(f)=\int f\,\dd\varphi.
	\]
	唯一性由上一节正测度表示的唯一性和 Jordan 分解的唯一性得到.
\end{proof}

\begin{exer}\label{exer:bv-function-positive-negative-measures}
	设局部有界变差函数 \(\varphi\) 定义 Radon 测度 \(\mu_\varphi\). 证明: 定理 \ref{thm:bv-decomposition-difference-increasing} 中构造的 \(\varphi^+\)、\(\varphi^-\) 以及全变差函数 \(V_\varphi\) 分别定义测度
	\[
	\mu_\varphi^+,
	\qquad
	\mu_\varphi^-,
	\qquad
	|\mu_\varphi|.
	\]
	换言之,
	\[
	(\mu_\varphi)^+=\mu_{\varphi^+},
	\qquad
	(\mu_\varphi)^-=\mu_{\varphi^-},
	\qquad
	|\mu_\varphi|=\mu_{V_\varphi}.
	\]
\end{exer}

\clearpage

\section{以 \texorpdfstring{\(\mu\)}{mu} 为基的测度}
\label{sec:measures-with-base-mu}

设 \(\mu\) 是一个 Radon 正测度. 和第二章中定义局部可积函数类似, 若函数 \(g\) 在每个紧集上关于 \(\mu\) 可积, 或等价地, 对任意 \(\varphi\in\C_0(\R^d)\) 都有
\[
\varphi g\in L^1_\mu,
\]
则称 \(g\) 为局部 \(\mu\)-可积函数, 记作 \(g\in L^1_{\operatorname{loc}}(\mu)\).

\begin{defn}[以 \(\mu\) 为基的测度]
	\label{def:measure-with-base-mu}
	设 \(\mu\) 是 Radon 正测度, \(g\in L^1_{\operatorname{loc}}(\mu)\). 定义
	\[
	(g\mu)(f)=\mu(fg)=\int f g\,\dd\mu,
	\qquad f\in\C_0(\R^d).
	\]
	则 \(g\mu\) 是一个 Radon 测度, 称为测度 \(\mu\) 与函数 \(g\) 的乘积. 任意形如 \(g\mu\) 的测度称为以 \(\mu\) 为基的测度.
\end{defn}

这个定义是合理的. 线性性显然. 若 \(f_n\in\C_0(\R^d)\) 都在同一个紧集 \(K\) 外为零, 且 \(f_n\to0\) 一致收敛, 则
\[
|(g\mu)(f_n)|\leq \|f_n\|_\infty\int_K |g|\,\dd\mu\to0,
\]
所以 \(g\mu\) 满足 Radon 测度的连续性条件.

\begin{rem}
	古典术语中, 形如 \(g\mu\) 的测度常称为关于 \(\mu\) 绝对连续的测度. 本书后面会用 Lebesgue--Radon--Nikodym 定理说明, 这与零集意义下的绝对连续性等价.
\end{rem}

\begin{lem}\label{lem:upper-integral-product-measure}
	设 \(g\geq0\) 且 \(g\in L^1_{\operatorname{loc}}(\mu)\). 则对任意非负函数 \(f\), 有
	\begin{equation}
		(g\mu)^*(f)=\mu^*(fg),
		\label{eq:upper-integral-product-measure}
	\end{equation}
	其中约定 \(0\cdot(+\infty)=0\).
\end{lem}

\begin{proof}
	若 \(f\in\C_0^+(\R^d)\), 则由定义直接有
	\[
	(g\mu)(f)=\mu(fg).
	\]
	由单调收敛性质, 等式先推广到 \(f\in\mathcal I^+\). 事实上, 若 \(0\leq f_n\uparrow f\), \(f_n\in\C_0^+\), 则
	\[
	(g\mu)^*(f)=\lim_n (g\mu)(f_n)
	=\lim_n \mu(f_n g)=\mu^*(fg).
	\]
	
	现在令 \(f\geq0\) 任意. 若 \(f\leq F\in\mathcal I^+\), 则由已经证明的情形,
	\[
	\mu^*(fg)\leq \mu^*(Fg)=(g\mu)^*(F).
	\]
	对所有这样的 \(F\) 取下确界, 得
	\[
	\mu^*(fg)\leq (g\mu)^*(f).
	\]
	反过来, 若 \(fg\leq H\in\mathcal I^+\), 在 \(g>0\) 的点上令 \(F=H/g\), 在 \(g=0\) 的点上令 \(F=+\infty\). 则 \(f\leq F\), 且 \(gF\leq H\). 以 \(\mathcal I^+\) 中函数从上逼近 \(F\) 后应用前一部分, 可得
	\[
	(g\mu)^*(f)\leq \mu^*(H).
	\]
	再对所有 \(H\in\mathcal I^+\), \(fg\leq H\) 取下确界, 得
	\[
	(g\mu)^*(f)\leq \mu^*(fg).
	\]
	两式合并即得 \eqref{eq:upper-integral-product-measure}.
\end{proof}

\begin{thm}\label{thm:l1-product-measure-characterization}
	设 \(g\geq0\) 且 \(g\in L^1_{\operatorname{loc}}(\mu)\). 则函数 \(f\in L^1_{g\mu}\) 的充分必要条件是
	\[
	fg\in L^1_\mu.
	\]
	在这种情形下,
	\begin{equation}
		(g\mu)(f)=\mu(fg).
		\label{eq:product-measure-integral-formula}
	\end{equation}
\end{thm}

\begin{proof}
	由引理 \ref{lem:upper-integral-product-measure}, 对任意非负函数 \(u\) 有
	\[
	(g\mu)^*(u)=\mu^*(ug).
	\]
	特别地, 对 \(u=|f|\), 得
	\[
	(g\mu)^*(|f|)=\mu^*(|fg|).
	\]
	因此 \(|f|\) 关于 \(g\mu\) 的上积分有限, 当且仅当 \(|fg|\) 关于 \(\mu\) 的上积分有限.
	
	同时, 零集和可测性的对应也由同一公式给出: 一个集合 \(E\) 关于 \(g\mu\) 为零集, 当且仅当 \(g=0\) 在 \(E\) 上 \(\mu\)-几乎处处成立. 因而 \(f\) 关于 \(g\mu\) 可测, 当且仅当 \(fg\) 关于 \(\mu\) 可测. 结合可积性判别定理, 即得
	\[
	f\in L^1_{g\mu}\quad\Longleftrightarrow\quad fg\in L^1_\mu.
	\]
	公式 \eqref{eq:product-measure-integral-formula} 先对 \(f\in\C_0\) 成立, 再由 \(L^1\) 逼近传递到一般可积函数.
\end{proof}

\begin{thm}\label{thm:l1-intersection-sum-measures}
	设 \(\mu\) 与 \(\nu\) 是 Radon 正测度. 则
	\[
	L^1_{\mu+\nu}=L^1_\mu\cap L^1_\nu.
	\]
	并且对任意 \(f\in L^1_{\mu+\nu}\), 有
	\begin{equation}
		(\mu+\nu)(f)=\mu(f)+\nu(f).
		\label{eq:sum-measures-integral-formula}
	\end{equation}
\end{thm}

\begin{proof}
	对任意非负函数 \(f\), 有
	\begin{equation}
		(\mu+\nu)^*(f)=\mu^*(f)+\nu^*(f).
		\label{eq:upper-integral-sum-measures}
	\end{equation}
	当 \(f\in\C_0^+\) 时, 这就是 \(\mu+\nu\) 的定义. 由单调收敛性质, 它推广到 \(f\in\mathcal I^+\), 再由上方控制函数取下确界, 推广到任意 \(f\geq0\).
	
	由 \eqref{eq:upper-integral-sum-measures} 可知, 一个集合关于 \(\mu+\nu\) 是零集, 当且仅当它同时关于 \(\mu\) 与 \(\nu\) 是零集. 因此一个函数关于 \(\mu+\nu\) 可测, 当且仅当它同时关于 \(\mu\) 与 \(\nu\) 可测.
	
	再由 \eqref{eq:upper-integral-sum-measures} 应用于 \(|f|\), 得
	\[
	(\mu+\nu)^*(|f|)<+\infty
	\]
	当且仅当
	\[
	\mu^*(|f|)<+\infty
	\quad\text{且}\quad
	\nu^*(|f|)<+\infty.
	\]
	结合可积性判别定理, 得
	\[
	L^1_{\mu+\nu}=L^1_\mu\cap L^1_\nu.
	\]
	最后, \eqref{eq:sum-measures-integral-formula} 对 \(f\in\C_0\) 由定义成立, 对一般 \(L^1_{\mu+\nu}\) 函数由逼近得到.
\end{proof}

下面研究第三节中定义的正部、负部和全变差在以 \(\mu\) 为基的测度上的具体形式.

\begin{thm}\label{thm:positive-negative-parts-density-measure}
	设 \(\mu\) 是 Radon 正测度, \(g\in L^1_{\operatorname{loc}}(\mu)\). 则
	\[
	(g\mu)^+=g^+\mu,
	\qquad
	(g\mu)^-=g^-\mu,
	\qquad
	|g\mu|=|g|\mu,
	\]
	其中
	\[
	g^+=\max(g,0),
	\qquad
	g^-=\max(-g,0).
	\]
\end{thm}

\begin{proof}
	令
	\[
	\nu_1=g^+\mu,
	\qquad
	\nu_2=g^-\mu.
	\]
	则 \(\nu_1,\nu_2\) 是 Radon 正测度, 且
	\[
	g\mu=\nu_1-\nu_2.
	\]
	设
	\[
	E^+=\{x:g(x)>0\},
	\qquad
	E^-=\{x:g(x)\leq0\}.
	\]
	则 \(E^+\) 关于 \(\nu_2\) 是零集, 而 \(E^-\) 关于 \(\nu_1\) 是零集. 因而 \(\nu_1\perp\nu_2\). 由 Jordan 分解的唯一性和最小性, 得
	\[
	(g\mu)^+=\nu_1=g^+\mu,
	\qquad
	(g\mu)^-=\nu_2=g^-\mu.
	\]
	于是
	\[
	|g\mu|=(g\mu)^++(g\mu)^-=(g^++g^-)\mu=|g|\mu.
	\]
\end{proof}

\begin{thm}\label{thm:zero-product-measure-iff-density-zero}
	测度 \(g\mu\) 为零测度的充分必要条件是 \(|g|\) 为 \(\mu\)-零函数, 即
	\[
	g=0\quad \mu\text{-几乎处处}.
	\]
\end{thm}

\begin{proof}
	若 \(|g|\) 是 \(\mu\)-零函数, 则对任意 \(f\in\C_0\), 有
	\[
	|(g\mu)(f)|\leq \|f\|_\infty\int |g|\,\dd\mu=0,
	\]
	所以 \(g\mu=0\).
	
	反过来, 若 \(g\mu=0\), 则由定理 \ref{thm:positive-negative-parts-density-measure}, 有
	\[
	|g|\mu=|g\mu|=0.
	\]
	这说明常数函数 \(1\) 关于测度 \(|g|\mu\) 的积分为零. 根据定理 \ref{thm:l1-product-measure-characterization}, 这等价于 \(|g|\) 关于 \(\mu\) 是零函数. 因此 \(g=0\) \(\mu\)-几乎处处.
\end{proof}

\begin{thm}\label{thm:monotone-density-dominated-measures}
	设 \(0\leq g_n\uparrow g\). 假定每个 \(g_n\) 都局部 \(\mu\)-可积, 并且存在 Radon 正测度 \(\nu\), 使得
	\[
	g_n\mu\leq \nu,
	\qquad n=1,2,\ldots.
	\]
	则 \(g\) 局部 \(\mu\)-可积.
\end{thm}

\begin{proof}
	任取 \(0\leq f\in\C_0(\R^d)\). 因为
	\[
	\int fg_n\,\dd\mu=(g_n\mu)(f)\leq \nu(f),
	\]
	且 \(fg_n\uparrow fg\), 由单调收敛定理得到
	\[
	\int fg\,\dd\mu\leq \nu(f)<+\infty.
	\]
	因此 \(fg\in L^1_\mu\) 对任意 \(f\in\C_0^+\) 成立. 由定义, \(g\) 局部 \(\mu\)-可积.
\end{proof}

\clearpage

\section{勒贝格分解, 勒贝格--Radon--Nikodym 定理}
\label{sec:lebesgue-radon-nikodym}

本节证明 Radon 测度之间最重要的结构分解. 固定一个 Radon 正测度 \(\mu\). 我们将说明, 任意 Radon 测度 \(\nu\) 都可以唯一地分解为两部分: 一部分以 \(\mu\) 为基, 另一部分与 \(\mu\) 互为奇异. 这就是 Lebesgue 分解. 进一步, 若 \(\nu\) 对 \(\mu\) 的零集不敏感, 则 \(\nu\) 必然可以写成 \(g\mu\), 这就是 Lebesgue--Radon--Nikodym 定理.

\begin{thm}[Lebesgue 分解]
	\label{thm:lebesgue-decomposition}
	设 \(\mu\) 为 Radon 正测度. 任意 Radon 测度 \(\nu\) 都可以唯一地写成
	\begin{equation}
		\nu=g\mu+\nu_s,
		\label{eq:lebesgue-decomposition}
	\end{equation}
	其中 \(g\) 是局部 \(\mu\)-可积函数, 而 \(\nu_s\perp\mu\). 若 \(\nu\) 是正测度, 则可以取
	\[
	g\geq0,\qquad \nu_s\geq0.
	\]
\end{thm}

证明这个定理之前, 先给出几个引理.

\begin{lem}\label{lem:max-mu-based-below-nu}
	设 \(\mu\) 和 \(\nu\) 为 Radon 正测度. 在所有以 \(\mu\) 为基并且不超过 \(\nu\) 的正测度中, 存在一个最大的测度.
\end{lem}

\begin{proof}
	取一个处处为正的连续函数 \(\eta\), 使得 \(\eta\in L^1_\nu\). 例如可以用一列紧集上的截断函数配以快速下降的系数来构造. 考虑所有满足
	\[
	0\leq h\mu\leq \nu
	\]
	的局部 \(\mu\)-可积函数 \(h\), 并令
	\[
	A=\sup (h\mu)(\eta).
	\]
	取一列 \(h_n\), 使得
	\[
	0\leq h_n\mu\leq\nu,
	\qquad
	(h_n\mu)(\eta)\to A.
	\]
	利用测度的最大值运算, 可以把 \(h_n\) 替换成
	\[
	\max(h_1,\ldots,h_n),
	\]
	从而不妨假定 \(h_n\uparrow\). 由上一节的单调密度定理可知, 极限
	\[
	h=\lim_{n\to\infty}h_n
	\]
	局部 \(\mu\)-可积, 并且
	\[
	h\mu\leq\nu.
	\]
	同时由单调收敛定理,
	\[
	(h\mu)(\eta)=\lim_{n\to\infty}(h_n\mu)(\eta)=A.
	\]
	
	现在证明 \(h\mu\) 是最大的. 若 \(k\mu\leq\nu\), 则
	\[
	\max(h,k)\mu\leq\nu.
	\]
	由 \(A\) 的定义,
	\[
	(\max(h,k)\mu)(\eta)\leq A=(h\mu)(\eta).
	\]
	因此
	\[
	\int (\max(h,k)-h)\eta\,\dd\mu=0.
	\]
	由于 \(\eta>0\), 得 \(\max(h,k)=h\) \(\mu\)-几乎处处, 即 \(k\leq h\) \(\mu\)-几乎处处. 所以
	\[
	k\mu\leq h\mu.
	\]
	这说明 \(h\mu\) 是所有此类测度中最大的一个.
\end{proof}

\begin{lem}\label{lem:dominated-positive-has-mu-based-submeasure}
	设 \(\rho\) 和 \(\mu\) 为 Radon 正测度, 且
	\[
	0\leq \rho\leq \mu,
	\qquad \rho\neq0.
	\]
	则存在一个非零的以 \(\mu\) 为基的 Radon 正测度 \(\sigma\), 使得
	\[
	0<\sigma\leq\rho.
	\]
\end{lem}

\begin{proof}
	若对每个 \(t>0\) 都有 \(\rho\leq t\mu\), 则令 \(t\downarrow0\) 可得 \(\rho=0\), 矛盾. 因此存在 \(t>0\), 使得
	\[
	\rho-t\mu\leq0
	\]
	不成立, 即 \((\rho-t\mu)^+\neq0\).
	
	对有符号测度 \(\rho-t\mu\) 作 Hahn 分解. 存在互余可测集 \(E^+\) 与 \(E^-\), 使得 \(E^+\) 关于 \((\rho-t\mu)^-\) 是零集, 而 \(E^-\) 关于 \((\rho-t\mu)^+\) 是零集. 记 \(\chi=\chi_{E^+}\). 则在测度意义下,
	\[
	\chi(\rho-t\mu)^-=0,
	\qquad
	(1-\chi)(\rho-t\mu)^+=0.
	\]
	由
	\[
	(\rho-t\mu)^+-(\rho-t\mu)^-=\rho-t\mu
	\]
	两边乘以 \(\chi\), 得
	\[
	\chi(\rho-t\mu)^+=\chi\rho-t\chi\mu.
	\]
	因此
	\[
	t\chi\mu\leq \chi\rho\leq \rho.
	\]
	另一方面, 因为 \(\rho\leq\mu\), 有
	\[
	\chi\mu\geq\chi\rho\geq\chi(\rho-t\mu)^+=(\rho-t\mu)^+\neq0.
	\]
	于是
	\[
	\sigma=t\chi\mu
	\]
	就是所需的非零以 \(\mu\) 为基的正测度.
\end{proof}

\begin{lem}\label{lem:non-singular-gives-mu-based-submeasure}
	设 \(\mu\) 和 \(\nu\) 为 Radon 正测度. 若
	\[
	\min(\mu,\nu)\neq0,
	\]
	即 \(\mu\) 与 \(\nu\) 不互为奇异, 则存在一个非零的以 \(\mu\) 为基的 Radon 正测度 \(\sigma\), 使得
	\[
	0<\sigma\leq\nu.
	\]
\end{lem}

\begin{proof}
	令
	\[
	\rho=\min(\mu,\nu).
	\]
	则 \(0<\rho\leq\mu\), 且 \(\rho\leq\nu\). 由引理 \ref{lem:dominated-positive-has-mu-based-submeasure}, 存在非零的以 \(\mu\) 为基的正测度 \(\sigma\), 使得
	\[
	0<\sigma\leq\rho\leq\nu.
	\]
\end{proof}

\begin{proof}[定理 \ref{thm:lebesgue-decomposition} 的证明]
	先设 \(\nu\geq0\). 由引理 \ref{lem:max-mu-based-below-nu}, 存在最大的以 \(\mu\) 为基且不超过 \(\nu\) 的正测度, 记为 \(g\mu\). 令
	\[
	\nu_s=\nu-g\mu.
	\]
	则 \(\nu_s\geq0\). 若 \(\nu_s\) 与 \(\mu\) 不互为奇异, 则由引理 \ref{lem:non-singular-gives-mu-based-submeasure}, 存在非零的以 \(\mu\) 为基的正测度 \(\sigma\), 使得
	\[
	0<\sigma\leq\nu_s.
	\]
	于是
	\[
	g\mu+\sigma\leq\nu,
	\]
	且 \(g\mu+\sigma\) 仍以 \(\mu\) 为基, 这与 \(g\mu\) 的最大性矛盾. 因此
	\[
	\nu_s\perp\mu.
	\]
	这就得到正测度情形的存在性.
	
	若 \(\nu\) 为一般 Radon 测度, 写出 Jordan 分解
	\[
	\nu=\nu^+-\nu^-.
	\]
	分别对 \(\nu^+\) 与 \(\nu^-\) 应用正测度情形, 得
	\[
	\nu^+=g_+\mu+\nu_s^+,
	\qquad
	\nu^-=g_-\mu+\nu_s^-.
	\]
	于是
	\[
	\nu=(g_+-g_-)\mu+(\nu_s^+-\nu_s^-).
	\]
	令
	\[
	g=g_+-g_-,
	\qquad
	\nu_s=\nu_s^+-\nu_s^-.
	\]
	因为 \(\nu_s^+\perp\mu\) 且 \(\nu_s^-\perp\mu\), 故 \(\nu_s\perp\mu\). 于是存在性得证.
	
	下面证明唯一性. 只需说明: 若
	\[
	g\mu+\nu_s=0,
	\qquad \nu_s\perp\mu,
	\]
	则 \(g\mu=0\) 且 \(\nu_s=0\). 由奇异性, 存在互余集合 \(E\) 与 \(F\), 使得 \(E\) 关于 \(|\nu_s|\) 是零集, 而 \(F\) 关于 \(\mu\) 是零集. 由
	\[
	g\mu=-\nu_s
	\]
	可得
	\[
	|g|\mu=|\nu_s|.
	\]
	在 \(E\) 上右端为零; 在 \(F\) 上左端为零. 因 \(E\cup F=\mathbb{R}^d\), 得
	\[
	|g|\mu=0.
	\]
	于是 \(g\mu=0\), 进而 \(\nu_s=0\). 唯一性得证.
\end{proof}

下面给出 Lebesgue--Radon--Nikodym 定理. 它刻画了一个正测度何时能写成 \(g\mu\) 的形式.

\begin{thm}[Lebesgue--Radon--Nikodym 定理]
	\label{thm:lebesgue-radon-nikodym}
	设 \(\mu\) 和 \(\nu\) 是 Radon 正测度. 下列条件等价:
	\begin{enumerate}
		\item \(\nu\) 以 \(\mu\) 为基, 即存在局部 \(\mu\)-可积函数 \(g\geq0\), 使得
		\[
		\nu=g\mu;
		\]
		\item 每个 \(\mu\)-零集都是 \(\nu\)-零集.
	\end{enumerate}
\end{thm}

\begin{proof}
	若 \(\nu=g\mu\), 则由上一节关于乘积测度的刻画可知: \(\mu\)-零集必为 \(g\mu\)-零集. 因而 \((1)\Rightarrow(2)\).
	
	反过来, 假设每个 \(\mu\)-零集都是 \(\nu\)-零集. 对 \(\nu\) 作关于 \(\mu\) 的 Lebesgue 分解:
	\[
	\nu=g\mu+\nu_s,
	\qquad
	\nu_s\perp\mu.
	\]
	由奇异性, 存在互余集合 \(E\) 与 \(F\), 使得 \(E\) 关于 \(\nu_s\) 是零集, 而 \(F\) 关于 \(\mu\) 是零集. 根据假设, \(F\) 也是 \(\nu\)-零集, 从而也是 \(\nu_s\)-零集. 于是 \(E\) 与 \(F\) 都是 \(\nu_s\)-零集, 其并为整个空间, 故
	\[
	\nu_s=0.
	\]
	因此 \(\nu=g\mu\), 即 \((2)\Rightarrow(1)\).
\end{proof}

\begin{cor}\label{cor:dominated-measure-radon-nikodym-density-bound}
	若 \(\mu\) 和 \(\nu\) 为 Radon 正测度, 且
	\[
	0\leq \nu\leq\mu,
	\]
	则存在局部 \(\mu\)-可积函数 \(g\), 使得
	\[
	\nu=g\mu,
	\qquad
	0\leq g\leq1\quad \mu\text{-几乎处处}.
	\]
\end{cor}

\begin{proof}
	由 \(\nu\leq\mu\) 可知, 每个 \(\mu\)-零集都是 \(\nu\)-零集. 由定理 \ref{thm:lebesgue-radon-nikodym}, 得 \(\nu=g\mu\). 又
	\[
	\mu-\nu=(1-g)\mu
	\]
	为正测度, 所以 \(1-g\geq0\) \(\mu\)-几乎处处, 即 \(g\leq1\) \(\mu\)-几乎处处. 由于 \(\nu\geq0\), 也有 \(g\geq0\) \(\mu\)-几乎处处.
\end{proof}

\begin{exer}\label{exer:common-base-for-two-measures}
	设 \(\mu\) 和 \(\nu\) 是两个 Radon 正测度. 证明存在一个 Radon 正测度 \(\lambda\), 使得 \(\mu\) 与 \(\nu\) 都以 \(\lambda\) 为基. 进一步, 把结论推广到可数多个给定的 Radon 正测度.
\end{exer}

\begin{exer}\label{exer:measure-geometric-mean}
	利用练习 \ref{exer:common-base-for-two-measures} 说明: 对任意两个 Radon 正测度 \(\mu\) 和 \(\nu\), 可以定义
	\[
	(\mu\nu)^{1/2},
	\qquad
	(\mu^2+\nu^2)^{1/2}
	\]
	等测度. 具体地说, 若 \(\mu=f\lambda\), \(\nu=g\lambda\), 可令
	\[
	(\mu\nu)^{1/2}=(fg)^{1/2}\lambda,
	\qquad
	(\mu^2+\nu^2)^{1/2}=(f^2+g^2)^{1/2}\lambda,
	\]
	并证明定义与共同基 \(\lambda\) 的选择无关.
	
	这些构造可用于定义可求长曲线的弧长测度. 例如曲线
	\[
	t\longmapsto x(t)=(x_1(t),\ldots,x_d(t))
	\]
	中每个坐标函数 \(x_j\) 为有界变差函数时, 可令
	\[
	\dd s=\left(\sum_{j=1}^d(\dd x_j)^2\right)^{1/2},
	\]
	从而把
	\[
	\int \varphi(t)\,\dd s(t)
	\]
	解释为 \(\varphi\) 在曲线上关于弧长的积分.
\end{exer}

\subsection{弥散部分与原子部分}

通常的 Lebesgue 分解还会把与基测度奇异的部分进一步分解为``连续的''部分和``跳跃的''部分. 在这里, 更合适的术语是弥散测度和原子测度.

\begin{defn}[弥散测度与原子测度]
	\label{def:diffuse-and-atomic-measure}
	设 \(\nu\) 为 Radon 测度.
	\begin{enumerate}
		\item 若每个点, 从而每个可数集合, 都是 \(|\nu|\)-零集, 则称 \(\nu\) 为弥散的.
		\item 若存在一个可数集合 \(A\), 使得 \(A^c\) 关于 \(|\nu|\) 是零集, 则称 \(\nu\) 为原子的.
	\end{enumerate}
\end{defn}

显然, 弥散测度与原子测度总是互为奇异的.

\begin{thm}\label{thm:diffuse-atomic-decomposition}
	任意 Radon 测度 \(\mu\) 都可以唯一地写成
	\begin{equation}
		\mu=\mu_1+\mu_2,
		\label{eq:diffuse-atomic-decomposition}
	\end{equation}
	其中 \(\mu_1\) 是弥散测度, \(\mu_2\) 是原子测度.
\end{thm}

\begin{proof}
	唯一性由弥散测度和原子测度互为奇异得到. 事实上, 若存在两个这样的分解, 将它们相减, 就得到一个既弥散又原子的测度, 它只能是零测度.
	
	下面证明存在性. 只需考虑 \(\mu\geq0\) 的情形; 一般情形对 \(|\mu|\) 或对正负部分分别处理即可. 使
	\[
	\mu(\{t\})\neq0
	\]
	的点 \(t\) 至多可数. 因为对任意紧集 \(K\) 和任意 \(\varepsilon>0\), 集合
	\[
	\{t\in K:\mu(\{t\})>\varepsilon\}
	\]
	只能是有限集, 否则会与 \(\mu(K)<+\infty\) 矛盾.
	
	把所有满足 \(\mu(\{t\})\neq0\) 的点排成序列 \(t_1,t_2,\ldots\). 定义
	\[
	\mu_2(f)=\sum_{n=1}^{\infty}\mu(\{t_n\})f(t_n),
	\qquad f\in C_0(\mathbb{R}^d).
	\]
	由于部分和不超过 \(\mu(f)\) 对 \(f\in C_0^+\) 成立, 上述级数收敛, 并且定义一个 Radon 正测度. 显然 \(\mu_2\) 集中在可数集 \(\{t_n\}\) 上, 因而是原子的.
	
	令
	\[
	\mu_1=\mu-\mu_2.
	\]
	则 \(\mu_1\geq0\). 对每个点 \(t\), 若 \(t=t_n\), 则 \(\mu_1(\{t\})=0\); 若 \(t\) 不在序列中, 则本来就有 \(\mu(\{t\})=0\), 也有 \(\mu_1(\{t\})=0\). 因而每个点关于 \(\mu_1\) 都是零集, \(\mu_1\) 弥散. 存在性得证.
\end{proof}

把 Lebesgue 分解和原子--弥散分解合在一起, 得到如下更细的分解.

\begin{thm}\label{thm:full-lebesgue-diffuse-atomic-decomposition}
	设 \(\mu\) 是弥散的 Radon 正测度. 则任意 Radon 测度 \(\nu\) 都可以唯一地写成
	\begin{equation}
		\nu=g\mu+\nu'+\nu'',
		\label{eq:full-lebesgue-diffuse-atomic-decomposition}
	\end{equation}
	其中 \(g\) 局部 \(\mu\)-可积, \(\nu'\) 是弥散测度且 \(\nu'\perp\mu\), 而 \(\nu''\) 是原子测度. 特别地, \(\nu''\perp\mu\).
\end{thm}

\begin{proof}
	先对 \(\nu\) 作关于 \(\mu\) 的 Lebesgue 分解:
	\[
	\nu=g\mu+\rho,
	\qquad
	\rho\perp\mu.
	\]
	再对 \(\rho\) 作弥散--原子分解:
	\[
	\rho=\nu'+\nu'',
	\]
	其中 \(\nu'\) 弥散, \(\nu''\) 原子. 因为 \(\rho\perp\mu\), 且 \(\nu',\nu''\) 都被 \(|\rho|\) 控制, 所以 \(\nu'\perp\mu\), \(\nu''\perp\mu\). 又因 \(\mu\) 弥散, 原子测度 \(\nu''\) 自动与 \(\mu\) 互为奇异.
	
	唯一性由三个部分所属类别两两分离得到: 以 \(\mu\) 为基的部分由 Radon--Nikodym 密度唯一确定; 与 \(\mu\) 奇异的部分由 Lebesgue 分解唯一确定; 在奇异部分中, 弥散部分和原子部分又由定理 \ref{thm:diffuse-atomic-decomposition} 唯一确定.
\end{proof}

\clearpage

\section{\texorpdfstring{\(L^p\)}{Lp} 上的连续线性泛函}
\label{sec:continuous-linear-functionals-on-lp}

本节固定一个 Radon 正测度 \(\mu\). 对 \(1\leq p<+\infty\), 记
\[
\|f\|_{p,\mu}
=
\left(\int |f|^p\,\dd\mu\right)^{1/p}.
\]
在不引起混淆时, 仍简记为 \(\|f\|_p\).

\begin{defn}[连续线性泛函]
	\label{def:continuous-linear-functional-lp}
	设 \(1\leq p<+\infty\). \(L^p_\mu\) 上的一个连续线性泛函, 是指一个线性映射
	\[
	\theta:L^p_\mu\longrightarrow \R,
	\]
	满足: 若 \(f_n\to f\) 于 \(L^p_\mu\), 即
	\[
	\|f_n-f\|_{p,\mu}\to0,
	\]
	则
	\[
	\theta(f_n)\to\theta(f).
	\]
\end{defn}

\begin{exer}\label{exer:continuous-linear-functional-bounded}
	证明: 定义 \ref{def:continuous-linear-functional-lp} 中的连续性等价于存在常数 \(C\geq0\), 使得
	\[
	|\theta(f)|\leq C\|f\|_{p,\mu},
	\qquad f\in L^p_\mu.
	\]
	满足上式的最小常数称为 \(\theta\) 的范数, 记作 \(\|\theta\|\).
\end{exer}

\begin{thm}[\(L^p\) 的对偶表示]
	\label{thm:lp-dual-representation-radon}
	设 \(1<p<+\infty\), 且 \(q\) 满足
	\[
	\frac1p+\frac1q=1.
	\]
	若 \(\theta\) 是 \(L^p_\mu\) 上的连续线性泛函, 则存在唯一的
	\(g\in L^q_\mu\), 使得
	\begin{equation}
		\theta(f)=\int f(x)g(x)\,\dd\mu(x),
		\qquad f\in L^p_\mu.
		\label{eq:lp-dual-representation-radon}
	\end{equation}
	并且
	\begin{equation}
		\|\theta\|=\|g\|_{q,\mu}.
		\label{eq:lp-dual-norm-equality-radon}
	\end{equation}
\end{thm}

\begin{proof}
	先说明 \eqref{eq:lp-dual-representation-radon} 右端确实给出一个连续线性泛函. 若 \(g\in L^q_\mu\), 则由 H\"older 不等式,
	\[
	\left|\int fg\,\dd\mu\right|
	\leq \|f\|_{p,\mu}\|g\|_{q,\mu}.
	\]
	因此 \(f\mapsto \int fg\,\dd\mu\) 是连续线性泛函, 且范数不超过
	\(\|g\|_{q,\mu}\). 由 H\"older 不等式的反向刻画, 这个范数实际上等于
	\(\|g\|_{q,\mu}\).
	
	下面证明任意连续线性泛函都具有这种形式. 设
	\[
	|\theta(f)|\leq C\|f\|_{p,\mu},
	\qquad f\in L^p_\mu.
	\]
	把 \(\theta\) 限制在 \(\C_0(\R^d)\) 上, 记为
	\[
	\nu(f)=\theta(f),
	\qquad f\in\C_0(\R^d).
	\]
	若 \(f_n\in\C_0(\R^d)\) 都在同一个紧集 \(K\) 外为零, 且 \(f_n\to0\) 一致收敛, 则
	\[
	\|f_n\|_{p,\mu}^p
	\leq \|f_n\|_\infty^p\mu^*(\chi_K)\to0.
	\]
	因此 \(\nu(f_n)\to0\). 这说明 \(\nu\) 是一个 Radon 测度.
	
	作 Jordan 分解
	\[
	\nu=\nu^+-\nu^-.
	\]
	对 \(0\leq f\in\C_0(\R^d)\), 由 \(\nu^+\) 的定义得
	\[
	\nu^+(f)
	=
	\sup_{0\leq h\leq f}\nu(h)
	\leq C\sup_{0\leq h\leq f}\|h\|_{p,\mu}
	\leq C\|f\|_{p,\mu}.
	\]
	同理
	\[
	\nu^-(f)\leq C\|f\|_{p,\mu},
	\qquad 0\leq f\in\C_0(\R^d).
	\]
	由单调收敛和上积分定义, 这些不等式可推广到任意非负函数:
	\begin{equation}
		(\nu^+)^*(f)
		\leq C\left(\mu^*(f^p)\right)^{1/p},
		\qquad f\geq0,
		\label{eq:positive-part-controlled-by-lp}
	\end{equation}
	并且 \(\nu^-\) 也满足同样估计.
	
	于是每个 \(\mu\)-零集都是 \(\nu^+\)-零集, 也是 \(\nu^-\)-零集. 由 Lebesgue--Radon--Nikodym 定理, 存在局部 \(\mu\)-可积函数 \(g_+,g_-\geq0\), 使得
	\[
	\nu^+=g_+\mu,
	\qquad
	\nu^-=g_-\mu.
	\]
	因此, 对 \(f\in\C_0(\R^d)\), 有
	\begin{equation}
		\theta(f)=\nu(f)=\int f(g_+-g_-)\,\dd\mu.
		\label{eq:theta-c0-density-before-lq}
	\end{equation}
	现在只需证明 \(g_+,g_-\in L^q_\mu\).
	
	我们只证明 \(g_+\in L^q_\mu\), 另一部分完全相同. 取一个处处正的连续函数 \(\eta\in L^q_\mu\). 例如可用紧集穷竭和快速下降系数构造. 令
	\[
	g_n=\min(g_+,n\eta),
	\qquad n=1,2,\ldots.
	\]
	则 \(g_n\in L^q_\mu\), 且 \(0\leq g_n\uparrow g_+\). 由于
	\[
	g_n\mu\leq g_+\mu=\nu^+,
	\]
	对任意 \(0\leq f\in L^p_\mu\), 由 \eqref{eq:positive-part-controlled-by-lp} 和 \(L^p\) 中 \(\C_0\) 的稠密性可得
	\[
	\int f g_n\,\dd\mu\leq C\|f\|_{p,\mu}.
	\]
	根据 H\"older 不等式的反向刻画,
	\[
	\|g_n\|_{q,\mu}\leq C,
	\qquad n=1,2,\ldots.
	\]
	令 \(n\to\infty\), 由单调收敛定理得到
	\[
	\|g_+\|_{q,\mu}\leq C.
	\]
	所以 \(g_+\in L^q_\mu\). 同理 \(g_-\in L^q_\mu\).
	
	令
	\[
	g=g_+-g_-.
	\]
	则 \(g\in L^q_\mu\). 由 \eqref{eq:theta-c0-density-before-lq} 知
	\eqref{eq:lp-dual-representation-radon} 对所有 \(f\in\C_0(\R^d)\) 成立. 又 \(\C_0(\R^d)\) 在 \(L^p_\mu\) 中稠密, 两边作为 \(L^p_\mu\) 上的连续线性泛函相等, 因而 \eqref{eq:lp-dual-representation-radon} 对所有 \(f\in L^p_\mu\) 成立.
	
	最后, 若 \(g_1,g_2\in L^q_\mu\) 都给出同一个泛函, 则
	\[
	\int f(g_1-g_2)\,\dd\mu=0,
	\qquad f\in L^p_\mu.
	\]
	取
	\[
	f=\operatorname{sgn}(g_1-g_2)|g_1-g_2|^{q-1},
	\]
	便得 \(g_1=g_2\) \(\mu\)-几乎处处. 唯一性得证. 范数等式
	\eqref{eq:lp-dual-norm-equality-radon} 已在证明开头说明.
\end{proof}

\begin{rem}
	定理 \ref{thm:lp-dual-representation-radon} 可以写成通常的形式
	\[
	(L^p_\mu)^*=L^q_\mu,
	\qquad 1<p<+\infty,
	\quad \frac1p+\frac1q=1.
	\]
	这里等号的含义是等距同构: 每个 \(g\in L^q_\mu\) 对应泛函
	\[
	f\longmapsto \int fg\,\dd\mu,
	\]
	并且这个对应保持范数.
\end{rem}

\clearpage

\section{古典情形的 Lebesgue 分解}
\label{sec:classical-lebesgue-decomposition}

在本节中, 我们把定理 \ref{thm:lebesgue-decomposition} 放回到最熟悉的情形中讨论: 基测度 \(\mu\) 取为 \(\R^d\) 上的 Lebesgue 测度. 设 \(\nu\) 是 \(\R^d\) 上的 Radon 测度, 并考虑它关于 Lebesgue 测度的 Lebesgue 分解
\begin{equation}
	\nu=g\mu+\nu_s,
	\label{eq:classical-lebesgue-decomposition}
\end{equation}
其中 \(g\in L^1_{\operatorname{loc}}(\R^d)\), 且 \(\nu_s\perp\mu\). 我们要说明, 这个分解与第二章中 Lebesgue 微分定理之间有直接联系: 密度 \(g\) 正是测度 \(\nu\) 对 Lebesgue 体积的微商.

若 \(S\) 是球, 记
\[
\nu(S)=\int_S \dd\nu.
\]
当 \(\nu\) 为有符号测度时, 这按其 Jordan 分解理解.

\begin{thm}\label{thm:classical-lebesgue-decomposition-derivative}
	设 \(\mu\) 为 \(\R^d\) 上的 Lebesgue 测度, \(\nu\) 为 Radon 测度, 并作分解 \eqref{eq:classical-lebesgue-decomposition}. 则对几乎所有的 \(x\in\R^d\), 有
	\begin{equation}
		g(x)=
		\lim_{\substack{\mu(S)\to0\\ x\in S}}
		\frac{\nu(S)}{\mu(S)},
		\label{eq:rn-density-as-derivative}
	\end{equation}
	其中 \(S\) 取遍包含 \(x\) 的球.
\end{thm}

\begin{proof}
	先看 \(g\mu\) 这一部分. 由 Lebesgue 微分定理, 对几乎所有 \(x\), 有
	\[
	\lim_{\substack{\mu(S)\to0\\x\in S}}
	\frac{1}{\mu(S)}\int_S g(y)\,\dd y=g(x).
	\]
	因此只需证明: 若 \(\nu\perp\mu\), 则
	\begin{equation}
		\lim_{\substack{\mu(S)\to0\\x\in S}}
		\frac{\nu(S)}{\mu(S)}=0
		\label{eq:singular-measure-derivative-zero}
	\end{equation}
	对几乎所有 \(x\) 成立. 对一般有符号测度, 对全变差测度 \(|\nu|\) 证明即可, 因为
	\[
	|\nu(S)|\leq |\nu|(S).
	\]
	所以以下设 \(\nu\) 为 Radon 正测度, 且 \(\nu\perp\mu\).
	
	首先回忆 Hardy 最大定理的测度版本. 对 Radon 正测度 \(\rho\), 定义
	\[
	\rho^\natural(x)=
	\sup_{S\ni x}\frac{\rho(S)}{\mu(S)},
	\]
	其中 \(S\) 取遍包含 \(x\) 的球. 若 \(\rho(\R^d)<+\infty\), 则对任意 \(s>0\), 有
	\begin{equation}
		\mu^*\{x:\rho^\natural(x)>s\}
		\leq
		\frac{4^d}{s}\rho(\R^d).
		\label{eq:measure-maximal-weak-type}
	\end{equation}
	这个证明与 Hardy 最大定理完全相同, 只是把 \(\int_S |f|\,\dd x\) 换成 \(\rho(S)\).
	
	由 \(\nu\perp\mu\), 存在互余集合 \(E_1,E_2\), 使得
	\[
	\mu(E_1)=0,
	\qquad
	\nu(E_2)=0.
	\]
	给定 \(s>0\). 令
	\[
	M_s=
	\left\{x:
	\limsup_{\substack{\mu(S)\to0\\x\in S}}
	\frac{\nu(S)}{\mu(S)}>s
	\right\}.
	\]
	取开集 \(O\supset E_2\), 且 \(\nu(O)<+\infty\). 令
	\[
	\nu_O=\chi_O\nu.
	\]
	因为 \(\nu(E_2)=0\), 可以进一步把 \(O\) 取得使 \(\nu(O)=\nu_O(\R^d)\) 任意小.
	
	若 \(x\in M_s\cap E_2\), 由于 \(O\) 是包含 \(E_2\) 的开集, 当包含 \(x\) 的球 \(S\) 足够小时有 \(S\subset O\). 因此
	\[
	\nu(S)=\nu_O(S)
	\]
	对足够小的 \(S\) 成立, 从而 \(\nu_O^\natural(x)>s\). 于是
	\[
	M_s\cap E_2\subset \{x:\nu_O^\natural(x)>s\}.
	\]
	又因 \(E_1\) 是 \(\mu\)-零集, 有
	\[
	\mu^*(M_s)=\mu^*(M_s\cap E_2).
	\]
	结合 \eqref{eq:measure-maximal-weak-type}, 得
	\[
	\mu^*(M_s)
	\leq
	\mu^*\{x:\nu_O^\natural(x)>s\}
	\leq
	\frac{4^d}{s}\nu(O).
	\]
	由于 \(O\supset E_2\) 可以取得使 \(\nu(O)\) 任意小, 得
	\[
	\mu^*(M_s)=0.
	\]
	令 \(s\) 取遍正有理数, 即得 \eqref{eq:singular-measure-derivative-zero}. 定理得证.
\end{proof}

\begin{rem}
	定理 \ref{thm:classical-lebesgue-decomposition-derivative} 的意思是: 在古典 Lebesgue 测度背景下, Radon--Nikodym 密度不是一个抽象存在的函数, 而是测度 \(\nu\) 对体积测度 \(\mu\) 的局部变化率. 奇异部分的局部体积密度几乎处处为零.
\end{rem}

在一维情形, Radon 测度可以用 Stieltjes 积分表示. 因此定理 \ref{thm:classical-lebesgue-decomposition-derivative} 给出有界变差函数几乎处处可微的另一种解释.

\begin{thm}\label{thm:bv-function-ae-derivative-from-decomposition}
	若 \(\varphi\) 是一维区间上的有界变差函数, 则它的微商 \(\varphi'(x)\) 几乎处处存在, 并且 \(\varphi'\) 是关于 Lebesgue 测度的局部可积函数.
\end{thm}

\begin{proof}
	设 \(\dd\varphi\) 是由 \(\varphi\) 定义的 Stieltjes 测度. 对 \(\dd\varphi\) 作关于 Lebesgue 测度的 Lebesgue 分解:
	\[
	\dd\varphi=g(x)\,\dd x+\dd\varphi_s.
	\]
	由定理 \ref{thm:classical-lebesgue-decomposition-derivative}, 对几乎所有的 \(x\), 有
	\[
	g(x)=
	\lim_{|I|\to0,\ x\in I}
	\frac{\dd\varphi(I)}{|I|}.
	\]
	而当 \(I=[a,b]\) 时,
	\[
	\dd\varphi(I)=\varphi(b)-\varphi(a)
	\]
	按端点处的左右连续约定理解. 因此上述极限正是 \(\varphi\) 的微商. 又 \(g\in L^1_{\operatorname{loc}}\), 所以 \(\varphi'\in L^1_{\operatorname{loc}}\).
\end{proof}

\begin{exer}\label{exer:classical-decomposition-total-variation-derivative}
	在定理 \ref{thm:classical-lebesgue-decomposition-derivative} 的假设下, 证明对几乎所有的 \(x\) 和任意实数 \(a\), 有
	\[
	\lim_{\substack{\mu(S)\to0\\x\in S}}
	\frac{|\nu-a\mu|(S)}{\mu(S)}
	=
	|g(x)-a|.
	\]
	提示: 先考察 \(a=g(x)\) 的情形.
\end{exer}

\begin{exer}\label{exer:approximate-identity-measure-density}
	设 \(h\in L^\infty(\R^d)\cap L^1(\R^d)\), 且 \(h\) 在无穷远处趋于零. 令
	\[
	h_\varepsilon(x)=\varepsilon^{-d}h(x/\varepsilon),
	\qquad \varepsilon>0.
	\]
	证明
	\[
	\lim_{\varepsilon\to0}
	\int h_\varepsilon(x-y)\,\dd\nu(y)
	=
	g(x)\int h(y)\,\dd y
	\]
	对几乎所有的 \(x\) 成立. 其中 \(\nu\) 和 \(g\) 的含义与定理 \ref{thm:classical-lebesgue-decomposition-derivative} 相同.
\end{exer}

\begin{exer}\label{exer:direct-proof-density-of-measure}
	仿照第二章第十节结尾处理单调函数微商的方法, 直接证明
	\[
	g(x)=
	\lim_{\substack{\mu(S)\to0\\x\in S}}
	\frac{1}{\mu(S)}\int_S\dd\nu
	\]
	几乎处处存在. 在 \(\mu\) 为 Lebesgue 测度的情形下, 这给出定理 \ref{thm:classical-lebesgue-decomposition-derivative} 的另一种证明.
\end{exer}

\clearpage

\section{各种各样的推广}
\label{sec:miscellaneous-generalizations}

到目前为止, 我们主要讨论的是实值函数的积分. 事实上, 前面建立的许多构造并不依赖于函数值必须是实数. 只要把实数值换成 Banach 空间中的向量值, 积分理论仍然可以相当自然地展开.

设 \(B\) 是一个实 Banach 空间. 记
\[
\C_0(\R^d,B)
\]
为所有从 \(\R^d\) 到 \(B\) 的连续函数全体, 并要求它们具有紧致支集. 若 \(f\in\C_0(\R^d,B)\), 则可以用 Riemann 和定义
\[
\int_{\R^d}f(x)\,\dd x\in B.
\]
具体地说, 取包含 \(\operatorname{supp}f\) 的有限区间 \(I\), 对 \(I\) 作剖分, 并在每个小区间 \(I_k\) 中取点 \(\xi_k\), 考察 Riemann 和
\[
\sum_k f(\xi_k)m(I_k).
\]
由于 \(f\) 在紧集上一致连续, 当剖分细度趋于零时, 这些 Riemann 和在 \(B\) 中收敛. 因为 \(B\) 是 Banach 空间, 极限存在并与剖分及取点方式无关. 这个极限就定义为 \(\int f\,\dd x\).

这个向量值积分具有与实值情形相同的线性性质:
\[
\int_{\R^d}(af+bg)(x)\,\dd x
=
a\int_{\R^d}f(x)\,\dd x
+
b\int_{\R^d}g(x)\,\dd x,
\]
其中 \(a,b\in\R\), \(f,g\in\C_0(\R^d,B)\). 此外还有估计
\[
\left\|\int_{\R^d}f(x)\,\dd x\right\|
\leq
\int_{\R^d}\|f(x)\|_B\,\dd x.
\]

由此可以定义 \(B\) 值的 \(L^1\) 函数. 若 \(f\) 是几乎处处定义的 \(B\) 值函数, 并且对任意 \(\varepsilon>0\), 存在 \(g\in\C_0(\R^d,B)\), 使得
\[
\int_{\R^d}^{*}\|f(x)-g(x)\|_B\,\dd x<\varepsilon,
\]
则称 \(f\) 是 \(B\) 值可积函数, 记作
\[
f\in L^1(\R^d,B).
\]
若 \(g_n\in\C_0(\R^d,B)\) 且
\[
\int_{\R^d}^{*}\|f(x)-g_n(x)\|_B\,\dd x\to0,
\]
则 \(\int g_n\,\dd x\) 在 \(B\) 中是 Cauchy 列, 因而收敛. 定义
\[
\int_{\R^d}f(x)\,\dd x
=
\lim_{n\to\infty}\int_{\R^d}g_n(x)\,\dd x.
\]
这个定义与逼近列的选择无关, 因为
\[
\left\|\int g_n\,\dd x-\int h_n\,\dd x\right\|
\leq
\int^{*}\|g_n-h_n\|_B\,\dd x.
\]

于是, \(B\) 值积分仍然满足线性性、三角不等式、控制收敛定理等形式. 例如, 若 \(f_n\to f\) 几乎处处, 且存在 \(G\in L^1(\R^d)\) 使得
\[
\|f_n(x)\|_B\leq G(x)
\]
几乎处处成立, 则
\[
\int_{\R^d}\|f_n(x)-f(x)\|_B\,\dd x\to0,
\]
从而
\[
\int_{\R^d}f_n(x)\,\dd x
\longrightarrow
\int_{\R^d}f(x)\,\dd x
\]
在 \(B\) 的范数意义下成立.

\begin{rem}
	向量值积分与实值积分的一个重要差别是: 一般 Banach 空间没有天然的序结构, 因而不能直接谈论 ``\(f\geq0\)'' 或 ``单调收敛''. 所以在向量值情形中, 很多定理要通过范数估计来表达. 例如, 正函数的单调收敛定理通常被控制收敛定理和 \(L^1\) 完备性所替代.
\end{rem}

上面的讨论还可以进一步推广. 实际上, 除了 \(\R^d\) 是一个具有可数基的局部紧致拓扑空间以外, 我们在建立积分理论时并没有使用太多欧氏空间的特殊性质. 这里所谓局部紧致, 是指每个点都有一个紧致邻域; 所谓具有可数基, 是指存在一个可数的开集族, 使得任意开集都可以写成其中某些开集的并.

因此, 若 \(X\) 是一个具有可数基的局部紧致 Hausdorff 空间, 可以令 \(\C_0(X)\) 表示 \(X\) 上具有紧致支集的连续实值函数全体. 只要给定一个正线性泛函
\[
\mu:\C_0(X)\longrightarrow\R,
\]
即
\[
\mu(af+bg)=a\mu(f)+b\mu(g),
\qquad
\mu(f)\geq0\quad(f\geq0),
\]
便可以仿照前面的方法定义上积分、零集、可测函数、可积函数以及关于 \(\mu\) 的积分. 这样得到的正线性泛函同样称为 \(X\) 上的 Radon 正测度.

更抽象地说, 甚至可以暂时把拓扑结构也放在一边. 设 \(E\) 是一个集合, \(\C_0\) 是定义在 \(E\) 上的一个实值函数线性空间, 并假定
\[
f\in\C_0
\quad\Longrightarrow\quad
|f|\in\C_0.
\]
于是对 \(f,g\in\C_0\), 也有
\[
\max(f,g)=\frac{f+g+|f-g|}{2}\in\C_0,
\]
以及
\[
\min(f,g)=\frac{f+g-|f-g|}{2}\in\C_0.
\]
也就是说, \(\C_0\) 是一个函数格.

若 \(\mu\) 是 \(\C_0\) 上的正线性泛函, 并且满足适当的从上连续性条件, 即当
\[
f_n\in\C_0^+,
\qquad
f_n\downarrow0
\]
时有
\[
\mu(f_n)\to0,
\]
则可以令 \(\mathcal I^+\) 为 \(\C_0^+\) 中上升序列极限所组成的函数类. 在这种抽象框架下, 仍然可以定义非负函数的上积分、零函数、零集合、可测函数和可积函数. 前面许多关于积分的基本性质仍然成立.

不过, 在这种完全抽象的形式中, 理论是否有内容, 取决于 \(\mathcal I^+\) 中函数类是否足够丰富, 以及它是否能有效刻画所研究的问题. 因此, 本书不再继续展开这个一般框架.

\begin{rem}
	概率论教材中常常从类似的抽象观点开始. 那里通常先给定一个集合 \(E\) 和一个由某些集合的特征函数生成的函数类, 例如由某个集合代数中集合的特征函数的有限线性组合构成的函数类, 然后在这个函数类上定义正线性泛函, 再逐步扩张为积分. 这与本书从 \(\C_0\) 和正线性泛函出发的路线在思想上是一致的.
\end{rem}

\begin{rem}
	本节的意义在于说明: 积分理论的核心并不局限于 \(\R^d\) 上的体积测度. 一旦抓住 ``正线性泛函'' 和 ``从简单函数类扩张到更大函数类'' 这两个结构, 就可以得到许多不同背景下的积分理论. 前面关于 Lebesgue 积分、Radon 测度、Stieltjes 积分和 Radon--Nikodym 定理的讨论, 都是这个共同思想的具体实现.
\end{rem}
	
	\appendix

\chapter{附录}
\label{app:basic-notation-topology}

众所周知, 空间 \(\mathbb R^d\) 是实数域 \(\mathbb R\) 上的一个向量空间. 本书前面各章中的积分理论, 都是在这个空间及其开子集上建立的. \(\mathbb R^d\) 上的每一个线性映射都可写成
\[
x\longmapsto Ax,
\]
其中 \(A\) 是一个 \(d\times d\) 实矩阵, 而 \(x\) 可以看成一个 \(d\times 1\) 列矩阵. 所有可逆 \(d\times d\) 实矩阵组成的群记为
\[
\operatorname{GL}(d,\mathbb R).
\]

若 \(x=(x_1,\ldots,x_d)\in\mathbb R^d\), 我们总用
\[
\|x\|=(x_1^2+\cdots+x_d^2)^{1/2}
\]
表示 \(x\) 的 Euclidean 范数.

\section{拓扑记号}
\label{app:topological-notation}

回忆 \(\mathbb R^d\) 中的开集定义. 若 \(O\subset\mathbb R^d\), 且对每个 \(x\in O\), 都存在 \(\delta>0\), 使得
\[
\{y:\|y-x\|<\delta\}\subset O,
\]
则称 \(O\) 为开集. 开集的任意并集以及有限交集仍是开集.

因此, 包含在给定集合 \(M\subset\mathbb R^d\) 中的所有开集的并仍是开集. 这个开集称为 \(M\) 的内部, 记作
\[
M^\circ.
\]
于是 \(x\in M^\circ\) 的充分必要条件是: 存在以 \(x\) 为中心的开球包含在 \(M\) 中.

若集合 \(F\subset\mathbb R^d\) 的余集
\[
F^c=\mathbb R^d\setminus F
\]
是开集, 则称 \(F\) 为闭集. 闭集的任意交集以及有限并集仍是闭集. 包含给定集合 \(M\) 的所有闭集的交称为 \(M\) 的闭包, 记作
\[
\overline M.
\]
也就是说, \(\overline M\) 是包含 \(M\) 的最小闭集. 点 \(x\in\overline M\) 的充分必要条件是: 每个以 \(x\) 为中心的开球都与 \(M\) 相交.

集合
\[
\partial M=\overline M\cap \overline{M^c}
\]
称为 \(M\) 的边界. 因此, \(x\in\partial M\) 当且仅当任意以 \(x\) 为中心的开球既与 \(M\) 相交, 又与 \(M^c\) 相交.

\begin{defn}[紧集]
	\label{def:compact-set-appendix}
	\(\mathbb R^d\) 中的一个闭且有界的集合称为紧集.
\end{defn}

根据 Heine--Borel 定理, 紧集还可以用开覆盖来刻画: 若 \(K\) 为紧集, 则覆盖 \(K\) 的任意一族开集必有有限子族仍覆盖 \(K\).

\begin{exer}
	\label{exer:closed-sequential-characterization}
	证明:
	\begin{enumerate}
		\item 集合 \(F\subset\mathbb R^d\) 为闭集的充分必要条件是: \(F\) 中任意收敛序列的极限仍属于 \(F\).
		\item 集合 \(K\subset\mathbb R^d\) 为紧集的充分必要条件是: \(K\) 中任意序列都有一个收敛子序列, 且该子序列的极限仍属于 \(K\).
	\end{enumerate}
\end{exer}

\section{函数和集合的基本记号}
\label{app:function-set-notation}

对于函数
\[
f:\mathbb R^d\longrightarrow\mathbb R,
\]
其支集 \(\operatorname{supp}f\) 定义为集合
\[
\{x:f(x)\neq0\}
\]
的闭包. 等价地, \((\operatorname{supp}f)^c\) 是使 \(f(x)=0\) 的点集的内部. 因为支集总是闭集, 所以 \(\operatorname{supp}f\) 为紧集的充分必要条件是: \(f\) 在某个有界集之外恒等于零.

我们把从 \(\mathbb R^d\) 到 \(\mathbb R\) 的所有连续函数组成的集合记为
\[
C=C(\mathbb R^d)=C(\mathbb R^d,\mathbb R),
\]
把其中具有紧致支集的函数组成的集合记为
\[
C_0=C_0(\mathbb R^d).
\]

对于任意函数列 \(f_n\), 可逐点定义
\[
|f|,\qquad f_1+f_2,\qquad f_1f_2,\qquad
\max(f_1,f_2),\qquad \min(f_1,f_2),
\]
以及
\[
\lim f_n,\qquad \sup_n f_n,\qquad \inf_n f_n.
\]
例如
\[
\bigl(\max(f_1,f_2)\bigr)(x)=\max\{f_1(x),f_2(x)\},
\]
而
\[
\bigl(\sup_n f_n\bigr)(x)=\sup_n f_n(x),
\qquad x\in\mathbb R^d.
\]

我们记
\[
f^+=\max(f,0)\geq0,
\qquad
f^-=\max(-f,0)\geq0.
\]
于是
\[
f=f^+-f^-,
\qquad
|f|=f^++f^-.
\]

若 \(E\subset\mathbb R^d\), 其特征函数记为
\[
\chi_E(x)=
\begin{cases}
	1, & x\in E,\\
	0, & x\in E^c=\mathbb R^d\setminus E.
\end{cases}
\]
于是
\[
\chi_{E^c}=1-\chi_E,
\]
并且对任意集合 \(E_1,E_2\subset\mathbb R^d\), 有
\[
\chi_{E_1\cup E_2}
=
\max(\chi_{E_1},\chi_{E_2})
=
\chi_{E_1}+\chi_{E_2}-\chi_{E_1}\chi_{E_2},
\]
\[
\chi_{E_1\cap E_2}
=
\min(\chi_{E_1},\chi_{E_2})
=
\chi_{E_1}\chi_{E_2},
\]
以及
\[
\chi_{E_1\triangle E_2}
=
|\chi_{E_1}-\chi_{E_2}|
=
\chi_{E_1}+\chi_{E_2}-2\chi_{E_1}\chi_{E_2}.
\]
这里
\[
E_1\triangle E_2=(E_1\cup E_2)\setminus(E_1\cap E_2)
\]
称为 \(E_1\) 和 \(E_2\) 的对称差, 即由属于 \(E_1\) 与 \(E_2\) 中恰好一个集合的点组成.

\begin{exer}
	\label{exer:indicator-continuity-boundary}
	证明: \(\chi_E\) 在点 \(x\) 连续的充分必要条件是
	\[
	x\notin\partial E.
	\]
\end{exer}

\section{极限、上下极限和上确界}
\label{app:limsup-sup-notation}

对于序列 \((a_n)\), 其中每个 \(a_n\) 可以是实数, 也可以是 \(+\infty\) 或 \(-\infty\), 若
\[
a_n\leq a_{n+1}
\qquad
(\text{或 }a_n\geq a_{n+1})
\]
且
\[
\lim_{n\to\infty}a_n=a,
\]
则分别用
\[
a_n\uparrow a,
\qquad
a_n\downarrow a
\]
表示.

对于任意序列 \((a_n)\), 优于它的最小下降序列是
\[
\left(\sup_{k\geq n}a_k\right)_{n=1}^{\infty}.
\]
因此定义
\[
\overline{\lim}_{n\to\infty}a_n
=
\lim_{n\to\infty}\sup_{k\geq n}a_k
=
\inf_n\sup_{k\geq n}a_k.
\]
类似地, 定义
\[
\underline{\lim}_{n\to\infty}a_n
=
\lim_{n\to\infty}\inf_{k\geq n}a_k
=
\sup_n\inf_{k\geq n}a_k.
\]
显然
\[
\underline{\lim}_{n\to\infty}a_n
\leq
\overline{\lim}_{n\to\infty}a_n,
\]
且等号成立当且仅当 \(\lim_{n\to\infty}a_n\) 存在.

设
\[
f:M\longrightarrow\mathbb R
\]
是任意函数. 我们分别用
\[
\sup_M f,
\qquad
\inf_M f
\]
表示值集 \(f(M)\) 的上确界和下确界. 若 \(f\) 上方无界, 则约定
\[
\sup_M f=+\infty.
\]
显然
\[
\inf_M f=-\sup_M(-f),
\]
并且
\[
\sup_M(f+g)
\leq
\sup_M f+
\sup_M g.
\]
取 \(M=\{k:k\geq n\}\), 便可得到: 若 \((a_n)\) 和 \((b_n)\) 为上方有界序列, 则
\[
\overline{\lim}_{n\to\infty}(a_n+b_n)
\leq
\overline{\lim}_{n\to\infty}a_n
+
\overline{\lim}_{n\to\infty}b_n.
\]
类似地, 若二者下方有界, 则
\[
\underline{\lim}_{n\to\infty}(a_n+b_n)
\geq
\underline{\lim}_{n\to\infty}a_n
+
\underline{\lim}_{n\to\infty}b_n.
\]

若 \(f\) 和 \(g\) 上方有界, 则由不等式
\[
f\leq g+|f-g|,
\qquad
g\leq f+|f-g|
\]
可得
\[
\left|\sup_M f-\sup_M g\right|
\leq
\sup_M|f-g|.
\]
当 \(M=\{1,2\}\) 时, 这给出
\[
\left|\max(a_1,a_2)-\max(b_1,b_2)\right|
\leq
\max\{|a_1-b_1|,|a_2-b_2|\}.
\]

\section{距离函数和截断函数}
\label{app:distance-function-cutoff}

对于 \(\mathbb R^d\) 的非空子集 \(M\), 定义点 \(x\) 到 \(M\) 的距离为
\[
d(x,M)=\inf_{y\in M}\|x-y\|
=
-\sup_{y\in M}\bigl(-\|x-y\|\bigr).
\]
注意
\[
d(x,M)=0
\]
的充分必要条件是
\[
x\in\overline M.
\]
其次, 由
\begin{align*}
	\left|d(x_1,M)-d(x_2,M)\right|
	&\leq
	\sup_{y\in M}\left|-\|x_1-y\|+\|x_2-y\|\right|\\
	&\leq \|x_1-x_2\|
\end{align*}
可知, 函数
\[
x\longmapsto d(x,M)
\]
是连续的, 事实上是一致连续的.

\begin{exer}
	\label{exer:distance-cutoff-open-closed}
	设 \(F\) 是开集 \(O\) 的一个闭子集. 证明函数
	\[
	f(x)=
	\frac{d(x,O^c)}{d(x,F)+d(x,O^c)},
	\qquad x\in\mathbb R^d,
	\]
	是连续的, 并且
	\[
	0\leq f\leq1,
	\qquad
	f(x)=1\quad(x\in F),
	\qquad
	\operatorname{supp}f\subset \overline O.
	\]
	其次证明: 若 \(K\subset O\) 是紧集, 则存在 \(f\in C_0(\mathbb R^d)\), 使得
	\[
	0\leq f\leq1,
	\qquad
	f=1\text{ 于 }K\text{ 的某个邻域内},
	\qquad
	\operatorname{supp}f\subset O.
	\]
\end{exer}

\chapter{综合练习}
\label{chap:comprehensive-exercises}

\begin{exer}
	设 \(f\) 是 \(\R^d\) 上的一个函数, 满足
	\[
	f(x+l)=f(x),\qquad x,l\in\R^d,
	\]
	其中 \(l\) 具有整数坐标. 证明: 若 \(f\) 在
	\[
	I=\{x:0\leq x_j\leq 1,\ j=1,\ldots,d\}
	\]
	上 Riemann 可积, 则
	\[
	\lim_{T\to\infty}\frac1T\int_0^T
	f(t\alpha_1,\ldots,t\alpha_d)\,\dd t
	=
	\int_I f(x)\,\dd x,
	\]
	其中 \(\alpha_1,\ldots,\alpha_d\) 在有理数域上线性无关. 也就是说, 若
	\[
	r_1\alpha_1+\cdots+r_d\alpha_d=0,
	\qquad r_j\in\mathbb Q,
	\]
	则必有 \(r_1=\cdots=r_d=0\).
\end{exer}

\begin{exer}
	设 \(f\) 是开集 \(O\subset\R^d\) 上的二次连续可微函数, 并令
	\[
	f_h(x)=f(x)-\langle x,h\rangle .
	\]
	证明: 对除一个零集外的所有 \(h\in\R^d\), 矩阵
	\[
	\left(\frac{\partial^2 f}{\partial x_j\partial x_k}\right)
	\]
	在 \(O\) 中使
	\[
	\frac{\partial f_h}{\partial x_j}=0,
	\qquad j=1,\ldots,d,
	\]
	的任意一点上都是可逆的.
\end{exer}

\begin{exer}
	在实数轴上定义函数 \(f\): 当 \(x\) 为无理数时 \(f(x)=0\), 并且
	\[
	f(0)=0,
	\qquad
	f(x)=\frac1q\quad\text{当 }x=\frac pq,
	\]
	其中 \(p,q\) 互素且 \(q>0\). 证明 \(f\) 是上方半连续函数.
\end{exer}

\begin{exer}
	从闭区间 \([0,1]\) 的中央删去一个长度为 \(1/p_1\) 的开区间. 再从余下各个闭子区间的中央同时删去一个开区间, 其长度是该子区间长度的 \(1/p_2\). 如此重复进行, 作到第 \(n\) 步得到闭集 \(F_n\). 证明
	\[
	m\left(\bigcap_{n=1}^{\infty}F_n\right)
	=
	\prod_{n=1}^{\infty}\left(1-\frac1{p_n}\right).
	\]
\end{exer}

\begin{exer}
	求所有在十进制记数法的表示中不出现数字 \(7\) 的实数所成集合的测度.
\end{exer}

\begin{exer}
	设 \(E\) 是一个实数集合. 当 \(E\) 中的数用十进制记数法表示时, 任意一个有限的数字序列都在 \(E\) 中出现. 证明 \(E^c\) 是一个零集.
\end{exer}

\begin{exer}
	设 \(E_1,E_2,\ldots\) 是 \(\R^d\) 的子集, 并且
	\[
	\sum_{j=1}^{\infty}m^*(E_j)<+\infty.
	\]
	证明集合
	\[
	\{x:x\in E_j\text{ 对无限多个 }j\text{ 成立}\}
	\]
	是零集.
\end{exer}

\begin{exer}
	设 \((a_n)_{n=1}^{\infty}\) 为实数序列, 并假定存在常数 \(K\), 使得任一长度为 \(1\) 的区间都至多包含 \(K\) 个 \(a_n\). 证明: 若 \(f\in L^1(\R)\), 则级数
	\[
	\sum_{n=1}^{\infty} f(x+a_n)
	\]
	对几乎所有的 \(x\) 绝对收敛.
\end{exer}

\begin{exer}
	证明: 若 \(f\in L^1(\R)\), 则级数
	\[
	\sum_{n=1}^{\infty}f(nx)
	\]
	对几乎所有的 \(x\) 绝对收敛.
\end{exer}

\begin{exer}
	设 \((a_n)_{n=1}^{\infty}\) 和 \((b_n)_{n=1}^{\infty}\) 为实数序列, 且 \(b_n\geq0\). 证明级数
	\[
	\sum_{n=1}^{\infty} b_n |x-a_n|^{-\alpha}
	\]
	当
	\[
	\sum_{n=1}^{\infty} b_n^{1/\alpha}<+\infty,
	\qquad \alpha>1,
	\]
	时对几乎所有的 \(x\) 收敛.
\end{exer}

\begin{exer}
	对于一个实数 \(a\), 令 \(E_a\) 为 \(\R^d\) 中所有满足下述条件的 \(x\) 之集合: 不等式
	\[
	|\langle l,x\rangle|
	=|l_1x_1+\cdots+l_dx_d|
	\geq \|l\|^{-a}
	\]
	对除有限个之外的所有整数向量 \(l=(l_1,\ldots,l_d)\) 成立. 证明: 若 \(a<d-1\), 则 \(E_a^c\) 是一个零集. 并且证明: 当 \(d>1\) 时, 对每个 \(x\in\R^d\), 有
	\[
	\liminf_{\|l\|\to\infty}
	\|l\|^{d-1}|\langle l,x\rangle|
	\leq C_d\|x\|,
	\]
	其中 \(C_d\) 是只依赖于 \(d\) 的常数.
\end{exer}

\begin{exer}
	设 \(E\) 为实数轴的一个子集. 用 \(M\) 表示所有与 \(E\) 中无穷多个点模 \(1\) 相等的点所组成的集合. 证明: 若 \(E\) 可测, 则 \(M\) 可测; 若 \(E\) 有有限外测度, 则 \(M\) 是一个零集.
\end{exer}

\begin{exer}
	设 \(E\) 为实数轴上的一个零集. 证明: 存在实数 \(a\), 使得对任意有理数 \(r\), 都有
	\[
	a+r\notin E.
	\]
\end{exer}

\begin{exer}
	设 \((E_k)_{k=1}^{\infty}\) 为 \(\R^d\) 中可测集合的一个序列. 对 \(i=0,1,2,\ldots\), 用 \(A_i\) 表示那些恰好属于 \(i\) 个集合 \(E_k\) 的点所组成的集合. 再令 \(A_\infty\) 表示属于无穷多个 \(E_k\) 的点所组成的集合. 证明 \(A_i\) 均可测, 并且在约定 \(+\infty\cdot m(A_\infty)\) 的意义下有
	\[
	\sum_{k=1}^{\infty}m(E_k)
	=
	\sum_{i=1}^{\infty} i m(A_i)+(+\infty)m(A_\infty).
	\]
	特别地, 若 \(m(A_\infty)=0\), 则
	\[
	\sum_{k=1}^{\infty}m(E_k)
	=
	\sum_{i=1}^{\infty} i m(A_i).
	\]
\end{exer}

\begin{exer}
	设 \(E\) 是 \((0,1)\) 中的一个可测集, 且 \(m(E)>1/2\). 证明存在 \(x,y\in E\), 使得
	\[
	x+y=1.
	\]
\end{exer}

\begin{exer}
	证明: 若实数轴上 \(f\in\mathcal I^+\), 则
	\[
	\int_{\R}^{*} f(x)\,\dd x
	\leq
	\liminf_{n\to\infty}\frac1n\sum_{k=-n^2}^{n^2}f\left(\frac{k}{n}\right).
	\]
\end{exer}

\begin{exer}
	设 \(f\in L^p(\R^d)\), \(g\in L^q(\R^d)\), 其中
	\[
	1\leq p,q\leq +\infty,
	\qquad
	\frac1p+\frac1q=1.
	\]
	证明函数
	\[
	x\longmapsto \int_{\R^d} f(x-y)g(y)\,\dd y
	\]
	是 \(x\) 的连续函数.
\end{exer}

\begin{exer}
	设 \(A,B\subset\R^d\) 为可测集合, 并假定对每个 \(x\in\R^d\), 存在一个零集 \(N_x\), 使得
	\[
	A+x\subset B\cup N_x.
	\]
	证明: 若 \(A\) 不是零集, 则 \(B^c\) 是零集.
\end{exer}

\begin{exer}
	假定 \(A\) 是一个可测且具有正测度的实数集合. 证明
	\[
	\{x-y:x,y\in A\}
	\]
	包含原点的一个邻域.
\end{exer}

\begin{exer}
	证明: 若实数轴上 \(f\in L^1\), 则
	\[
	\lim_{n\to\infty}
	\sum_{k=-n^2}^{n^2}
	\left|
	\int_{k/n}^{(k+1)/n}f(x)\,\dd x
	\right|
	=
	\int_{-\infty}^{+\infty}|f(x)|\,\dd x.
	\]
\end{exer}

\begin{exer}
	设 \(f\in L^1(\R)\), \(g\in L^\infty(\R)\), 并且 \(g\) 是周期为 \(T\) 的函数. 证明
	\[
	\lim_{n\to\infty}\int_{\R} f(x)g(nx)\,\dd x
	=
	\left(\int_{\R}f(x)\,\dd x\right)
	\left(\frac1T\int_0^T g(x)\,\dd x\right).
	\]
\end{exer}

\begin{exer}
	证明: 若级数
	\[
	\sum_{n=1}^{\infty}a_n\sin nx
	\]
	在一个正测度集合中的所有 \(x\) 处绝对收敛, 则
	\[
	\sum_{n=1}^{\infty}|a_n|<+\infty.
	\]
\end{exer}

\begin{exer}
	证明: 若 \(f\in L^1(\R)\), 则
	\[
	\lim_{T\to\infty}
	\int_{-\infty}^{+\infty}
	\left|
	\frac1{2T}\int_{-T}^{T} f(s+t)\,\dd t
	\right|\dd s
	=
	\left|\int_{-\infty}^{+\infty}f(t)\,\dd t\right|.
	\]
\end{exer}

\begin{exer}
	设 \(f_1,f_2,\ldots\) 是区间 \((0,1)\) 上的可测函数, 并且
	\[
	|f_n(x)|\leq1.
	\]
	证明: 若
	\[
	\int_0^y f_n(x)\,\dd x\to0,
	\qquad n\to\infty,
	\]
	对所有 \(y\in(0,1)\) 成立, 则对所有 \(g\in L^1(0,1)\), 有
	\[
	\int_0^1 g(x)f_n(x)\,\dd x\to0,
	\qquad n\to\infty.
	\]
\end{exer}

\begin{exer}
	设 \(f\in L^1\). 证明: 对可测集 \(E\), 当 \(m(E)\to0\) 时,
	\[
	\int_E f(x)\,\dd x\to0.
	\]
	更强地, 有
	\[
	\int_E |f(x)|\,\dd x\to0.
	\]
\end{exer}

\begin{exer}
	设 \(0\leq f_n\in L^1\). 假定对几乎所有的 \(x\), 有
	\[
	f_n(x)\to f(x),
	\]
	并且
	\[
	\int f_n(x)\,\dd x\to A,
	\qquad n\to\infty.
	\]
	证明 \(f\in L^1\), 并且
	\[
	\int |f_n-f|\,\dd x\to A-\int f\,\dd x,
	\qquad n\to\infty.
	\]
\end{exer}

\begin{exer}
	设 \(f,f_1,f_2,\ldots\in L^p\), 其中 \(1\leq p<+\infty\). 假定对几乎所有的 \(x\), 有
	\[
	f_n(x)\to f(x),
	\]
	并且
	\[
	\|f_n\|_p\to \|f\|_p.
	\]
	证明
	\[
	\|f_n-f\|_p\to0.
	\]
\end{exer}

\begin{exer}
	设 \(f\in L^p\), \(g_n\in L^q\), \(n=1,2,\ldots\), 其中
	\[
	1<p<+\infty,
	\qquad
	\frac1p+\frac1q=1.
	\]
	证明: 若 \(\|g_n\|_q\leq C\), 且 \(g_n(x)\to g(x)\) 几乎处处成立, 则 \(g\in L^q\), 并且
	\[
	\int f g_n\,\dd x\to \int f g\,\dd x.
	\]
\end{exer}

\begin{exer}
	设 \(f_n\in L^p\), \(g_n\in L^q\), 其中
	\[
	1<p<+\infty,
	\qquad
	\frac1p+\frac1q=1.
	\]
	证明: 若
	\[
	|f_n|\leq F\in L^p\quad\text{对所有 }n,
	\qquad
	\|g_n\|_q\leq C,
	\]
	且
	\[
	f_n(x)\to f(x),
	\qquad
	g_n(x)\to g(x)
	\]
	几乎处处成立, 则 \(f\in L^p\), \(g\in L^q\),
	\[
	\|f-f_n\|_p\to0,
	\]
	并且
	\[
	\int f_n g_n\,\dd x\to \int f g\,\dd x.
	\]
\end{exer}

\begin{exer}
	设 \(f_n\) 为非负可测函数序列, 且
	\[
	\int f_n(x)\,\dd x=1
	\]
	对每个 \(n\) 成立, 并且 \(f_n(x)\to0\) 对所有 \(x\) 成立. 试问函数
	\[
	f=\sup_n f_n
	\]
	是否一定可测? 是否一定可积?
\end{exer}

\begin{exer}
	设 \(f(x,y)\) 对 \((x,y)\in\R^2\) 有定义. 假定 \(f(x,y)\) 对任一固定 \(y\) 是 \(x\) 的连续函数, 并且对任一固定 \(x\) 是 \(y\) 的可测函数. 证明 \(f\) 是 \(\R^2\) 上的可测函数.
\end{exer}

\begin{exer}
	试问 \((0,1)\) 中一个开子集的边界是否必定是零测集?
\end{exer}

\begin{exer}
	设 \(E\) 是一个实数集合, 并且存在常数 \(k<1\), 使得对任意区间 \(I\), 有
	\[
	m^*(E\cap I)\leq k m(I).
	\]
	证明 \(E\) 是零集.
\end{exer}

\begin{exer}
	构造一个可测集合 \(E\subset\R^d\), 使得对任一测度为有限值的非空开集 \(O\), 都有
	\[
	0<m(E\cap O)<m(O).
	\]
\end{exer}

\begin{exer}
	设 \(E\) 是区间 \(I\subset\R^d\) 的一个子集, 并令 \(0<\varepsilon<1\). 若 \(I_1,\ldots,I_N\) 是区间 \(I\) 的一个剖分, 用 \(I'\) 表示所有满足
	\[
	m^*(I_j\cap E)>\varepsilon m(I_j)
	\]
	的 \(I_j\) 的并. 证明当剖分的细度趋于零时,
	\[
	m(I')\to m^*(E).
	\]
	再用同样方法定义 \(I''\), 但将 \(E\) 替换为 \(I\setminus E\). 证明 \(E\) 可测的充分必要条件是
	\[
	m(I'\cap I'')\to0
	\]
	随剖分细度趋于零成立.
\end{exer}

\begin{exer}
	设 \(f_n\) 是在 \((-\infty,+\infty)\) 中取值的可测函数序列. 证明: 若对每个 \(\varepsilon>0\), 有
	\[
	m\{x:|f_m(x)-f_n(x)|\geq\varepsilon\}\to0,
	\qquad n,m\to\infty,
	\]
	则可以选出一个子序列, 使它对几乎所有的 \(x\) 收敛.
\end{exer}

\begin{exer}
	设非负可测函数 \(f\) 和 \(g\) 定义在可测集 \(E\subset\R^d\) 上. 令
	\[
	E_y=\{x\in E:g(x)\geq y\},
	\qquad
	h(y)=\int_{E_y}f(x)\,\dd x.
	\]
	证明
	\[
	\int_E f(x)g(x)\,\dd x
	=
	\int_0^{\infty}h(y)\,\dd y.
	\]
\end{exer}

\begin{exer}
	假定 \(f\) 在区间 \([0,1]\) 上连续, 而 \(g\) 是 \([0,1]\) 上的可测函数, 且
	\[
	0\leq g(x)\leq1.
	\]
	证明极限
	\[
	\lim_{n\to\infty}\int_0^1 f\bigl(g(x)^n\bigr)\,\dd x
	\]
	存在, 并给出其值.
\end{exer}

\begin{exer}
	设 \(f\) 在实数轴上可积. 证明对几乎所有的 \(a\), 有
	\[
	\lim_{n\to\infty}\int_0^1 f(x^n+a)\,\dd x=f(a).
	\]
\end{exer}

\begin{exer}
	设 \(E\) 是 \(\R\) 的子集. 证明
	\[
	m^*(E)
	=
	\inf_{f\in K}
	\left\{
	\lim_{x\to+\infty}\bigl(f(x)-f(-x)\bigr)
	\right\},
	\]
	其中 \(K\) 是所有在 \(\R\) 上递增, 并且在 \(E\) 的每一点处微商存在且不小于 \(1\) 的函数 \(f\) 所组成的集合.
\end{exer}

\begin{exer}
	设 \(f\) 是 \((0,1)\) 上的一个增函数. 证明
	\[
	\int_0^1 f'(x)\,\dd x
	\]
	最大不超过 \(f\) 的值域的测度.
\end{exer}

\chapter{综合练习的提示}
\label{chap:comprehensive-exercise-hints}

\begin{enumerate}
	\item 根据 Weierstrass 逼近定理, 一个连续周期函数可以用三角多项式
	\[
	\sum c_l \exp(2\pi i\langle l,x\rangle)
	\]
	一致逼近, 其中 \(l\) 的坐标为整数. 修改式 \((1.1.5)\), 使其中的
	\(h_1\) 和 \(h_2\) 为三角多项式. 再思考
	\[
	f(x)=\exp(2\pi i\langle l,x\rangle)
	\]
	这种特征函数情形下结论的含义.
	
	\item 对映射 \(f'\) 应用 Morse--Sard 定理. 注意, 若使用 Lebesgue 零集, 则不必限制在紧集 \(K\) 上.
	
	\item 对任意 \(a>0\), 集合
	\[
	\{x:f(x)\geq a\}
	\]
	是离散集, 因而是闭集. 其次, 当 \(a\leq0\) 时, 该集合为全体实数集.
	
	\item 使用 Beppo Levi 定理.
	
	\item 这只是练习 4 的一个变体. 这里每次删去区间长度的 \(1/10\), 只是删去的不一定是正中间的那个子区间. 所得集合的测度为零.
	
	\item 当只考虑一个固定的有限数字串时, 问题与练习 5 基本相同. 最后注意, 有限数字串只有可数多个.
	
	\item 考虑特征函数之和.
	
	\item 与定理 \ref{thm:riesz-fischer-l1} 前面的练习比较. 证明
	\[
	\sum_{n=1}^{\infty}|f(x+a_n)|
	\]
	是局部可积的.
	
	\item 考虑
	\[
	\int g(x)\sum_{n=1}^{\infty}|f(nx)|\,\dd x,
	\]
	其中 \(g\in C_0^+\), 且 \(g\) 在靠近 \(0\) 的地方恒等于零.
	
	\item 如果函数 \(|x-a_n|^{-\alpha}\) 是可积的, 那么这基本上是练习 8. 若单独考虑满足
	\[
	|x-a_n|\leq b_n^{1/\alpha}
	\]
	对无穷多个 \(n\) 成立的那些点 \(x\), 则先证明这些点构成一个零集, 然后排除这些奇性并像练习 8 那样进行证明.
	
	\item 估计所有满足
	\[
	\|x\|\leq M,
	\qquad
	|l_1x_1+\cdots+l_dx_d|\leq \|l\|^{-a}
	\]
	的 \(x\) 的测度. 判断
	\[
	\sum_{l\ne0}\|l\|^{-b}
	\]
	在什么条件下收敛.
	
	若不用积分理论直接证明第二个论断会比较复杂. 当然, 若存在 \(l\ne0\) 使得
	\[
	\langle l,x\rangle=0,
	\]
	则结论平凡. 在其他情形, 考虑所有纯量积
	\[
	\langle l,x\rangle,
	\qquad \|l\|\leq R.
	\]
	它们是互不相同的, 并且都位于区间
	\[
	(-R\|x\|,R\|x\|)
	\]
	之内. 因为这样的整数向量 \(l\) 的个数至少为
	\[
	C(R-\sqrt d)^d,
	\]
	其中 \(C\) 是单位球的体积, 所以可以找到 \(l_1(R)\) 和 \(l_2(R)\), 使得
	\[
	0<
	|\langle x,l_1(R)\rangle-\langle x,l_2(R)\rangle|
	\leq
	2R\|x\|
	\bigl(C(R-\sqrt d)^d-1\bigr)^{-1}\to0.
	\]
	于是
	\[
	\liminf_{R\to\infty}
	\|l(R)\|^{d-1}|\langle x,l(R)\rangle|
	\leq
	2^dC^{-1}\|x\|,
	\]
	其中
	\[
	l(R)=l_1(R)-l_2(R),
	\qquad \|l(R)\|\to\infty.
	\]
	
	\item 若 \(f\) 是 \(E\) 或 \(E^c\) 中某个可积集合的特征函数, 则有练习 8 的论断.
	
	\item 可数多个零集之并仍是零集.
	
	\item 引进特征函数.
	
	\item 否则 \(E\) 和
	\[
	E_1=\{x:1-x\in E\}
	\]
	就应当是分离的. 这样它们的并集包含于 \((0,1)\), 但测度大于 \(1\), 矛盾.
	
	\item 若 \(f\) 为连续函数, 则等号成立. 对一般的 \(f\in\mathcal I^+\), 用连续函数从下方近似 \(f\).
	
	\item 由 H\"older 不等式有
	\[
	\left|
	\int f(x-y)g(y)\,\dd y
	\right|
	\leq
	\|f\|_p\|g\|_q.
	\]
	再按 \(L^p\) 或 \(L^q\) 范数, 用连续函数近似 \(f\) 或 \(g\).
	
	\item 引进 \(A\) 和 \(B^c\) 的可积子集的特征函数, 并利用 Fubini 定理后面的练习.
	
	\item 若 \(f\) 为 \(E\) 的特征函数, 则所考虑的集合包含所有满足
	\[
	\int f(x+y)f(y)\,\dd y\neq0
	\]
	的点 \(x\). 为什么? 可以假定 \(E\) 具有有限测度. 由练习 17, 上述积分是 \(x\) 的连续函数, 并且在原点处不等于零.
	
	\item 当 \(f\in C_0\) 时结论显然. 一般情形根据 \(L^1\) 的定义作近似.
	
	\item 可以先假定 \(g\geq0\). 若 \(f\in C_0\), 结论显然成立: 将积分区间分解成 \(nx\) 落在 \(g\) 的若干周期中的部分, 然后应用积分中值定理. 对一般 \(f\in L^1\), 用 \(C_0\) 函数逼近.
	
	当 \(g\) 为正弦函数或余弦函数时, 这个特殊情形称为 Riemann--Lebesgue 引理.
	
	\item 各项绝对值之和在某个正测度集合上是有界的. 在这个集合上积分, 并利用练习 21.
	
	\item 当 \(f\in C_0\) 时结论容易. 一般情形用 \(C_0\) 函数逼近.
	
	\item 首先考虑阶梯函数 \(g\) 的情形; 这些阶梯函数在 \(L^1\) 中稠密.
	
	\item 用一个 \(C_0\) 中的函数近似 \(f\).
	
	\item 利用 Fatou 引理, 并考虑
	\[
	g_n=\min(f,f_n).
	\]
	注意 \(A=\int f\,\dd x\) 的特殊情形.
	
	\item 根据 Lusin 定理和练习 25, 可以找到紧集 \(K\), 使得
	\[
	\int_{K^c}|f|^p\,\dd x
	\]
	任意小, 同时还使 \(f_n\) 在 \(K\) 上一致收敛到 \(f\). 再利用练习 26.
	
	\item 利用 Lusin 定理、练习 25 和 H\"older 不等式.
	
	\item 利用练习 28.
	
	\item \begin{enumerate}
		\item \(f=\sup_n f_n\) 显然是可测函数.
		\item \(f\) 不一定可积. 否则由 Lebesgue 控制收敛定理可推出
		\[
		\int f_n\,\dd x\to0,
		\]
		这与 \(\int f_n\,\dd x=1\) 矛盾.
	\end{enumerate}
	
	\item 令
	\[
	X_k(x)=\frac jk,
	\qquad
	\frac jk\leq x<\frac{j+1}{k},
	\]
	其中 \(j\) 为整数. 则
	\[
	f(x,y)=\lim_{k\to\infty}f(X_k(x),y),
	\]
	因而 \(f\) 是可测函数.
	
	\item 不一定. 可以由练习 4 得到一个反例.
	
	\item 先把不等式推广到开集 \(I\). 可以假定
	\[
	m^*(E)<+\infty.
	\]
	于是可得
	\[
	m^*(E)\leq k\,m^*(E),
	\]
	从而 \(m^*(E)=0\).
	
	\item 注意, 只需关于一个开集基
	\[
	O_1,O_2,\ldots
	\]
	来作构造即可. 逐次构造 \(E_k\), 使它们关于
	\[
	O_1,\ldots,O_k
	\]
	满足要求, 并且
	\[
	m(E_k\triangle E_{k+1})
	\]
	快速趋于零.
	
	\item 可以假定 \(E\) 是可测集. 先选取有限个区间的并 \(A\), 使得
	\[
	m(E\triangle A)
	\]
	很小. 取 \(\delta<\min(\varepsilon,1-\varepsilon)\). 假定
	\[
	m(I_j\cap(A\triangle E))>\delta m(I_j)
	\]
	不成立. 于是有
	\[
	m(I_j\cap A)>(\varepsilon+\delta)m(I_j)
	\Longrightarrow
	m(I_j\cap E)>\varepsilon m(I_j),
	\]
	以及
	\[
	m(I_j\cap E)>\varepsilon m(I_j)
	\Longrightarrow
	m(I_j\cap A)>(\varepsilon-\delta)m(I_j).
	\]
	这些区间的测度有多大?
	
	\item 参见定理 \ref{thm:riesz-fischer-l1} 的证明.
	
	\item 使用 Fubini 定理. 不要忘记检验可测性.
	
	\item Lebesgue 控制收敛定理给出极限值
	\[
	f(0)m\{x:g(x)<1\}+f(1)m\{x:g(x)=1\}.
	\]
	
	\item 若
	\[
	F(t)=\int_0^t f(a+x)\,\dd x,
	\]
	则
	\[
	\frac{F(t)}{t}\to f(a)
	\]
	对几乎所有的 \(a\) 成立. 可将积分写成
	\[
	\int_0^1 F'(t)\,\dd t^{1/n}
	=
	\frac{F(1)}{n}
	+
	\frac1n\left(1-\frac1n\right)
	\int_0^1 F(t)t^{1/n-2}\,\dd t.
	\]
	先考虑原点附近的积分, 在那里
	\[
	\left|\frac{F(t)}{t}-f(a)\right|<\varepsilon.
	\]
	然后再考虑这个邻域外的积分.
	
	\item 方法一: 令 \(f\) 是包含 \(E\) 的一个开集之特征函数的积分.
	
	方法二: 利用
	\[
	\lim_{x\to+\infty}\bigl(f(x)-f(-x)\bigr)
	\geq
	\lim_{x\to+\infty}\int_{-x}^{x}f'(t)\,\dd t,
	\]
	以及在一个包含 \(E\) 的可测集合上 \(f'\) 存在并且 \(f'\geq1\).
	
	\item 值域的测度等于
	\[
	f(1)-f(0)
	\]
	减去所有不连续点处跳跃度之和. 也就是说, 它等于
	\[
	\int_0^1 f'(x)\,\dd x+\int_0^1 \dd s(x),
	\]
	其中 \(\dd s\) 是 \(\dd f\) 的奇异弥散部分.
\end{enumerate}
	
\end{document}